% Fonctions dérivées — Mathématiques 1re Tech — Cours signé OnBuch+
% Programme officiel MINESEC. Compiler avec : tectonic N04-fonctions-derivees.tex
\newcommand{\DOCMATIERE}{Mathématiques}
\newcommand{\DOCNIVEAU}{1re Tech}
\newcommand{\DOCMODULE}{Mathématiques – classe de Première technique}
\newcommand{\DOCLECON}{N04}
\newcommand{\DOCTITRE}{Fonctions dérivées}
% OnBuch+ — préambule « cours signés » ENRICHI (compilé avec tectonic / XeLaTeX)
% Système visuel « Pop » d'OnBuch+ : fond crème (jamais blanc), bordures encre
% épaisses, ombres « dures » décalées, titres Archivo Black en capitales.
% Version MATHÉMATIQUES : graphes (pgfplots/tikz), boîtes pédagogiques
% (définition, propriété, méthode, attention, le savais-tu, exercice résolu,
% exercice, TP, bilan), page de garde et signature OnBuch+ ancrée en bas.
%
% Macros attendues dans main.tex AVANT \input{preamble} :
%   \DOCMATIERE  \DOCNIVEAU  \DOCMODULE  \DOCLECON (numéro)  \DOCTITRE
\documentclass[11pt]{article}

\usepackage{fontspec}
\usepackage[french]{babel}
\usepackage[a4paper,margin=2cm,top=3.4cm,bottom=2.8cm,headheight=1.4cm,footskip=1.3cm]{geometry}
\usepackage{amsmath,amssymb,amsfonts}
\usepackage{mathtools} % \xmapsto, \coloneqq... souvent utilisés par les modèles
\usepackage{cancel} % simplifications barrées
% \abs{z} (module, valeur absolue) et \norm{u} (norme), utilisés par les modèles
\providecommand{\abs}{}\let\abs\relax\DeclarePairedDelimiter{\abs}{\lvert}{\rvert}
\providecommand{\norm}{}\let\norm\relax\DeclarePairedDelimiter{\norm}{\lVert}{\rVert}
\usepackage[table]{xcolor} % [table] : fournit aussi \rowcolors (alternance de lignes)
\usepackage{pagecolor}
\usepackage{tikz}
\usetikzlibrary{arrows.meta,positioning,calc,shapes.geometric,shapes.misc,shapes.arrows,decorations.pathmorphing,decorations.pathreplacing,decorations.markings,patterns,shadows,mindmap,trees,fit,backgrounds,matrix,intersections,angles,quotes}
\usepackage{pgfplots}
\usepackage{pgf-pie} % diagrammes circulaires \pie (statistiques)
\pgfplotsset{compat=1.18}
\usepgfplotslibrary{fillbetween,groupplots} % aires entre courbes (calcul intégral)
\usepackage{enumitem}
\usepackage{fancyhdr}
% [most] charge toutes les bibliothèques tcolorbox (listings, minted, raster,
% documentation...) et gonfle inutilement la mémoire TeX sur un document long
% (~300 boîtes) : on ne charge que ce qui est réellement utilisé.
\usepackage[skins,breakable]{tcolorbox}
\usepackage{graphicx}
\usepackage{colortbl}
\usepackage{booktabs}
\usepackage{tabularx}
\usepackage{diagbox} % case d'en-tête diagonale des tableaux à double entrée
% Colonnes extensibles centrée / gauche / droite (C, L, R), souvent utilisées par les modèles
\newcolumntype{C}{>{\centering\arraybackslash}X}
\newcolumntype{L}{>{\raggedright\arraybackslash}X}
\newcolumntype{R}{>{\raggedleft\arraybackslash}X}
\usepackage{array}
\usepackage{multicol}
\usepackage{siunitx}
% parse-numbers=false : accepte aussi \SI{2,5 \times 10^{-3}}{...}
\sisetup{locale=FR,output-decimal-marker={,},group-minimum-digits=4,parse-numbers=false}
\DeclareSIUnit\ohm{\ensuremath{\symup{\Omega}}} % U+2126 absent de la police
\usepackage{qrcode}
% \degree : commande fréquemment utilisée par les modèles pour un angle ou une
% température en dehors de \SI{}{} (non fournie par les paquets chargés).
\providecommand{\degree}{^{\circ}}
% \qed (fin de démonstration), utilisé par les modèles sans amsthm
\providecommand{\qed}{\hfill$\square$}
% \barre : double barre des valeurs interdites dans un tableau de variations (tkz-tab)
\providecommand{\barre}{\ensuremath{\|}}
% environnement proof (amsthm non chargé) utilisé par les modèles
\ifdefined\proof\else\newenvironment{proof}[1][Démonstration]{\par\noindent\textit{#1.}\ }{\qed\par}\fi
\usepackage{lastpage}
\usepackage{hyperref}
\hypersetup{hidelinks}

% ============================================================
% Palette « Pop » OnBuch+ — mêmes hex que la classe P (plus_ui.dart)
% ============================================================
\definecolor{poporange}{HTML}{F59321}
\definecolor{popdark}{HTML}{E07A0C}
\definecolor{popcream}{HTML}{FFF6E8}
\definecolor{popink}{HTML}{1C1714}
\definecolor{poppurple}{HTML}{7A5AE0}
\definecolor{popgreen}{HTML}{2E9E6A}
\definecolor{poppink}{HTML}{E0457B}
\definecolor{popblue}{HTML}{2F7DE1}
\definecolor{popgold}{HTML}{C98A12}
\definecolor{popmuted}{HTML}{8A7F76}
\definecolor{popline}{HTML}{EADFD2}
\definecolor{popgray}{HTML}{8A7F76}
\colorlet{popgrey}{popgray}
% Alias tolérants (le modèle nomme parfois les couleurs différemment)
\colorlet{popwhite}{white}
\colorlet{popblack}{popink}
\colorlet{popred}{poppink}
\colorlet{popyellow}{popgold}
% Teintes claires pour fonds de tableaux / zones
\colorlet{popblueL}{popblue!10!white}
\colorlet{poporangeL}{poporange!14!white}
\colorlet{popgreenL}{popgreen!12!white}
\colorlet{poppurpleL}{poppurple!12!white}
\colorlet{poppinkL}{poppink!12!white}
\colorlet{popgoldL}{popgold!14!white}

\pagecolor{popcream}

% ============================================================
% Typographie
% ============================================================
\defaultfontfeatures{Ligatures=TeX}
\setmainfont{PlusJakartaSans-Medium.ttf}[Path=../../../fonts/,BoldFont=PlusJakartaSans-Bold.ttf]
% Maths en sans-serif (Fira Math) avec chiffres et lettres droites de Plus
% Jakarta, pour que les formules se fondent dans le texte.
\usepackage[math-style=ISO,bold-style=ISO]{unicode-math}
\setmathfont{FiraMath-Regular.otf}[Path=../../../fonts/]
\setmathfont{PlusJakartaSans-Medium.ttf}[Path=../../../fonts/,range={up/num,up/{Latin,latin},\mathup}]
\usepackage{icomma} % virgule décimale sans espace en mode math
% Caractères mathématiques Unicode absents des polices de texte (Plus Jakarta,
% Archivo Black) : rendus par leur équivalent mathématique, en texte comme en
% formule — sinon ils s'affichent en carré vide (ex. « Arithmétique dans ℤ »).
\usepackage{newunicodechar}
\newunicodechar{ℝ}{\ensuremath{\mathbb{R}}}
\newunicodechar{ℤ}{\ensuremath{\mathbb{Z}}}
\newunicodechar{ℂ}{\ensuremath{\mathbb{C}}}
\newunicodechar{ℕ}{\ensuremath{\mathbb{N}}}
\newunicodechar{ℚ}{\ensuremath{\mathbb{Q}}}
\newunicodechar{𝔻}{\ensuremath{\mathbb{D}}}
\newunicodechar{α}{\ensuremath{\alpha}}
\newunicodechar{β}{\ensuremath{\beta}}
\newunicodechar{γ}{\ensuremath{\gamma}}
\newunicodechar{δ}{\ensuremath{\delta}}
\newunicodechar{ε}{\ensuremath{\varepsilon}}
\newunicodechar{θ}{\ensuremath{\theta}}
\newunicodechar{λ}{\ensuremath{\lambda}}
\newunicodechar{μ}{\ensuremath{\mu}}
\newunicodechar{σ}{\ensuremath{\sigma}}
\newunicodechar{τ}{\ensuremath{\tau}}
\newunicodechar{φ}{\ensuremath{\varphi}}
\newunicodechar{ψ}{\ensuremath{\psi}}
\newunicodechar{ω}{\ensuremath{\omega}}
\newunicodechar{Σ}{\ensuremath{\Sigma}}
% majuscules produites par \MakeUppercase dans les titres (α-aminés -> Α)
\newunicodechar{Α}{\ensuremath{\alpha}}
\newunicodechar{Β}{\ensuremath{\beta}}
\newunicodechar{Γ}{\ensuremath{\Gamma}}
\newunicodechar{Δ}{\ensuremath{\Delta}}
\newunicodechar{Ε}{\ensuremath{\varepsilon}}
\newunicodechar{Θ}{\ensuremath{\Theta}}
\newunicodechar{Λ}{\ensuremath{\Lambda}}
\newunicodechar{Μ}{\ensuremath{\mu}}
\newunicodechar{Τ}{\ensuremath{\tau}}
\newunicodechar{Φ}{\ensuremath{\Phi}}
\newunicodechar{Ψ}{\ensuremath{\Psi}}
\newunicodechar{Ω}{\ensuremath{\Omega}}
\newunicodechar{↦}{\ensuremath{\mapsto}}
\newunicodechar{∈}{\ensuremath{\in}}
\newunicodechar{∉}{\ensuremath{\notin}}
\newunicodechar{∧}{\ensuremath{\wedge}}
\newunicodechar{⇔}{\ensuremath{\Leftrightarrow}}
\newunicodechar{⟺}{\ensuremath{\Longleftrightarrow}}
\newunicodechar{⇒}{\ensuremath{\Rightarrow}}
\newunicodechar{⟹}{\ensuremath{\Longrightarrow}}
\newunicodechar{∘}{\ensuremath{\circ}}
\newunicodechar{∪}{\ensuremath{\cup}}
\newunicodechar{∩}{\ensuremath{\cap}}
\newunicodechar{≡}{\ensuremath{\equiv}}
\newunicodechar{⊂}{\ensuremath{\subset}}
\newunicodechar{⊥}{\ensuremath{\perp}}
\newunicodechar{∃}{\ensuremath{\exists}}
\newunicodechar{∀}{\ensuremath{\forall}}
\newunicodechar{∛}{\ensuremath{\sqrt[3]{\,}}}
\newunicodechar{∜}{\ensuremath{\sqrt[4]{\,}}}
\newunicodechar{⃗}{\ensuremath{{}^{\to}}}
\newunicodechar{ⁿ}{\ifmmode^{n}\else\textsuperscript{\ensuremath{n}}\fi}
\newunicodechar{ˣ}{\ifmmode^{x}\else\textsuperscript{\ensuremath{x}}\fi}
\newunicodechar{ᵇ}{\ifmmode^{b}\else\textsuperscript{\ensuremath{b}}\fi}
\newunicodechar{ʳ}{\ifmmode^{r}\else\textsuperscript{\ensuremath{r}}\fi}
\newunicodechar{ᵘ}{\ifmmode^{u}\else\textsuperscript{\ensuremath{u}}\fi}
\newunicodechar{ʸ}{\ifmmode^{y}\else\textsuperscript{\ensuremath{y}}\fi}
\newunicodechar{ᵃ}{\ifmmode^{a}\else\textsuperscript{\ensuremath{a}}\fi}
\newunicodechar{ᵉ}{\ifmmode^{e}\else\textsuperscript{\ensuremath{e}}\fi}
\newunicodechar{ᵏ}{\ifmmode^{k}\else\textsuperscript{\ensuremath{k}}\fi}
\newunicodechar{ᵖ}{\ifmmode^{p}\else\textsuperscript{\ensuremath{p}}\fi}
\newunicodechar{ᴺ}{\ifmmode^{N}\else\textsuperscript{\ensuremath{N}}\fi}
\newunicodechar{ᶜ}{\ifmmode^{c}\else\textsuperscript{\ensuremath{c}}\fi}
\newunicodechar{ʰ}{\ifmmode^{h}\else\textsuperscript{\ensuremath{h}}\fi}
\newunicodechar{ᵠ}{\ifmmode^{q}\else\textsuperscript{\ensuremath{q}}\fi}
\newunicodechar{ₙ}{\ifmmode_{n}\else\textsubscript{\ensuremath{n}}\fi}
\newunicodechar{ₐ}{\ifmmode_{a}\else\textsubscript{\ensuremath{a}}\fi}
\newunicodechar{ᵢ}{\ifmmode_{i}\else\textsubscript{\ensuremath{i}}\fi}
\newunicodechar{ᵣ}{\ifmmode_{r}\else\textsubscript{\ensuremath{r}}\fi}
\newunicodechar{ₖ}{\ifmmode_{k}\else\textsubscript{\ensuremath{k}}\fi}
\newunicodechar{ᵦ}{\ifmmode_{\beta}\else\textsubscript{\ensuremath{\beta}}\fi}
\newunicodechar{ₓ}{\ifmmode_{x}\else\textsubscript{\ensuremath{x}}\fi}
\newfontfamily\popdisplay{ArchivoBlack-Regular.ttf}[Path=../../../fonts/]
\newfontfamily\poplogo{SpaceGrotesk-Bold.ttf}[Path=../../../fonts/]
\newfontfamily\popsemi{PlusJakartaSans-SemiBold.ttf}[Path=../../../fonts/]
\newfontfamily\popxbold{PlusJakartaSans-ExtraBold.ttf}[Path=../../../fonts/]
\newfontfamily\popxboldls{PlusJakartaSans-ExtraBold.ttf}[Path=../../../fonts/,LetterSpace=3.0]

\setlist[enumerate]{leftmargin=1.7em,itemsep=5pt,topsep=4pt}
\setlist[itemize]{leftmargin=1.7em,itemsep=4pt,topsep=4pt,label={\color{poporange}\textbullet}}
\setlength{\parindent}{0pt}
\setlength{\parskip}{5pt}
\renewcommand{\arraystretch}{1.25}

% pgfplots : style Pop par défaut
\pgfplotsset{
  every axis/.append style={axis line style={popink,line width=1pt},tick style={popink},
    grid style={popline,line width=0.5pt},label style={font=\small},tick label style={font=\footnotesize}},
  popcurve/.style={poporange,line width=1.8pt,no markers,smooth},
}
% Même style disponible dans un \draw TikZ simple (hors pgfplots)
\tikzset{popcurve/.style={poporange,line width=1.8pt}}

% ============================================================
% Pill / squiggle
% ============================================================
\newcommand{\poppill}[3]{%
  \tikz[baseline=-0.55ex]\node[draw=popink,line width=1.2pt,rounded corners=2.6pt,%
    fill=#2,inner xsep=6.5pt,inner ysep=3pt]{%
    \popxboldls\fontsize{7}{7}\selectfont\color{#3}\MakeUppercase{#1}};%
}
\newcommand{\popsquiggle}[1]{%
  \tikz[baseline=-0.4ex]{
    \draw[#1,line width=2.6pt,line cap=round]
      plot [smooth] coordinates {(0,0.09) (0.34,0.01) (0.68,0.21) (1.02,0.01) (1.36,0.10)};
  }%
}
% Mot-clé surligné (marqueur orange sous le texte)
\newcommand{\cle}[1]{{\popxbold\color{popdark}#1}}

% ============================================================
% En-tête / pied
% ============================================================
\pagestyle{fancy}
\fancyhf{}
\renewcommand{\headrulewidth}{0pt}
\renewcommand{\footrulewidth}{0pt}
\lhead{\raisebox{-0.15ex}{\poplogo\fontsize{13}{13}\selectfont\color{popink}OnBuch\color{poporange}+}}
\rhead{\poppill{\DOCMATIERE\ \textbullet\ \DOCNIVEAU\ \textbullet\ Leçon \DOCLECON}{white}{popink}}
\lfoot{\popxbold\fontsize{7.6}{7.6}\selectfont\color{popmuted}\MakeUppercase{Cours signé OnBuch+}\ \textbullet\ {\popsemi\DOCTITRE}}
\rfoot{\tikz[baseline=-0.6ex]\node[rounded corners=8.5pt,fill=popink,minimum height=17pt,minimum width=17pt,inner xsep=5pt,inner sep=0]{\popxbold\fontsize{8}{8}\selectfont\color{popcream}\thepage\,/\,\pageref*{LastPage}};}

% ============================================================
% Titres
% ============================================================
\newcommand{\chaptitre}[2]{%
  \begin{tcolorbox}[enhanced,colback=poporange,colframe=popink,boxrule=2.6pt,arc=11pt,
      drop shadow=popink,left=15pt,right=15pt,top=12pt,bottom=11pt,before skip=3pt,after skip=18pt]
    \poppill{#2}{popink}{popcream}\par\vspace{9pt}
    {\raggedright\hyphenpenalty=10000\exhyphenpenalty=10000\popdisplay\fontsize{22}{24}\selectfont\color{popcream}\MakeUppercase{#1}\par}
  \end{tcolorbox}
}
% Section numérotée : pastille orange avec le numéro + titre Archivo
\newcounter{popsec}
\newcommand{\coursec}[1]{%
  \stepcounter{popsec}\setcounter{popsub}{0}%
  \par\vspace{16pt}\noindent
  \begin{minipage}{\linewidth}
  \tikz[baseline=-0.7ex]\node[draw=popink,line width=1.6pt,fill=poporange,rounded corners=4pt,minimum size=22pt,inner sep=0]{\popdisplay\fontsize{12}{12}\selectfont\color{popcream}\thepopsec};%
  \hspace{8pt}{\popdisplay\fontsize{15}{17}\selectfont\color{popink}\MakeUppercase{#1}}\par
  \vspace{4pt}\noindent{\color{poporange}\rule{36pt}{3.6pt}}
  \end{minipage}\par\nopagebreak\vspace{8pt}
}
\newcounter{popsub}
\newcommand{\courssub}[1]{%
  \stepcounter{popsub}%
  \par\vspace{9pt}\noindent{\popxbold\fontsize{11.5}{13}\selectfont\color{popink}{\color{poporange}\thepopsec.\thepopsub}\hspace{6pt}#1}\par\nopagebreak\vspace{4pt}
}

% ============================================================
% Boîtes pédagogiques — même recette : fond blanc, bordure épaisse colorée,
% ombre dure encre, badge plein. Titre optionnel [#1].
% ============================================================
\tcbset{popbase/.style={breakable,enhanced,colback=white,boxrule=2.3pt,arc=9pt,
  shadow={3pt}{-3pt}{0pt}{popink},left=14pt,right=14pt,top=11pt,bottom=11pt,before skip=11pt,after skip=15pt,
  fontupper=\popsemi\fontsize{10.6}{14.5}\selectfont\color{popink},
  fontlower=\popsemi\fontsize{10.6}{14.5}\selectfont\color{popink}}}
% Coche dessinée (le glyphe ✓ n'existe pas dans les polices du document)
\newcommand{\popcheck}[1]{\tikz[baseline=-0.5ex]\draw[#1,line width=1.6pt,line cap=round,line join=round](-0.13,0.0)--(-0.04,-0.09)--(0.13,0.1);}
\newcommand{\popboxhead}[3]{\poppill{#1}{#2}{white}\ifx\relax#3\relax\else\hspace{6pt}{\popxbold\fontsize{10.6}{12}\selectfont\color{popink}#3}\fi\par\vspace{7pt}}

\newtcolorbox{aretenir}[1][]{popbase,colframe=popgreen,before upper={\popboxhead{À retenir}{popgreen}{#1}}}
\newtcolorbox{exemplebox}[1][]{popbase,colframe=poporange,before upper={\popboxhead{Exemple}{poporange}{#1}}}
\newtcolorbox{definition}[1][]{popbase,colframe=popblue,before upper={\popboxhead{Définition}{popblue}{#1}}}
\newtcolorbox{propriete}[1][]{popbase,colframe=poppurple,before upper={\popboxhead{Propriété}{poppurple}{#1}}}
\newtcolorbox{methode}[1][]{popbase,colframe=popgold,before upper={\popboxhead{Méthode}{popgold}{#1}}}
\newtcolorbox{attention}[1][]{popbase,colframe=poppink,before upper={\popboxhead{Attention}{poppink}{#1}}}
\newtcolorbox{experience}[1][]{popbase,colframe=popblue,colback=popblueL,before upper={\popboxhead{Expérience}{popblue}{#1}}}
% « Le savais-tu ? » : carte violette pleine (contexte, culture, Cameroun)
\newtcolorbox{savaistu}[1][]{popbase,colback=poppurple,colframe=popink,
  fontupper=\popsemi\fontsize{10.4}{14.2}\selectfont\color{white},
  before upper={\poppill{Le savais-tu\,?}{white}{poppurple}\ifx\relax#1\relax\else\hspace{6pt}{\popxbold\color{white}#1}\fi\par\vspace{7pt}}}
% Exercice résolu : énoncé en haut, \tcblower puis la solution sur fond crème
\newcounter{popexo}\newcounter{popexr}
\newtcolorbox{exoresolu}[1][]{popbase,colframe=poporange,colbacklower=popcream,
  segmentation style={popink,line width=1.2pt,dashed},
  before upper={\stepcounter{popexr}\popboxhead{Exercice résolu \thepopexr}{poporange}{#1}},
  before lower={\poppill{Solution}{popink}{popcream}\par\vspace{6pt}}}
% Exercice (non résolu en place) — la correction est dans la partie Corrigés
\newtcolorbox{exercice}[1][]{popbase,colframe=popink,
  before upper={\stepcounter{popexo}\popboxhead{Exercice \thepopexo}{popink}{#1}}}
% Corrigé
\newtcolorbox{corrige}[1][]{popbase,colframe=popgreen,colback=popgreenL,
  before upper={\popboxhead{Corrigé}{popgreen}{#1}}}
% Objectifs / prérequis / bilan
\newtcolorbox{objectifs}{popbase,colframe=popink,colback=poporangeL,
  before upper={\popboxhead{Objectifs de la leçon}{popink}{}}}
\newtcolorbox{prerequis}{popbase,colframe=popink,colback=popblueL,
  before upper={\popboxhead{Prérequis}{popblue}{}}}
\newtcolorbox{bilan}{popbase,colframe=popink,colback=white,boxrule=2.8pt,
  before upper={\popboxhead{Fiche bilan}{poporange}{L'essentiel à connaître pour l'examen}}}
% Figure encadrée avec légende
\newtcolorbox{popfigure}[1][]{enhanced,breakable=false,colback=white,colframe=popline,boxrule=1.4pt,arc=8pt,
  left=10pt,right=10pt,top=10pt,bottom=8pt,before skip=10pt,after skip=12pt,halign=center,
  lower separated=false,halign lower=center,
  fontlower=\popsemi\fontsize{9}{11.5}\selectfont\color{popmuted},
  #1}
\newcommand{\legende}[1]{\par\vspace{6pt}{\popsemi\fontsize{9}{11.5}\selectfont\color{popmuted}#1}}

% ============================================================
% Signature OnBuch+
% ============================================================
\newcommand{\poplogoinline}[1]{{\poplogo\fontsize{#1}{#1}\selectfont\color{popink}OnBuch\color{poporange}+}}
\newcommand{\signatureonbuch}{%
  \begin{tcolorbox}[enhanced,colback=popink,colframe=popink,boxrule=2.2pt,arc=10pt,
      drop shadow=poporange,left=14pt,right=14pt,top=10pt,bottom=10pt,before skip=0pt,after skip=0pt]
    \tikz[baseline=-0.7ex]\node[circle,fill=popgreen,draw=popcream,line width=1.3pt,minimum size=22pt,inner sep=0]{\popcheck{white}};%
    \hspace{9pt}\begin{minipage}[c]{0.62\linewidth}
      {\popxbold\fontsize{12}{13}\selectfont\color{popcream}Signé\ \textbullet\ {\poplogo OnBuch\color{poporange}+}}\par\vspace{2pt}
      {\raggedright\popsemi\fontsize{8}{10.5}\selectfont\color{popline}\MakeUppercase{Équipe pédagogique OnBuch+}\\\MakeUppercase{Conforme au programme officiel MINESEC}\par}
    \end{minipage}\hfill
    {\poplogo\fontsize{22}{22}\selectfont\color{popcream}OnBuch\color{poporange}+}
  \end{tcolorbox}%
}

% ============================================================
% Page de garde
% #1 titre · #2 matière · #3 niveau · #4 module · #5 n° leçon · #6 durée ·
% #7 description / objectifs (texte ou itemize)
% ============================================================
\newcommand{\pagegarde}[7]{%
  \thispagestyle{empty}
  \newgeometry{a4paper,margin=1.7cm,top=1.6cm,bottom=1.4cm}%
  \begin{tikzpicture}[remember picture,overlay]
    \fill[poporange!18] ([xshift=-3.2cm,yshift=-3.6cm]current page.north east) circle (4.6cm);
    \fill[poppurple!14] ([xshift=2.4cm,yshift=7.6cm]current page.south west) circle (3.8cm);
  \end{tikzpicture}%
  {\poplogo\fontsize{30}{30}\selectfont\color{popink}OnBuch\color{poporange}+}\hfill
  \raisebox{0.6ex}{\poppill{Cours signé \textbullet\ Premium}{popink}{popcream}}\par
  \vspace{1pt}\popsquiggle{poppurple}\par
  \vspace{8pt}
  \begin{tcolorbox}[enhanced,colback=poporange,colframe=popink,boxrule=2.8pt,arc=13pt,
      drop shadow=popink,left=18pt,right=18pt,top=12pt,bottom=13pt,after skip=10pt]
    \poppill{#2 \textbullet\ #3}{popink}{popcream}\hspace{5pt}\poppill{#4}{white}{popink}\par\vspace{8pt}
    {\popdisplay\fontsize{14}{14}\selectfont\color{popink}LEÇON #5}\par\vspace{4pt}
    {\raggedright\hyphenpenalty=10000\exhyphenpenalty=10000\popdisplay\fontsize{25}{27}\selectfont\color{popcream}\MakeUppercase{#1}\par}
    \vspace{8pt}
    {\popxbold\fontsize{9.5}{9.5}\selectfont\color{popink}\MakeUppercase{Durée conseillée : #6}}
  \end{tcolorbox}
  \begin{tcolorbox}[enhanced,colback=white,colframe=popink,boxrule=2.2pt,arc=10pt,
      drop shadow=popink,left=16pt,right=16pt,top=10pt,bottom=10pt,after skip=8pt]
    \poppill{Dans cette leçon}{poppurple}{white}\par\vspace{5pt}
    \popsemi\fontsize{9.3}{12.6}\selectfont\color{popink}#7
  \end{tcolorbox}
  \noindent\begin{minipage}[t]{0.487\linewidth}
    \begin{tcolorbox}[enhanced,colback=popgreen,colframe=popink,boxrule=2.2pt,arc=10pt,
        drop shadow=popink,left=11pt,right=11pt,top=10pt,bottom=10pt,height=3.1cm,valign=center]
      \tcbox[colback=white,colframe=popink,boxrule=1.3pt,arc=5pt,boxsep=0pt,nobeforeafter,
        left=3pt,right=3pt,top=3pt,bottom=3pt,valign=center]{\qrcode[height=1.9cm]{https://play.google.com/store/apps/details?id=com.luvvix.onbuch&hl=fr}}%
      \hspace{8pt}\begin{minipage}[c]{0.5\linewidth}
        {\popdisplay\fontsize{11}{12.5}\selectfont\color{white}\MakeUppercase{Télécharge OnBuch}}\par\vspace{4pt}
        {\raggedright\popsemi\fontsize{8.4}{11}\selectfont\color{white}Gratuit \textbullet\ Google Play\\Tuteur IA, annales, concours, résultats\par}
      \end{minipage}
    \end{tcolorbox}
  \end{minipage}\hfill
  \begin{minipage}[t]{0.487\linewidth}
    \begin{tcolorbox}[enhanced,colback=poppurple,colframe=popink,boxrule=2.2pt,arc=10pt,
        drop shadow=popink,left=11pt,right=11pt,top=10pt,bottom=10pt,height=3.1cm,valign=center]
      \tikz[baseline=-0.7ex]\node[circle,fill=white,draw=popink,line width=1.3pt,minimum size=1.6cm,inner sep=0]{\popdisplay\fontsize{24}{24}\selectfont\color{poppurple}+};%
      \hspace{8pt}\begin{minipage}[c]{0.6\linewidth}
        {\popdisplay\fontsize{11}{12.5}\selectfont\color{white}\MakeUppercase{Passe à OnBuch+}}\par\vspace{4pt}
        {\raggedright\popsemi\fontsize{8.4}{11}\selectfont\color{white}Tous les cours signés, Tuteur IA illimité, fiches PDF — onglet OnBuch+ de l'app\par}
      \end{minipage}
    \end{tcolorbox}
  \end{minipage}
  \vspace{8pt}
  \popcontenu
  \vfill
  \signatureonbuch
  \restoregeometry
  \newpage
}

% Rangée « Ce cours contient » (page de garde)
\newcommand{\poptile}[3]{% #1 couleur #2 gros texte #3 légende
  \begin{tcolorbox}[enhanced,colback=#1,colframe=popink,boxrule=2pt,arc=9pt,shadow={2.5pt}{-2.5pt}{0pt}{popink},
      left=6pt,right=6pt,top=9pt,bottom=9pt,halign=center,valign=center,height=1.75cm,width=\linewidth,nobeforeafter]
    {\popdisplay\fontsize{12.5}{13}\selectfont\color{white}\MakeUppercase{#2}}\par\vspace{3pt}
    {\centering\popsemi\fontsize{7.8}{9.5}\selectfont\color{white}#3\par}
  \end{tcolorbox}}
\newcommand{\popcontenu}{%
  \par\noindent{\popxboldls\fontsize{7.5}{7.5}\selectfont\color{popmuted}CE COURS SIGNÉ CONTIENT}\par\vspace{7pt}
  \noindent\begin{minipage}[t]{0.235\linewidth}\poptile{popblue}{Cours}{complet, illustré, pas à pas}\end{minipage}\hfill
  \begin{minipage}[t]{0.235\linewidth}\poptile{poporange}{Méthodes}{et exercices résolus}\end{minipage}\hfill
  \begin{minipage}[t]{0.235\linewidth}\poptile{poppink}{Exercices}{type examen + corrigés}\end{minipage}\hfill
  \begin{minipage}[t]{0.235\linewidth}\poptile{popgreen}{Fiche bilan}{l'essentiel à retenir}\end{minipage}\par}

% Sommaire
\newtcolorbox{sommaire}{enhanced,colback=white,colframe=popink,boxrule=2.2pt,arc=10pt,
  drop shadow=popink,left=18pt,right=18pt,top=13pt,bottom=13pt,before skip=6pt,after skip=20pt,
  before upper={\poppill{Sommaire}{poppurple}{white}\par\vspace{9pt}\popsemi\fontsize{10.6}{14.5}\selectfont\color{popink}}}

% Signature de fin de cours — poussée tout en bas de la dernière page
\newcommand{\findecours}{%
  \par\vfill\nopagebreak
  \begin{center}\popsquiggle{poporange}\end{center}\vspace{4pt}
  \signatureonbuch
}
\AtBeginDocument{\def\llbracket{\mathopen{[\mkern-3.5mu[}}\def\rrbracket{\mathclose{]\mkern-3.5mu]}}} % intervalles d'entiers (glyphe ⟦ absent de la police)
% Glyphes absents de Fira Math : redéfinis à partir de glyphes disponibles
\AtBeginDocument{%
  \def\setminus{\mathbin{\backslash}}\let\smallsetminus\setminus
  \def\vdots{\mathord{\vbox{\baselineskip3.2pt\lineskiplimit0pt\kern4pt\hbox{.}\hbox{.}\hbox{.}}}}%
  \def\ddots{\mathinner{\mkern1mu\raise6.5pt\hbox{.}\mkern2mu\raise3.6pt\hbox{.}\mkern2mu\raise0.7pt\hbox{.}\mkern1mu}}%
  \def\triangle{\mathord{\scalebox{0.9}{$\symup{\Delta}$}}}%
  \def\nabla{\mathord{\raisebox{\depth}{\scalebox{1}[-1]{$\symup{\Delta}$}}}}%
  \def\checkmark{\popcheck{popdark}}%
  \def\star{\mathbin{\ast}}\def\bigstar{\mathord{\ast}}%
  \def\diamond{\mathbin{\scalebox{0.6}{$\Diamond$}}}%
  \def\bigcup{\mathop{\mathchoice{\raisebox{-0.3em}{\scalebox{1.7}{$\cup$}}}{\scalebox{1.25}{$\cup$}}{\cup}{\cup}}\displaylimits}%
  \def\bigcap{\mathop{\mathchoice{\raisebox{-0.3em}{\scalebox{1.7}{$\cap$}}}{\scalebox{1.25}{$\cap$}}{\cap}{\cap}}\displaylimits}%
  \def\lesseqqgtr{\mathrel{\lessgtr}}\def\bigoplus{\mathop{\oplus}\displaylimits}%
}
\providecommand{\ssi}{si et seulement si} % abréviation fréquente
\AtBeginDocument{\providecommand{\wideparen}{\overparen}} % arc de cercle

%% FIGFIX-BEGIN v3 — correctifs globaux des figures (docs/onbuch-generation/tools/figfix)
\usepackage{adjustbox}\usepackage{etoolbox}
% toute figure TikZ/pgfplots plus large que la ligne est réduite à la largeur disponible
\BeforeBeginEnvironment{tikzpicture}{\begin{adjustbox}{max width=\linewidth}}
\AfterEndEnvironment{tikzpicture}{\end{adjustbox}}
% légendes pgfplots lisibles par-dessus les courbes
\pgfplotsset{every axis/.append style={legend style={fill=white,fill opacity=0.92,text opacity=1,draw=black!15,inner sep=3pt}}}
% étiquettes d'arêtes (syntaxe edge["..."]) avec fond blanc
\tikzset{every edge quotes/.append style={fill=white,inner sep=1.2pt}}
% bulles de cartes mentales : texte à la ligne dans la bulle, bulle un peu plus grande
\tikzset{every concept/.append style={text width=2.0cm,align=center,minimum size=2.3cm,inner sep=2pt}}
%% FIGFIX-END
\begin{document}

\pagegarde{Fonctions dérivées}{Mathématiques}{1re Tech}{Mathématiques – classe de Première technique}{N04}{4 séances (fiche de progression harmonisée)}{Cette leçon introduit la notion de nombre dérivé comme pente de la tangente à une courbe, puis construit la fonction dérivée et les formules de dérivation usuelles. Elle aboutit à l'équation de la tangente en un point et à des applications concrètes (vitesse, coût marginal, optimisation simple).
\vspace{4pt}
\begin{multicols}{2}\small
\begin{itemize}
\item À la fin de cette leçon, je sais interpréter graphiquement le nombre dérivé d'une fonction en un point comme le coefficient directeur de la tangente
\item À la fin de cette leçon, je sais calculer un nombre dérivé à partir de la définition par le taux d'accroissement
\item À la fin de cette leçon, je sais dériver les fonctions usuelles (constante, puissance, racine carrée, inverse, polynômes)
\item À la fin de cette leçon, je sais appliquer les règles de dérivation d'une somme, d'un produit par un scalaire, d'un produit et d'un quotient
\item À la fin de cette leçon, je sais déterminer l'équation réduite de la tangente à une courbe en un point d'abscisse donnée
\item À la fin de cette leçon, je sais lire graphiquement un nombre dérivé et le signe d'une dérivée sur une courbe donnée
\end{itemize}\end{multicols}}

\begin{sommaire}
\begin{multicols}{2}
\begin{enumerate}
\item Approche intuitive du nombre dérivé
\item Tangente à une courbe en un point
\item Fonction dérivée et dérivées des fonctions usuelles
\item Opérations sur les fonctions dérivées
\item Applications concrètes de la dérivation
\item Synthèse, compléments et vers la Terminale
\item Activité d'intégration
\item Méthodes et exercices résolus
\item Exercices
\item Corrigés des exercices
\item Fiche bilan
\end{enumerate}
\end{multicols}
\end{sommaire}

\begin{objectifs}
\begin{itemize}
\item À la fin de cette leçon, je sais interpréter graphiquement le nombre dérivé d'une fonction en un point comme le coefficient directeur de la tangente
\item À la fin de cette leçon, je sais calculer un nombre dérivé à partir de la définition par le taux d'accroissement
\item À la fin de cette leçon, je sais dériver les fonctions usuelles (constante, puissance, racine carrée, inverse, polynômes)
\item À la fin de cette leçon, je sais appliquer les règles de dérivation d'une somme, d'un produit par un scalaire, d'un produit et d'un quotient
\item À la fin de cette leçon, je sais déterminer l'équation réduite de la tangente à une courbe en un point d'abscisse donnée
\item À la fin de cette leçon, je sais lire graphiquement un nombre dérivé et le signe d'une dérivée sur une courbe donnée
\item À la fin de cette leçon, je sais modéliser une situation concrète (vitesse instantanée, coût marginal, recette) à l'aide de la dérivée
\item À la fin de cette leçon, je sais rédiger et justifier une démarche complète : calcul de dérivée, calcul de f(a) et f'(a), équation de tangente, conclusion
\end{itemize}
\end{objectifs}

\begin{prerequis}
\begin{itemize}
\item Équation réduite d'une droite et coefficient directeur (classe de Seconde)
\item Fonctions de référence : x ↦ x², x ↦ x³, x ↦ 1/x, x ↦ √x et leurs courbes
\item Notion de fonction, image d'un nombre, lecture graphique d'une courbe
\item Calcul algébrique : développement, factorisation, quantité conjuguée pour les racines
\item Taux d'accroissement moyen d'une fonction entre deux points (début d'année de Première)
\end{itemize}
\end{prerequis}

\newpage

\coursec{Approche intuitive du nombre dérivé}

\courssub{Une situation concrète : la vitesse d'un taxi-brousse sur l'axe Douala--Yaoundé}

Imaginez un chauffeur de taxi-brousse qui effectue le trajet Douala--Yaoundé (\SI{300}{km}). Il part à 6\,h et arrive à 10\,h. Sa \textbf{vitesse moyenne} est facile à calculer :
\[
v_{\text{moyenne}} = \frac{\text{distance parcourue}}{\text{durée du trajet}} = \frac{300}{4} = \SI{75}{km/h}.
\]
Mais cette vitesse moyenne ne dit rien de ce qui s'est passé \emph{à un instant précis} : a-t-il roulé à \SI{100}{km/h} dans la descente de Nkolbisson ? A-t-il ralenti à \SI{30}{km/h} dans les travaux à Édéa ? Pour connaître sa \textbf{vitesse instantanée} à 8\,h\,15, il faut regarder ce qui se passe sur un intervalle de temps \emph{très court} autour de cet instant.

\medskip
\noindent C'est exactement le même problème que celui de Mama Ngo Bell au marché Mokolo : elle veut savoir comment sa recette varie \emph{au moment où} elle change le prix de \SI{10}{FCFA}. Ou celui de l'économiste de la SABC : à partir de quelle quantité le coût marginal (la variation instantanée du coût) devient-il trop élevé ? Toutes ces questions reviennent à mesurer la \textbf{pente instantanée} d'une courbe. L'outil mathématique qui répond à cette question est le \textbf{nombre dérivé}, et sa généralisation est la \textbf{fonction dérivée}.

\courssub{Taux d'accroissement et sécante : la pente sur un intervalle}

Soit $f$ une fonction définie sur un intervalle $I$ et $a \in I$. On se place en $A(a\,;\,f(a))$ sur la courbe $\mathcal{C}_f$. On considère un point $M$ d'abscisse $a+h$ ($h \neq 0$, assez petit pour que $a+h \in I$). Le point $M$ a pour coordonnées $(a+h\,;\,f(a+h))$.

\begin{definition}[Taux d'accroissement (ou taux de variation)]
Le \cle{taux d'accroissement} de $f$ en $a$ pour un écart $h \neq 0$ est le nombre réel :
\[
\tau(h) = \frac{f(a+h) - f(a)}{h}.
\]
\end{definition}

\begin{propriete}[Interprétation géométrique]
$\tau(h)$ est le \textbf{coefficient directeur} de la droite $(AM)$ (la \cle{sécante} à la courbe $\mathcal{C}_f$ passant par $A$ et $M$).
\[
\tau(h) = \frac{y_M - y_A}{x_M - x_A} = \frac{f(a+h)-f(a)}{(a+h)-a}.
\]
\end{propriete}

\begin{popfigure}
\begin{tikzpicture}[scale=0.9, >=Stealth]
\begin{axis}[
    width=\linewidth, height=8cm,
    axis lines=middle,
    xlabel=$x$, ylabel=$y$,
    xmin=-0.5, xmax=3.5,
    ymin=-0.5, ymax=4.5,
    xtick={1,2,3}, ytick={1,2,3,4},
    tick label style={font=\small},
    clip=false,
    domain=-0.5:3.5, samples=100,
]
% Parabole y = x^2
\addplot[popcurve, thick, popblue] {x^2};
\node[popblue, anchor=south west] at (axis cs:3.2, 3.5) {$\mathcal{C}_f : y=x^2$};

% Point A(1,1)
\addplot[only marks, mark=*, mark size=2.5, popdark] coordinates {(1,1)};
\node[popdark, anchor=south east] at (axis cs:1,1) {$A(1;1)$};

% Sécantes pour h = 2, 1, 0.5
% h=2 -> M(3,9) hors cadre, on prend h=1.5 -> M(2.5, 6.25) hors cadre aussi. 
% Restons dans le cadre : h=1 -> M(2,4), h=0.5 -> M(1.5, 2.25), h=0.2 -> M(1.2, 1.44)
% On trace les droites complètes (sécantes)
% Equation sécante passant par A(1,1) et M(a+h, (a+h)^2) : pente = 2+h. y - 1 = (2+h)(x-1) -> y = (2+h)x -1 -h

% h = 1 (orange)
\addplot[domain=-0.5:3.5, samples=2, poporange, dashed, thick] {(2+1)*x - 1 - 1};
\addplot[only marks, mark=*, mark size=2, poporange] coordinates {(2,4)};
\node[poporange, anchor=south west] at (axis cs:2,4) {$M_1(2;4)$};
\node[poporange, anchor=north, font=\small] at (axis cs:2.5, 3.5) {$h=1$};

% h = 0.5 (green)
\addplot[domain=-0.5:3.5, samples=2, popgreen, dashed, thick] {(2+0.5)*x - 1 - 0.5};
\addplot[only marks, mark=*, mark size=2, popgreen] coordinates {(1.5,2.25)};
\node[popgreen, anchor=south west] at (axis cs:1.5,2.25) {$M_{0,5}$};
\node[popgreen, anchor=north, font=\small] at (axis cs:2, 2.5) {$h=0,5$};

% h = 0.2 (purple)
\addplot[domain=-0.5:3.5, samples=2, poppurple, dashed, thick] {(2+0.2)*x - 1 - 0.2};
\addplot[only marks, mark=*, mark size=2, poppurple] coordinates {(1.2,1.44)};
\node[poppurple, anchor=south west] at (axis cs:1.2,1.44) {$M_{0,2}$};

% Tangente en A : y = 2x - 1
\addplot[domain=-0.5:3.5, samples=2, poppink, very thick] {2*x - 1};
\node[poppink, anchor=south west, font=\small\bfseries] at (axis cs:2.5, 4) {Tangente ($h\to0$)};
\end{axis}
\end{tikzpicture}
\legende{Parabole $y=x^2$ : les sécantes $(AM_h)$ pour $h=1$, $h=0,5$, $h=0,2$ pivotent vers la tangente en $A(1;1)$ quand $h$ tend vers $0$.}
\end{popfigure}

\courssub{Passage à la limite : la tangente comme position limite de la sécante}

Sur la figure ci-dessus, on voit que lorsque $h$ devient très petit (le point $M$ se rapproche de $A$), la sécante $(AM)$ \textbf{pivote} autour de $A$ et se rapproche d'une position limite : c'est la \cle{tangente} à la courbe en $A$. La pente de cette tangente est la limite des pentes des sécantes.

\begin{definition}[Nombre dérivé d'une fonction en un point]
Soit $f$ une fonction définie sur un intervalle $I$ et $a$ un point intérieur à $I$. On dit que $f$ est \cle{dérivable en $a$} si le taux d'accroissement $\tau(h) = \dfrac{f(a+h)-f(a)}{h}$ admet une \textbf{limite finie} quand $h$ tend vers $0$.
Cette limite, notée $f'(a)$, est appelée le \textbf{nombre dérivé} de $f$ en $a$ :
\[
f'(a) = \lim_{h \to 0} \frac{f(a+h)-f(a)}{h}.
\]
Géométriquement, $f'(a)$ est le coefficient directeur de la tangente à la courbe $\mathcal{C}_f$ au point $A(a\,;\,f(a))$.
\end{definition}

\begin{aretenir}[Vocabulaire et notations]
\begin{itemize}
    \item Si la limite existe et est finie : $f$ est \textbf{dérivable en $a$}.
    \item Si la limite n'existe pas (forme indéterminée non levée, limite infinie, ou pas de limite) : $f$ n'est \textbf{pas dérivable en $a$}.
    \item Notations équivalentes pour le nombre dérivé : $f'(a)$, $\dfrac{df}{dx}(a)$, $Df(a)$.
    \item L'équation de la tangente en $A(a\,;\,f(a))$ sera étudiée en détail dans la section 2 : $y = f'(a)(x-a) + f(a)$.
\end{itemize}
\end{aretenir}

\courssub{Exemple de fonction non dérivable : la fonction valeur absolue}

La fonction $f(x) = |x|$ est continue en $0$ (son graphe n'a pas de « trou »), mais elle n'est \textbf{pas dérivable} en $0$. Pourquoi ? Parce que le graphe a un \textbf{point anguleux} (une « cassure ») en $0$.

\begin{popfigure}
\begin{tikzpicture}[scale=1.2, >=Stealth]
\begin{axis}[
    width=0.8\linewidth, height=7cm,
    axis lines=middle,
    xlabel=$x$, ylabel=$y$,
    xmin=-2.5, xmax=2.5,
    ymin=-0.5, ymax=2.5,
    xtick={-2,-1,1,2}, ytick={1,2},
    tick label style={font=\small},
    clip=false,
    domain=-2.5:2.5, samples=100,
]
% y = |x|
\addplot[popcurve, very thick, popblue] {abs(x)};
\node[popblue, anchor=south east] at (axis cs:2.3, 2.3) {$\mathcal{C}_f : y=|x|$};

% Point anguleux
\addplot[only marks, mark=*, mark size=3, popdark] coordinates {(0,0)};
\node[popdark, anchor=north east] at (axis cs:0,0) {$O(0;0)$};

% Demi-tangentes (droites directrices)
% Pour x>0 : y=x (pente 1)
\addplot[domain=0:2.5, samples=2, poporange, dashed, thick] {x};
\node[poporange, anchor=south west, font=\small] at (axis cs:2, 2) {Pente $+1$ (branche droite)};

% Pour x<0 : y=-x (pente -1)
\addplot[domain=-2.5:0, samples=2, popgreen, dashed, thick] {-x};
\node[popgreen, anchor=south east, font=\small] at (axis cs:-2, 2) {Pente $-1$ (branche gauche)};

% Flèches indiquant l'absence de tangente unique
\draw[poppink, very thick, <->] (axis cs:-0.8, 1.5) -- (axis cs:0.8, 1.5) node[midway, above, font=\small, poppink] {Pas de tangente unique};
\end{axis}
\end{tikzpicture}
\legende{La fonction $x \mapsto |x|$ admet deux demi-tangentes de pentes $-1$ et $+1$ en $0$, mais pas de tangente unique : elle n'est pas dérivable en $0$.}
\end{popfigure}

\begin{exoresolu}[Non-dérivabilité de $x \mapsto |x|$ en $0$]
Montrer que $f(x)=|x|$ n'est pas dérivable en $0$.
\tcblower
On calcule le taux d'accroissement $\tau(h) = \frac{|0+h|-|0|}{h} = \frac{|h|}{h}$.
\begin{itemize}
    \item Si $h > 0$ : $\tau(h) = \frac{h}{h} = 1$. La limite à droite est $1$.
    \item Si $h < 0$ : $\tau(h) = \frac{-h}{h} = -1$. La limite à gauche est $-1$.
\end{itemize}
Les limites à gauche et à droite sont différentes ($-1 \neq 1$), donc la limite quand $h \to 0$ n'existe pas. $f$ n'est pas dérivable en $0$.
\end{exoresolu}

\courssub{Méthode : calculer un nombre dérivé par la définition}

\begin{methode}[Calcul de $f'(a)$ par la définition (4 étapes)]
Pour calculer le nombre dérivé $f'(a)$ d'une fonction $f$ en un point $a$ :
\begin{enumerate}
    \item \textbf{Écrire} le taux d'accroissement : $\tau(h) = \dfrac{f(a+h)-f(a)}{h}$ ($h \neq 0$).
    \item \textbf{Simplifier} l'expression algébrique de $\tau(h)$ pour faire disparaître le $h$ au dénominateur (factorisation, mise au même dénominateur, conjugués, identités remarquables).
    \item \textbf{Faire tendre $h$ vers $0$} dans l'expression simplifiée (on remplace $h$ par $0$).
    \item \textbf{Conclure} : $f'(a) = \lim_{h\to 0} \tau(h) = \text{valeur trouvée}$.
\end{enumerate}
\end{methode}

\begin{attention}[Piège majeur : la forme indéterminée « $0/0$ »]
\emph{Ne jamais écrire} : $f'(a) = \frac{f(a+0)-f(a)}{0} = \frac{0}{0}$.
Le taux d'accroissement $\tau(h)$ n'est \textbf{défini que pour $h \neq 0$}. Il faut \textbf{obligatoirement simplifier} $\tau(h)$ \emph{avant} de faire tendre $h$ vers $0$. Si vous ne simplifiez pas, vous obtenez une forme indéterminée qui ne veut rien dire.
\end{attention}

\courssub{Calculs dirigés : exemples classiques au programme}

\begin{exoresolu}[Dérivé de $f(x)=x^2$ en $a=1$ (puis en $a$ quelconque)]
Calculer $f'(1)$ pour $f(x)=x^2$ par la définition.
\tcblower
\textbf{Étape 1 :} $\tau(h) = \frac{f(1+h)-f(1)}{h} = \frac{(1+h)^2 - 1^2}{h}$.
\textbf{Étape 2 :} Simplification.
\[
(1+h)^2 - 1 = 1 + 2h + h^2 - 1 = 2h + h^2 = h(2+h).
\]
Donc $\tau(h) = \frac{h(2+h)}{h} = 2+h$ \quad (valable pour $h \neq 0$).
\textbf{Étape 3 :} $\lim_{h\to 0} (2+h) = 2$.
\textbf{Étape 4 :} $f'(1) = 2$.
\medskip
\noindent \textbf{Généralisation en $a$ quelconque :}
$\tau(h) = \frac{(a+h)^2 - a^2}{h} = \frac{a^2+2ah+h^2-a^2}{h} = \frac{h(2a+h)}{h} = 2a+h \xrightarrow[h\to0]{} 2a$.
Donc \textbf{$f'(a) = 2a$} pour tout réel $a$.
\end{exoresolu}

\begin{exoresolu}[Dérivé de $f(x)=\frac{1}{x}$ en $a=2$ (puis en $a \neq 0$)]
Calculer $f'(2)$ pour $f(x)=\frac{1}{x}$ par la définition.
\tcblower
\textbf{Étape 1 :} $\tau(h) = \frac{f(2+h)-f(2)}{h} = \frac{\frac{1}{2+h} - \frac{1}{2}}{h}$.
\textbf{Étape 2 :} Simplification (mise au même dénominateur au numérateur).
\[
\frac{1}{2+h} - \frac{1}{2} = \frac{2 - (2+h)}{2(2+h)} = \frac{-h}{2(2+h)}.
\]
Donc $\tau(h) = \frac{\frac{-h}{2(2+h)}}{h} = \frac{-h}{2(2+h)} \times \frac{1}{h} = \frac{-1}{2(2+h)}$ \quad ($h \neq 0$, $h \neq -2$).
\textbf{Étape 3 :} $\lim_{h\to 0} \frac{-1}{2(2+h)} = \frac{-1}{2\times 2} = -\frac{1}{4}$.
\textbf{Étape 4 :} $f'(2) = -\frac{1}{4}$.
\medskip
\noindent \textbf{Généralisation en $a \neq 0$ :}
$\tau(h) = \frac{\frac{1}{a+h}-\frac{1}{a}}{h} = \frac{a-(a+h)}{a(a+h)} \times \frac{1}{h} = \frac{-h}{a(a+h)h} = \frac{-1}{a(a+h)} \xrightarrow[h\to0]{} -\frac{1}{a^2}$.
Donc \textbf{$f'(a) = -\frac{1}{a^2}$} pour tout $a \in \mathbb{R}^*$.
\end{exoresolu}

\begin{savaistu}[Naissance du calcul différentiel : Newton vs Leibniz]
Au XVII\textsuperscript{e} siècle, deux génies ont inventé le calcul différentiel \textbf{indépendamment} :
\begin{itemize}
    \item \textbf{Isaac Newton} (1643--1727), physicien anglais : il cherchait à décrire le mouvement (vitesse, accélération). Il appelait cela la « méthode des fluxions » ($\dot{x}$ pour la dérivée par rapport au temps).
    \item \textbf{Gottfried Wilhelm Leibniz} (1646--1716), philosophe et mathématicien allemand : il abordait le problème par la géométrie (tangentes, quadratures). Il a inventé la notation $\frac{dy}{dx}$ que nous utilisons encore aujourd'hui.
\end{itemize}
Une violente querelle de priorité a opposé les mathématiciens anglais (partisans de Newton) aux mathématiciens du continent (partisans de Leibniz) pendant des décennies. Aujourd'hui, on reconnaît les deux comme co-inventeurs. La notation de Leibniz ($f'(x)$, $\frac{df}{dx}$) a triomphé car elle est plus commode pour les calculs (règle de la chaîne, dérivées d'ordre supérieur).
\end{savaistu}

\begin{exercice}[Entraînement : calcul par la définition]
Calculer par la définition le nombre dérivé en $a$ des fonctions suivantes :
\begin{enumerate}
    \item $f(x) = 3x^2 - 2x$ en $a = 1$.
    \item $f(x) = \sqrt{x}$ en $a = 4$ (indication : utiliser le conjugué).
    \item $f(x) = \frac{1}{x+1}$ en $a = 0$.
\end{enumerate}
\end{exercice}

\begin{corrige}[Exercice n°1]
\begin{enumerate}
    \item $f(x)=3x^2-2x$, $a=1$.
    $\tau(h) = \frac{3(1+h)^2-2(1+h) - (3-2)}{h} = \frac{3(1+2h+h^2)-2-2h -1}{h} = \frac{3+6h+3h^2-3-2h}{h} = \frac{4h+3h^2}{h} = 4+3h \to 4$.
    $f'(1)=4$.
    \item $f(x)=\sqrt{x}$, $a=4$.
    $\tau(h) = \frac{\sqrt{4+h}-2}{h} = \frac{(\sqrt{4+h}-2)(\sqrt{4+h}+2)}{h(\sqrt{4+h}+2)} = \frac{4+h-4}{h(\sqrt{4+h}+2)} = \frac{h}{h(\sqrt{4+h}+2)} = \frac{1}{\sqrt{4+h}+2} \to \frac{1}{2+2} = \frac{1}{4}$.
    $f'(4)=\frac{1}{4}$.
    \item $f(x)=\frac{1}{x+1}$, $a=0$.
    $\tau(h) = \frac{\frac{1}{h+1}-1}{h} = \frac{1-(h+1)}{h(h+1)} = \frac{-h}{h(h+1)} = \frac{-1}{h+1} \to -1$.
    $f'(0)=-1$.
\end{enumerate}
\end{corrige}

\coursec{Tangente à une courbe en un point}

\courssub{De la sécante à la tangente : intuition géométrique}

Dans la section précédente, nous avons découvert le \cle{nombre dérivé} $f'(a)$ comme la limite du taux d'accroissement $\frac{f(x)-f(a)}{x-a}$ lorsque $x$ tend vers $a$. Géométriquement, ce taux d'accroissement est le \cle{coefficient directeur} de la \cle{sécante} $(AB)$ joignant les points $A(a\,;\,f(a))$ et $B(x\,;\,f(x))$ de la courbe $\mathcal{C}_f$.

\begin{popfigure}
\begin{tikzpicture}[scale=1.1, every node/.style={font=\small}]
  % Axes
  \draw[->] (-0.5,0) -- (4.5,0) node[right] {$x$};
  \draw[->] (0,-1.5) -- (0,4.5) node[above] {$y$};
  % Courbe f(x) = x^2 - 2x
  \draw[popblue, thick, domain=-0.2:4.2, samples=100] plot (\x, {\x*\x - 2*\x});
  % Point A(2,0)
  \coordinate (A) at (2,0);
  \fill[poporange] (A) circle (2pt) node[below right] {$A(2\,;\,0)$};
  % Sécantes pour x = 3, 2.5, 2.2, 2.1
  \foreach \xval/\col in {3/popgreen, 2.5/poppurple, 2.2/popink, 2.1/popgold}{
    \coordinate (B) at (\xval, {\xval*\xval - 2*\xval});
    \draw[\col, dashed] (A) -- (B);
    \fill[\col] (B) circle (1.5pt);
  }
  % Tangente y = 2x - 4
  \draw[popdark, thick, domain=0.5:3.5] plot (\x, {2*\x - 4});
  \node[popdark] at (3.8,3.2) {Tangente};
  % Légende
  \node[align=left, anchor=north west] at (-0.3,4.2) {
    \textcolor{popblue}{\rule{1em}{1pt}} $f(x)=x^2-2x$\par
    \textcolor{popgreen}{\rule{1em}{1pt}} Sécante $x=3$\par
    \textcolor{poppurple}{\rule{1em}{1pt}} Sécante $x=2,5$\par
    \textcolor{popink}{\rule{1em}{1pt}} Sécante $x=2,2$\par
    \textcolor{popgold}{\rule{1em}{1pt}} Sécante $x=2,1$\par
    \textcolor{popdark}{\rule{1em}{1pt}} Tangente en $A$
  };
\end{tikzpicture}
\legende{Approche de la tangente par les sécantes : quand $B$ se rapproche de $A$, la sécante $(AB)$ tend vers une position limite, la tangente.}
\end{popfigure}

\begin{definition}[Tangente à une courbe en un point]
Soit $f$ une fonction dérivable en $a$ et $\mathcal{C}_f$ sa courbe représentative dans un repère orthonormé. La \cle{tangente} à $\mathcal{C}_f$ au point $A(a\,;\,f(a))$ est la droite passant par $A$ dont le coefficient directeur est le nombre dérivé $f'(a)$. C'est la position limite des sécantes $(AB)$ lorsque le point $B$ tend vers $A$ sur la courbe.
\end{definition}

\courssub{Équation de la tangente : théorème et démonstration}

\begin{propriete}[Équation de la tangente — Théorème fondamental]
Soit $f$ une fonction dérivable en $a$. La tangente $\mathcal{T}_A$ à la courbe $\mathcal{C}_f$ au point $A(a\,;\,f(a))$ a pour équation cartésienne :
\[
\boxed{y = f'(a)(x - a) + f(a)}
\]
\end{propriete}

\begin{experience}[Démonstration guidée : pourquoi cette formule ?]
\emph{Objectif : retrouver l'équation de la tangente en utilisant seulement la définition de la droite passant par un point avec un coefficient directeur connu.}

\begin{enumerate}
  \item \textbf{Ce que l'on sait :} La tangente $\mathcal{T}_A$ passe par $A(a\,;\,f(a))$ et son coefficient directeur est $m = f'(a)$.
  \item \textbf{Équation réduite d'une droite :} Toute droite de coefficient directeur $m$ passant par un point $(x_0,y_0)$ a pour équation $y - y_0 = m(x - x_0)$.
  \item \textbf{Application :} On remplace $x_0$ par $a$, $y_0$ par $f(a)$ et $m$ par $f'(a)$ :
        \[ y - f(a) = f'(a)(x - a) \]
  \item \textbf{Forme explicite :} On isole $y$ :
        \[ y = f'(a)(x - a) + f(a) \]
\end{enumerate}
\emph{Conclusion :} La formule n'est rien d'autre que l'équation cartésienne d'une droite dont on connaît un point et la pente. La nouveauté, c'est que cette pente \textbf{est le nombre dérivé} $f'(a)$.
\end{experience}

\begin{aretenir}[Lecture graphique du nombre dérivé]
Si la tangente $\mathcal{T}_A$ est tracée sur le graphique, on peut \textbf{lire graphiquement} $f'(a)$ comme le coefficient directeur de cette tangente :
\[
f'(a) = \frac{\text{variation de } y}{\text{variation de } x} \text{ sur la tangente}
\]
On choisit deux points lisibles $M$ et $N$ sur la tangente, on repère leurs coordonnées et on calcule $\frac{y_N - y_M}{x_N - x_M}$.
\end{aretenir}

\begin{aretenir}[Formules de dérivation usuelles (rappel pour la classe de Première technique)]
Pour toute fonction dérivable sur son domaine :
\begin{itemize}
  \item $f(x) = k \quad (k \in \mathbb{R}) \quad \Rightarrow \quad f'(x) = 0$
  \item $f(x) = x^n \quad (n \in \mathbb{N}^*) \quad \Rightarrow \quad f'(x) = n x^{n-1}$
  \item $f(x) = \sqrt{x} = x^{1/2} \quad (x > 0) \quad \Rightarrow \quad f'(x) = \frac{1}{2\sqrt{x}}$
  \item $f(x) = \frac{1}{x} = x^{-1} \quad (x \neq 0) \quad \Rightarrow \quad f'(x) = -\frac{1}{x^2}$
  \item $(u+v)' = u' + v'$ \quad ; \quad $(\lambda u)' = \lambda u'$ \quad ; \quad $(uv)' = u'v + uv'$
  \item $\left(\frac{u}{v}\right)' = \frac{u'v - uv'}{v^2} \quad (v \neq 0)$
\end{itemize}
\end{aretenir}

\courssub{Méthode systématique : déterminer l'équation d'une tangente}

\begin{methode}[Trouver l'équation de la tangente en un point d'abscisse $a$]
\begin{enumerate}
  \item \textbf{Calculer l'ordonnée du point de contact :} $f(a)$.
  \item \textbf{Calculer la fonction dérivée $f'(x)$} (à l'aide des formules usuelles), \textbf{puis sa valeur en $a$ :} $f'(a)$.
  \item \textbf{Écrire l'équation de la tangente} avec la formule : $y = f'(a)(x - a) + f(a)$.
  \item \textbf{Réduire l'expression} pour obtenir la forme $y = mx + p$ (coefficient directeur $m$ et ordonnée à l'origine $p$).
\end{enumerate}
\end{methode}

\courssub{Exemple chiffré détaillé}

\begin{exemplebox}[Tangente à la courbe de $f(x)=x^2-2x$ au point d'abscisse $a=2$]
Soit $f(x)=x^2-2x$ définie sur $\mathbb{R}$. Déterminer l'équation de la tangente $\mathcal{T}$ à $\mathcal{C}_f$ au point $A$ d'abscisse $2$.

\textbf{Étape 1 :} $f(2) = 2^2 - 2\times 2 = 4 - 4 = 0$. Le point de contact est $A(2\,;\,0)$.

\textbf{Étape 2 :} $f'(x) = 2x - 2$ (dérivée de $x^2$ est $2x$, dérivée de $-2x$ est $-2$).
Donc $f'(2) = 2\times 2 - 2 = 2$. Le coefficient directeur de la tangente est $2$.

\textbf{Étape 3 :} Équation de la tangente :
\[ y = f'(2)(x - 2) + f(2) = 2(x - 2) + 0 = 2x - 4 \]

\textbf{Étape 4 :} La forme réduite est déjà obtenue : $\boxed{y = 2x - 4}$.
\end{exemplebox}

\begin{popfigure}
\begin{tikzpicture}[scale=1.1, every node/.style={font=\small}]
  % Axes
  \draw[->] (-0.5,0) -- (4.5,0) node[right] {$x$};
  \draw[->] (0,-1.5) -- (0,4.5) node[above] {$y$};
  % Courbe f(x) = x^2 - 2x
  \draw[popblue, thick, domain=-0.2:4.2, samples=100] plot (\x, {\x*\x - 2*\x});
  \node[popblue] at (4.3, 3.5) {$\mathcal{C}_f : y=x^2-2x$};
  % Point A(2,0)
  \coordinate (A) at (2,0);
  \fill[poporange] (A) circle (2.5pt) node[below right] {$A(2\,;\,0)$};
  % Tangente y = 2x - 4
  \draw[poporange, thick, domain=0.5:3.5] plot (\x, {2*\x - 4});
  \node[poporange] at (3.8, 3.0) {$\mathcal{T}_A : y=2x-4$};
  % Triangle de pente sur la tangente
  \coordinate (M) at (1, -2);
  \coordinate (N) at (3, 2);
  \draw[popdark, dashed] (M) -- (N) -- (3, -2) -- cycle;
  \node[popdark, left] at (1, -2) {$M(1\,;\,-2)$};
  \node[popdark, right] at (3, 2) {$N(3\,;\,2)$};
  \node[popdark, below] at (2, -2) {$\Delta x = 2$};
  \node[popdark, right] at (3, 0) {$\Delta y = 4$};
  % Annotation f'(a)
  \node[popdark, align=center] at (0.8, 3.5) {$f'(2) = \frac{\Delta y}{\Delta x} = \frac{4}{2} = 2$};
\end{tikzpicture}
\legende{Tangente à $f(x)=x^2-2x$ en $A(2;0)$ : lecture graphique du coefficient directeur $f'(2)=2$.}
\end{popfigure}

\courssub{Cas particuliers et applications}

\begin{propriete}[Tangente horizontale]
La tangente à $\mathcal{C}_f$ au point $A(a\,;\,f(a))$ est \textbf{horizontale} (parallèle à l'axe des abscisses) \textbf{si et seulement si} $f'(a) = 0$.
Son équation est alors simplement $y = f(a)$.
\end{propriete}

\begin{popfigure}
\begin{tikzpicture}[scale=1.1, every node/.style={font=\small}]
  % Axes
  \draw[->] (-0.5,0) -- (5.5,0) node[right] {$x$};
  \draw[->] (0,-2.5) -- (0,3.5) node[above] {$y$};
  % Courbe f(x) = x^3 - 3x (exemple avec deux extrema)
  \draw[popblue, thick, domain=-2.2:2.2, samples=100] plot (\x, {\x*\x*\x - 3*\x});
  \node[popblue] at (5.2, 2.5) {$\mathcal{C}_f : y=x^3-3x$};
  % Points où f'(x)=0 : x = -1 et x = 1
  \coordinate (A) at (-1, 2);
  \coordinate (B) at (1, -2);
  \fill[poporange] (A) circle (2.5pt) node[above left] {$A(-1\,;\,2)$};
  \fill[poporange] (B) circle (2.5pt) node[below right] {$B(1\,;\,-2)$};
  % Tangentes horizontales
  \draw[poporange, dashed, thick] (-2.2, 2) -- (2.2, 2);
  \draw[poporange, dashed, thick] (-2.2, -2) -- (2.2, -2);
  \node[poporange] at (2.5, 2) {$y=2$};
  \node[poporange] at (2.5, -2) {$y=-2$};
  % Annotation f'(a)=0
  \node[popdark, align=center] at (-2.2, 3) {$f'(-1)=0 \Rightarrow$ tangente horizontale};
  \node[popdark, align=center] at (-2.2, -2.8) {$f'(1)=0 \Rightarrow$ tangente horizontale};
\end{tikzpicture}
\legende{Courbe admettant deux tangentes horizontales aux points où $f'(a)=0$ (extrema locaux).}
\end{popfigure}

\begin{exemplebox}[Tangente parallèle à une droite donnée]
Soit $f(x)=x^3-3x$. Trouver le(s) point(s) de $\mathcal{C}_f$ où la tangente est parallèle à la droite $\mathcal{D} : y = 6x + 1$.

\textbf{Raisonnement :} Deux droites sont parallèles $\iff$ elles ont le même coefficient directeur.
La droite $\mathcal{D}$ a pour coefficient directeur $6$. On cherche donc $a$ tel que $f'(a) = 6$.

$f'(x) = 3x^2 - 3$.
On résout $3a^2 - 3 = 6 \iff 3a^2 = 9 \iff a^2 = 3 \iff a = \sqrt{3}$ ou $a = -\sqrt{3}$.

Les points sont :
$A(\sqrt{3}\,;\, f(\sqrt{3})) = (\sqrt{3}\,;\, 3\sqrt{3} - 3\sqrt{3}) = (\sqrt{3}\,;\, 0)$
$B(-\sqrt{3}\,;\, f(-\sqrt{3})) = (-\sqrt{3}\,;\, -3\sqrt{3} + 3\sqrt{3}) = (-\sqrt{3}\,;\, 0)$

Les équations des tangentes :
En $A$ : $y = 6(x - \sqrt{3}) + 0 = 6x - 6\sqrt{3}$
En $B$ : $y = 6(x + \sqrt{3}) + 0 = 6x + 6\sqrt{3}$
\end{exemplebox}

\courssub{Exercices résolus et pièges à éviter}

\begin{exoresolu}[Tangente à une fonction racine]
Déterminer l'équation de la tangente à la courbe de $f(x)=\sqrt{x}$ au point d'abscisse $a=4$.

\tcblower
\textbf{Solution :}
\begin{enumerate}
  \item $f(4) = \sqrt{4} = 2$. Point $A(4\,;\,2)$.
  \item $f(x) = x^{1/2} \Rightarrow f'(x) = \frac{1}{2}x^{-1/2} = \frac{1}{2\sqrt{x}}$.
        $f'(4) = \frac{1}{2\sqrt{4}} = \frac{1}{4}$.
  \item Équation : $y = \frac{1}{4}(x - 4) + 2 = \frac{1}{4}x - 1 + 2 = \frac{1}{4}x + 1$.
\end{enumerate}
$\boxed{y = \frac{1}{4}x + 1}$
\end{exoresolu}

\begin{attention}[Erreurs fréquentes — Pièges du Probatoire]
\begin{enumerate}
  \item \textbf{Confondre $f'(a)$ et $f'(x)$ :} $f'(a)$ est un \textbf{nombre} (la pente en $a$), $f'(x)$ est une \textbf{fonction}. Dans la formule $y = f'(a)(x-a)+f(a)$, on utilise le \textbf{nombre} $f'(a)$.
  \item \textbf{Oublier le terme $(x-a)$ :} Écrire $y = f'(a)x + f(a)$ est \textbf{faux} (sauf si $a=0$). La tangente passe par $A(a\,;\,f(a))$, pas par $(0\,;\,f(a))$.
  \item \textbf{Calculer $f'(x)$ mais ne pas l'évaluer en $a$ :} Il faut remplacer $x$ par $a$ dans $f'(x)$ pour obtenir le coefficient directeur.
  \item \textbf{Ne pas réduire l'équation finale :} Au Probatoire, on demande souvent la forme réduite $y = mx + p$. Pensez à développer et simplifier.
  \item \textbf{Domaine de définition :} Vérifiez que $a$ appartient bien au domaine de définition de $f$ \textbf{et} de $f'$ (ex. $f(x)=\sqrt{x}$ n'est pas dérivable en $0$).
\end{enumerate}
\end{attention}

\begin{exercice}[Entraînement : tangentes et parallélisme]
Soit $f(x) = \dfrac{1}{x}$ définie sur $\mathbb{R}^*$.
\begin{enumerate}
  \item Déterminer l'équation de la tangente $\mathcal{T}_1$ à $\mathcal{C}_f$ au point d'abscisse $1$.
  \item Déterminer l'équation de la tangente $\mathcal{T}_2$ à $\mathcal{C}_f$ au point d'abscisse $-2$.
  \item Ces deux tangentes sont-elles parallèles ? Justifier.
  \item Trouver le(s) point(s) de $\mathcal{C}_f$ où la tangente est parallèle à la droite $y = -4x + 5$.
\end{enumerate}
\end{exercice}

\begin{corrige}[Exercice : Tangentes à $f(x)=1/x$]
\begin{enumerate}
  \item $f(1)=1$. $f'(x)=-\frac{1}{x^2} \Rightarrow f'(1)=-1$.
        $\mathcal{T}_1 : y = -1(x-1)+1 = -x+2$.
  \item $f(-2)=-\frac{1}{2}$. $f'(-2)=-\frac{1}{4}$.
        $\mathcal{T}_2 : y = -\frac{1}{4}(x+2)-\frac{1}{2} = -\frac{1}{4}x - \frac{1}{2} - \frac{1}{2} = -\frac{1}{4}x - 1$.
  \item Coefficients directeurs : $-1$ et $-\frac{1}{4}$. Différents $\Rightarrow$ \textbf{non parallèles}.
  \item On cherche $a$ tel que $f'(a) = -4 \iff -\frac{1}{a^2} = -4 \iff a^2 = \frac{1}{4} \iff a = \frac{1}{2}$ ou $a = -\frac{1}{2}$.
        Points : $(\frac{1}{2}\,;\,2)$ et $(-\frac{1}{2}\,;\,-2)$.
        Tangentes : $y = -4x + 4$ et $y = -4x - 4$.
\end{enumerate}
\end{corrige}

\courssub{Applications concrètes : métiers techniques}

\begin{savaistu}[La tangente sur le terrain : pente de route, profil en long]
Dans les travaux publics (BTP), le \emph{profil en long} d'une route est la courbe donnant l'altitude $z$ en fonction de l'abscisse curviligne $x$. La \textbf{pente} (en \%) en un point est exactement $100 \times f'(x)$ où $f(x)$ est l'altitude.
\begin{itemize}
  \item Une tangente horizontale ($f'(x)=0$) correspond à un \textbf{sommet} ou un \textbf{creux} de la route (point haut ou point bas).
  \item La tangente à la courbe en un point donne la \textbf{droite de plus grande pente} locale : c'est la direction qu'une bille suivrait si on la posait sur la route à cet endroit.
  \item Pour un électricien : la courbe $I(U)$ d'une diode (caractéristique courant-tension). La tangente en un point de fonctionnement donne la \textbf{résistance dynamique} $r_d = \frac{1}{f'(U)}$ (inverse de la pente), paramètre crucial pour les petits signaux.
\end{itemize}
\end{savaistu}

\begin{methode}[Résumé : comment aborder un exercice de tangente au Probatoire]
\begin{enumerate}
  \item \textbf{Lire l'énoncé :} identifier $f$, $a$, et ce qu'on demande (équation, abscisse(s) pour tangente parallèle/horizontale, etc.).
  \item \textbf{Vérifier la dérivabilité :} $a \in D_f$ et $f$ dérivable en $a$.
  \item \textbf{Appliquer la méthode en 4 étapes} (calculer $f(a)$, $f'(x)$, $f'(a)$, écrire et réduire).
  \item \textbf{Pour "tangente parallèle à..." :} égaler $f'(a)$ au coefficient directeur de la droite donnée, résoudre pour $a$, puis faire l'étape 1 à 4 pour chaque $a$ trouvé.
  \item \textbf{Pour "tangente horizontale" :} résoudre $f'(a)=0$.
  \item \textbf{Relire :} l'équation finale est-elle sous la forme $y=mx+p$ ? Le point $A$ appartient-il bien à la tangente ? (Test rapide : remplacer $x$ par $a$, on doit retrouver $f(a)$).
\end{enumerate}
\end{methode}

\coursec{Fonction dérivée et dérivées des fonctions usuelles}

Dans la section précédente, nous avons appris à déterminer le \cle{nombre dérivé} $f'(a)$ d'une fonction en un point $a$, et à l'interpréter comme le coefficient directeur de la tangente à la courbe au point d'abscisse $a$. Mais refaire le calcul de la limite du taux d'accroissement à chaque point serait fastidieux. L'idée géniale de cette section est la suivante : au lieu de calculer $f'(a)$ pour une valeur fixée de $a$, on calcule $f'(x)$ pour \emph{n'importe quel} $x$. On obtient ainsi une nouvelle fonction, la \cle{fonction dérivée}, qui donne d'un seul coup tous les nombres dérivés. C'est exactement ce que fait un mécanicien qui, au lieu de mesurer la vitesse d'un véhicule à un instant précis, installe un compteur qui affiche la vitesse à \emph{tout} instant.

\courssub{De la notion ponctuelle à la fonction dérivée}

\begin{definition}[Fonction dérivée]
Soit $f$ une fonction définie sur un intervalle $I$. On dit que $f$ est \cle{dérivable sur $I$} lorsque $f$ est dérivable en tout point $x$ de $I$, c'est-à-dire lorsque le taux d'accroissement
\[
\tau(h) = \frac{f(x+h)-f(x)}{h}
\]
admet une limite finie quand $h$ tend vers $0$, pour tout $x$ de $I$.
Dans ce cas, la fonction qui à tout $x$ de $I$ associe le nombre dérivé $f'(x)$ est appelée \cle{fonction dérivée} de $f$ sur $I$. On la note $f'$ :
\[
f' : x \longmapsto f'(x).
\]
\end{definition}

\begin{aretenir}[À retenir]
Le \cle{nombre dérivé} $f'(a)$ est un \emph{nombre} (le coefficient directeur de la tangente au point d'abscisse $a$). La \cle{fonction dérivée} $f'$ est une \emph{fonction} : elle permet de calculer tous les nombres dérivés d'un coup, par simple remplacement de $x$ par la valeur souhaitée.
\end{aretenir}

\begin{exemplebox}[Une image concrète]
Un chauffeur de taxi-moto à Douala parcourt une distance $f(t)$ (en km) en fonction du temps $t$ (en h). La fonction dérivée $f'$ est sa \emph{vitesse instantanée} : $f'(t)$ donne la vitesse affichée au compteur à l'instant $t$. Connaître la fonction $f'$ revient à connaître la vitesse à chaque instant du trajet, sans refaire de calcul.
\end{exemplebox}

\courssub{Les dérivées des fonctions usuelles}

Grâce au calcul de la limite du taux d'accroissement, on établit une fois pour toutes les dérivées des fonctions de référence. Ces résultats sont à connaître \textbf{par c\oe ur} : ils sont la boîte à outils de tout le chapitre.

\begin{propriete}[Dérivées des fonctions usuelles — à connaître par c\oe ur]
\begin{center}
\begin{tabularx}{\linewidth}{|l|X|X|}
\hline
\rowcolor{popblueL}
Fonction $f$ & Fonction dérivée $f'$ & $f$ dérivable sur \\
\hline
$f(x) = k$ (constante) & $f'(x) = 0$ & $\mathbb{R}$ \\
\hline
$f(x) = x$ & $f'(x) = 1$ & $\mathbb{R}$ \\
\hline
$f(x) = x^2$ & $f'(x) = 2x$ & $\mathbb{R}$ \\
\hline
$f(x) = x^3$ & $f'(x) = 3x^2$ & $\mathbb{R}$ \\
\hline
$f(x) = x^n$ \ ($n$ entier $\geq 1$) & $f'(x) = n\,x^{n-1}$ & $\mathbb{R}$ \\
\hline
$f(x) = \dfrac{1}{x}$ & $f'(x) = -\dfrac{1}{x^2}$ & $]-\infty\,;\,0[$ et $]0\,;\,+\infty[$ \\
\hline
$f(x) = \sqrt{x}$ & $f'(x) = \dfrac{1}{2\sqrt{x}}$ & $]0\,;\,+\infty[$ \\
\hline
\end{tabularx}
\end{center}
\end{propriete}

\begin{attention}[Erreurs fréquentes]
\begin{itemize}
\item \textbf{Signe oublié} : $(1/x)' = -\dfrac{1}{x^2}$, et non $\dfrac{1}{x^2}$. Le signe moins est la faute la plus fréquente au Probatoire.
\item \textbf{Ensemble de dérivabilité} : la formule $(\sqrt{x})' = \dfrac{1}{2\sqrt{x}}$ n'a de sens que pour $x > 0$. Écrire $f'(0)$ pour $f(x)=\sqrt{x}$ est une erreur : $\sqrt{x}$ n'est \emph{pas} dérivable en $0$ (voir plus loin).
\item \textbf{Exposant} : $(x^n)' = n\,x^{n-1}$ : on \emph{descend} l'exposant en facteur, puis on le \emph{diminue de 1}. Ainsi $(x^5)' = 5x^4$, et non $5x^5$ ni $x^4$.
\end{itemize}
\end{attention}

\courssub{Démonstrations des formules usuelles}

Certaines formules se démontrent directement à partir de la définition du nombre dérivé. Ces démonstrations sont formatrices : elles montrent que les formules du tableau ne tombent pas du ciel.

\begin{experience}[Démonstration guidée : dérivée de $f(x) = x^2$]
Nous avons déjà rencontré ce calcul dans la section 1 ; reprenons-le méthodiquement. Soit $f(x) = x^2$ et $x$ un réel quelconque.
\begin{enumerate}
\item On forme le taux d'accroissement :
\[
\tau(h) = \frac{f(x+h)-f(x)}{h} = \frac{(x+h)^2 - x^2}{h}.
\]
\item On développe le numérateur : $(x+h)^2 = x^2 + 2xh + h^2$, donc
\[
\tau(h) = \frac{x^2 + 2xh + h^2 - x^2}{h} = \frac{2xh + h^2}{h}.
\]
\item On simplifie par $h$ (licite car $h \neq 0$) : $\tau(h) = 2x + h$.
\item On fait tendre $h$ vers $0$ : $\tau(h) \to 2x$.
\end{enumerate}
Conclusion : pour tout réel $x$, $f'(x) = 2x$. La fonction carré est dérivable sur $\mathbb{R}$ et $(x^2)' = 2x$.
\end{experience}

\begin{experience}[Démonstration guidée : dérivée de $f(x) = \dfrac{1}{x}$]
Soit $f(x) = \dfrac{1}{x}$, définie pour $x \neq 0$.
\begin{enumerate}
\item Taux d'accroissement en $x \neq 0$ :
\[
\tau(h) = \frac{\dfrac{1}{x+h} - \dfrac{1}{x}}{h}.
\]
\item On réduit le numérateur au même dénominateur :
\[
\frac{1}{x+h} - \frac{1}{x} = \frac{x - (x+h)}{x(x+h)} = \frac{-h}{x(x+h)}.
\]
\item On divise par $h$ :
\[
\tau(h) = \frac{-h}{h\,x(x+h)} = \frac{-1}{x(x+h)}.
\]
\item Quand $h \to 0$, $x + h \to x$, donc $\tau(h) \to -\dfrac{1}{x^2}$.
\end{enumerate}
Conclusion : pour tout $x \neq 0$, $\left(\dfrac{1}{x}\right)' = -\dfrac{1}{x^2}$. Remarquez que cette dérivée est toujours \emph{négative} : la fonction inverse est décroissante sur chacun des intervalles de son domaine, ce que confirme sa courbe (hyperbole).
\end{experience}

\begin{propriete}[Dérivée de la fonction racine carrée — démonstration exigible]
Soit $f(x) = \sqrt{x}$. Alors $f$ est dérivable sur $]0\,;\,+\infty[$ et
\[
f'(x) = \frac{1}{2\sqrt{x}} \quad \text{pour tout } x > 0.
\]
\end{propriete}

\begin{experience}[Démonstration : la ruse de la quantité conjuguée]
Soit $x > 0$ fixé. Le taux d'accroissement est
\[
\tau(h) = \frac{\sqrt{x+h} - \sqrt{x}}{h}.
\]
Problème : si l'on fait $h \to 0$, on obtient la forme indéterminée $\frac{0}{0}$. L'astuce consiste à multiplier en haut et en bas par la \cle{quantité conjuguée} $\sqrt{x+h} + \sqrt{x}$, afin d'utiliser l'identité $(a-b)(a+b) = a^2 - b^2$ :
\begin{align*}
\tau(h) &= \frac{\big(\sqrt{x+h} - \sqrt{x}\big)\big(\sqrt{x+h} + \sqrt{x}\big)}{h\big(\sqrt{x+h} + \sqrt{x}\big)} \\
&= \frac{(x+h) - x}{h\big(\sqrt{x+h} + \sqrt{x}\big)} \\
&= \frac{h}{h\big(\sqrt{x+h} + \sqrt{x}\big)} \\
&= \frac{1}{\sqrt{x+h} + \sqrt{x}}.
\end{align*}
Quand $h \to 0$, $\sqrt{x+h} \to \sqrt{x}$, donc le dénominateur tend vers $2\sqrt{x}$ (qui est non nul car $x > 0$). Ainsi :
\[
\tau(h) \longrightarrow \frac{1}{2\sqrt{x}}.
\]
Conclusion : $f$ est dérivable en tout $x > 0$ et $f'(x) = \dfrac{1}{2\sqrt{x}}$.
\end{experience}

\begin{attention}[Pourquoi exclure $x = 0$ ?]
Dans la démonstration, le dénominateur $2\sqrt{x}$ doit être non nul : il faut donc $x > 0$. En $0$, le taux d'accroissement vaut
\[
\tau(h) = \frac{\sqrt{h} - 0}{h} = \frac{1}{\sqrt{h}} \quad (h > 0),
\]
qui tend vers $+\infty$ quand $h \to 0^+$ : la limite n'est pas finie. La fonction racine carrée \textbf{n'est pas dérivable en 0}. Graphiquement, sa courbe admet en l'origine une \cle{demi-tangente verticale} : la pente devient infiniment grande.
\end{attention}

\begin{popfigure}
\begin{center}
\begin{tikzpicture}
\begin{axis}[
  width=0.85\linewidth, height=6.5cm,
  axis lines=middle,
  xlabel={$x$}, ylabel={$y$},
  xmin=-0.4, xmax=4.4, ymin=-0.5, ymax=2.5,
  xtick={1,2,3,4}, ytick={1,2},
  samples=200, thick
]
\addplot[popcurve,domain=0:4.2]{sqrt(max(0,x))};
\addplot[poporange,very thick,domain=0:0.02] coordinates {(0,0) (0,2.3)};
\node[poporange,anchor=west] at (axis cs:0.12,1.9) {\small demi-tangente verticale};
\node[popblue,anchor=north west] at (axis cs:2.2,1.2) {\small $y = \sqrt{x}$};
\fill[popdark] (axis cs:0,0) circle (1.6pt);
\node[anchor=north east,popdark] at (axis cs:0,0) {\small $O$};
\end{axis}
\end{tikzpicture}
\end{center}
\legende{La courbe de $x \mapsto \sqrt{x}$ démarre à l'origine avec une demi-tangente verticale : la pente y est infinie, donc la fonction n'est pas dérivable en $0$.}
\end{popfigure}

\begin{aretenir}[Formule admise : dérivée de $x^n$]
Pour tout entier $n \geq 1$ et tout réel $x$ :
\[
\left(x^n\right)' = n\,x^{n-1}.
\]
Cette formule généralise $(x)' = 1$, $(x^2)' = 2x$, $(x^3)' = 3x^2$. Elle se démontre par récurrence sur $n$ ; cette démonstration est \textbf{hors exigibilité} en Première technique : la formule est admise et doit simplement être sue et bien appliquée.
\end{aretenir}

\courssub{Entraînement : calculer des nombres dérivés à l'aide des formules}

\begin{exemplebox}[Exemple chiffré détaillé]
Soit $f(x) = x^3$. D'après la formule $(x^n)' = n\,x^{n-1}$ avec $n = 3$ :
\[
f'(x) = 3x^2.
\]
On en déduit le nombre dérivé en $-2$ :
\[
f'(-2) = 3 \times (-2)^2 = 3 \times 4 = 12.
\]
Interprétation géométrique : la tangente à la courbe de $f$ au point d'abscisse $-2$ a pour \cle{coefficient directeur} $12$. La tangente est très « pentue » : elle monte vite.
\end{exemplebox}

\begin{exoresolu}[Calculs directs de dérivées]
\textbf{Énoncé.} Pour chaque fonction, donner la fonction dérivée en précisant l'ensemble de dérivabilité, puis calculer le nombre dérivé demandé.
\begin{enumerate}
\item $f(x) = x^5$ ; calculer $f'(1)$ et $f'(-1)$.
\item $g(x) = \dfrac{1}{x}$ ; calculer $g'(2)$.
\item $h(x) = \sqrt{x}$ ; calculer $h'(9)$.
\end{enumerate}
\tcblower
\textbf{Solution détaillée.}
\begin{enumerate}
\item $f(x) = x^5$ est dérivable sur $\mathbb{R}$ et $f'(x) = 5x^4$.
\[
f'(1) = 5 \times 1^4 = 5 \quad ; \quad f'(-1) = 5 \times (-1)^4 = 5.
\]
\item $g(x) = \dfrac{1}{x}$ est dérivable sur $]-\infty\,;\,0[$ et sur $]0\,;\,+\infty[$, et $g'(x) = -\dfrac{1}{x^2}$.
\[
g'(2) = -\frac{1}{2^2} = -\frac{1}{4}.
\]
La tangente à l'hyperbole au point d'abscisse $2$ descend (coefficient directeur négatif).
\item $h(x) = \sqrt{x}$ est dérivable sur $]0\,;\,+\infty[$ et $h'(x) = \dfrac{1}{2\sqrt{x}}$.
\[
h'(9) = \frac{1}{2\sqrt{9}} = \frac{1}{2 \times 3} = \frac{1}{6}.
\]
\end{enumerate}
\end{exoresolu}

\begin{exercice}[À vous de jouer]
\begin{enumerate}
\item Soit $f(x) = x^4$. Déterminer $f'(x)$, puis calculer $f'(2)$ et $f'(0)$.
\item Soit $g(x) = \dfrac{1}{x}$. Calculer $g'(-3)$. Que vaut le coefficient directeur de la tangente à la courbe de $g$ au point d'abscisse $-3$ ?
\item Soit $h(x) = \sqrt{x}$. Calculer $h'(4)$ et $h'(0{,}25)$. La fonction $h$ admet-elle un nombre dérivé en $0$ ? Justifier.
\item Un technicien mesure la pente d'une toiture modélisée par $f(x) = x^2$ pour $x \in [0\,;\,5]$. Quelle est la pente de la tangente en $x = 3$ ?
\end{enumerate}
\end{exercice}

\begin{corrige}[Exercice « À vous de jouer »]
\begin{enumerate}
\item $f'(x) = 4x^3$ sur $\mathbb{R}$. $f'(2) = 4 \times 8 = 32$ ; $f'(0) = 0$ (tangente horizontale à l'origine).
\item $g'(x) = -\dfrac{1}{x^2}$, donc $g'(-3) = -\dfrac{1}{9}$. Le coefficient directeur de la tangente en $-3$ est $-\dfrac{1}{9}$.
\item $h'(x) = \dfrac{1}{2\sqrt{x}}$ pour $x > 0$. $h'(4) = \dfrac{1}{2 \times 2} = \dfrac{1}{4}$ ; $h'(0{,}25) = \dfrac{1}{2 \times 0{,}5} = 1$. En $0$, le taux d'accroissement $\dfrac{1}{\sqrt{h}}$ tend vers $+\infty$ : pas de limite finie, donc $h$ n'est pas dérivable en $0$ (demi-tangente verticale).
\item $f'(x) = 2x$, donc $f'(3) = 6$ : la pente de la tangente en $x = 3$ vaut $6$.
\end{enumerate}
\end{corrige}

\begin{savaistu}[Trois notations pour une même idée]
La notation $f'$ que nous utilisons est due au mathématicien français \textbf{Joseph-Louis Lagrange} (1736-1813). Mais ses illustres prédécesseurs avaient chacun leur écriture : \textbf{Isaac Newton}, qui pensait la dérivée comme une vitesse, notait $\dot{x}$ (un point sur la lettre) ; \textbf{Gottfried Wilhelm Leibniz} notait $\dfrac{\mathrm{d}f}{\mathrm{d}x}$, notation encore omniprésente aujourd'hui en physique (vitesse $v = \dfrac{\mathrm{d}x}{\mathrm{d}t}$) et en économie (coût marginal). Trois génies, trois écritures, une même révolution : mesurer l'instantané.
\end{savaistu}

\begin{aretenir}[L'essentiel de la section]
\begin{itemize}
\item Si $f$ est dérivable en tout point d'un intervalle $I$, on définit sa \cle{fonction dérivée} $f' : x \mapsto f'(x)$ sur $I$.
\item Les formules usuelles sont à connaître par c\oe ur : $(k)'=0$ ; $(x)'=1$ ; $(x^2)'=2x$ ; $(x^n)' = n\,x^{n-1}$ ; $\left(\dfrac{1}{x}\right)' = -\dfrac{1}{x^2}$ ; $(\sqrt{x})' = \dfrac{1}{2\sqrt{x}}$.
\item Toujours préciser l'\cle{ensemble de dérivabilité} : $\dfrac{1}{x}$ n'est dérivable que pour $x \neq 0$, et $\sqrt{x}$ seulement pour $x > 0$.
\item La dérivée de $\sqrt{x}$ se démontre par la \cle{quantité conjuguée} : $\tau(h) = \dfrac{1}{\sqrt{x+h}+\sqrt{x}} \to \dfrac{1}{2\sqrt{x}}$.
\end{itemize}
\end{aretenir}

Dans la section suivante, nous verrons comment dériver des fonctions plus élaborées — sommes, produits par un nombre, produits et quotients de fonctions — en combinant les formules usuelles que nous venons d'établir.

\coursec{Opérations sur les fonctions dérivées}

Dans la section précédente, nous avons établi les dérivées des fonctions usuelles (puissances, racine carrée, etc.). En pratique, les fonctions rencontrées en technique (modélisation d'un coût, d'une trajectoire, d'un signal électrique) sont le plus souvent des \textbf{combinaisons} de ces fonctions de base : sommes, produits, quotients. Cette section donne les règles pour dériver ces combinaisons sans repartir de la définition à chaque fois.

\courssub{Dérivée d'une somme et d'un produit par un scalaire}

L'intuition est simple : la variation d'une somme est la somme des variations ; multiplier une fonction par une constante $k$ multiplie sa variation par $k$.

\begin{propriete}[Linéarité de la dérivation -- Admis au Probatoire]
Soient $u$ et $v$ deux fonctions dérivables sur un intervalle $I$ et $k \in \mathbb{R}$.
\begin{itemize}
    \item \textbf{Somme :} La fonction $u+v$ est dérivable sur $I$ et $\quad (u+v)' = u' + v'$.
    \item \textbf{Produit par un scalaire :} La fonction $ku$ est dérivable sur $I$ et $\quad (ku)' = k\,u'$.
\end{itemize}
Par récurrence, la dérivée d'une somme finie est la somme des dérivées : $(u_1 + \dots + u_n)' = u_1' + \dots + u_n'$.
\end{propriete}

\begin{exemplebox}[Dérivation d'un polynôme terme à terme]
Soit $f(x) = 3x^4 - 5x^2 + 2x - 7$ sur $\mathbb{R}$.
On dérive chaque terme en utilisant $(x^n)' = nx^{n-1}$ et la linéarité :
\[
f'(x) = 3(4x^3) - 5(2x) + 2(1) - 0 = 12x^3 - 10x + 2.
\]
\end{exemplebox}

\courssub{Dérivée d'un produit de deux fonctions}

Contrairement à la somme, la dérivée d'un produit n'est \textbf{pas} le produit des dérivées. L'intuition géométrique : l'aire d'un rectangle de côtés $u$ et $v$ varie quand $u$ change ($v$ fixe) et quand $v$ change ($u$ fixe).

\begin{experience}[Démonstration guidée de la formule du produit]
Soient $u$ et $v$ dérivables en $a$. On note $\Delta u = u(a+h)-u(a)$, $\Delta v = v(a+h)-v(a)$.
L'accroissement du produit est :
\[
\Delta(uv) = u(a+h)v(a+h) - u(a)v(a).
\]
\textbf{Étape 1 :} Ajouter et soustraire $u(a+h)v(a)$ (ou $u(a)v(a+h)$) :
\[
\Delta(uv) = \underbrace{u(a+h)v(a+h) - u(a+h)v(a)}_{\text{variation due à }v} + \underbrace{u(a+h)v(a) - u(a)v(a)}_{\text{variation due à }u}.
\]
\textbf{Étape 2 :} Factoriser :
\[
\Delta(uv) = u(a+h)\,\Delta v + v(a)\,\Delta u.
\]
\textbf{Étape 3 :} Diviser par $h$ et faire tendre $h \to 0$ :
\[
\frac{\Delta(uv)}{h} = u(a+h)\frac{\Delta v}{h} + v(a)\frac{\Delta u}{h}.
\]
Puisque $u$ est dérivable en $a$, elle est continue : $\lim_{h\to 0} u(a+h) = u(a)$.
On obtient finalement :
\[
(uv)'(a) = u(a)v'(a) + u'(a)v(a).
\]
\end{experience}

\begin{propriete}[Dérivée d'un produit -- Démonstration exigible au Probatoire]
Soient $u$ et $v$ dérivables sur $I$. La fonction $uv$ est dérivable sur $I$ et :
\[
(uv)' = u'v + uv'.
\]
\end{propriete}

\begin{attention}[Piège classique : $(uv)' \neq u'v'$]
Prenons $u(x)=x$ et $v(x)=x$. Alors $uv = x^2$.
\begin{itemize}
    \item Mauvaise formule : $u'v' = 1 \times 1 = 1$.
    \item Bonne formule : $(x^2)' = 2x$.
\end{itemize}
L'erreur vient du fait qu'on oublie que \emph{les deux} fonctions varient simultanément. \textbf{La formule exacte fait deux termes :} un où $u$ dérive et $v$ reste, l'autre où $v$ dérive et $u$ reste.
\end{attention}

\begin{exemplebox}[{Produit : $f(x) = (x^2+1)\sqrt{x}$ sur $]0,+\infty[$}]
On pose $u(x)=x^2+1$, $v(x)=\sqrt{x}=x^{1/2}$.
$u'(x)=2x$, $v'(x)=\frac{1}{2}x^{-1/2} = \frac{1}{2\sqrt{x}}$.
\[
f'(x) = u'v + uv' = 2x\sqrt{x} + (x^2+1)\frac{1}{2\sqrt{x}}.
\]
On simplifie souvent en mettant au même dénominateur $2\sqrt{x}$ :
\[
f'(x) = \frac{4x^2 + x^2 + 1}{2\sqrt{x}} = \frac{5x^2+1}{2\sqrt{x}}.
\]
\end{exemplebox}

\courssub{Dérivée d'un quotient et cas particulier de l'inverse}

Pour $f = u/v$ avec $v(x) \neq 0$, on peut soit admettre la formule, soit la redémontrer en écrivant $f = u \times \frac{1}{v}$ et en utilisant la formule du produit + la dérivée de l'inverse.

\begin{experience}[Dérivée de l'inverse $1/v$]
Soit $v$ dérivable et non nulle sur $I$. Posons $w = 1/v$. Alors $v \cdot w = 1$.
En dérivant (produit) : $v'w + vw' = 0 \implies vw' = -v'w \implies w' = -\frac{v'w}{v} = -\frac{v'}{v^2}$.
\end{experience}

\begin{propriete}[Dérivée de l'inverse -- Admise ou déduite du produit]
Si $v$ est dérivable et $v(x) \neq 0$ sur $I$, alors $1/v$ est dérivable et :
\[
\left(\frac{1}{v}\right)' = -\frac{v'}{v^2}.
\]
\end{propriete}

\begin{propriete}[Dérivée d'un quotient -- Admise au Probatoire]
Soient $u,v$ dérivables sur $I$ avec $v(x) \neq 0$. La fonction $u/v$ est dérivable sur $I$ et :
\[
\left(\frac{u}{v}\right)' = \frac{u'v - uv'}{v^2}.
\]
\textit{Mnémotechnique : « Dérivée du haut $\times$ bas moins haut $\times$ dérivée du bas, tout sur le bas au carré. »}
\end{propriete}

\begin{popfigure}
\begin{tikzpicture}[
    node distance=1.2cm and 1.5cm,
    every node/.style={font=\small},
    func/.style={rectangle, draw=popblue, thick, fill=popblueL, rounded corners, minimum width=2.5cm, minimum height=0.9cm, align=center},
    op/.style={diamond, draw=poporange, thick, fill=poporangeL, aspect=2, inner sep=2pt, align=center},
    rule/.style={rectangle, draw=popgreen, thick, fill=popgreenL, rounded corners, minimum width=3cm, minimum height=0.8cm, align=center, text width=2.8cm},
    arrow/.style={-Stealth, thick, popdark}
]
% Nœuds fonctions de base
\node[func] (u) {$u(x)$};
\node[func, right=of u] (v) {$v(x)$};

% Nœuds opérations
\node[op, below=of u] (sum) {$+$};
\node[op, below=of v] (prod) {$\times$};
\node[op, right=1.5cm of prod] (quot) {$/$};

% Flèches vers opérations
\draw[arrow] (u) -- (sum);
\draw[arrow] (v) -- (sum);
\draw[arrow] (u) -- (prod);
\draw[arrow] (v) -- (prod);
\draw[arrow] (u) -- (quot);
\draw[arrow] (v) -- (quot);

% Règles associées
\node[rule, below=of sum] (rsum) {$(u+v)' = u' + v'$};
\node[rule, below=of prod] (rprod) {$(uv)' = u'v + uv'$};
\node[rule, below=of quot] (rquot) {$\left(\frac{u}{v}\right)' = \frac{u'v-uv'}{v^2}$};

\draw[arrow] (sum) -- (rsum);
\draw[arrow] (prod) -- (rprod);
\draw[arrow] (quot) -- (rquot);

% Scalaires
\node[func, left=1.5cm of u] (k) {$k \in \mathbb{R}$};
\node[op, below=of k] (smult) {$\times$};
\node[rule, below=of smult] (rsmult) {$(ku)' = ku'$};
\draw[arrow] (k) -- (smult);
\draw[arrow] (u) -- (smult);
\draw[arrow] (smult) -- (rsmult);
\end{tikzpicture}
\legende{Arbre de décision : identifier la structure de la fonction pour choisir la règle de dérivation.}
\end{popfigure}

\begin{exoresolu}[Dérivation pas à pas d'une fonction rationnelle]
\textbf{Énoncé :} Dériver $f(x) = \dfrac{2x+1}{x-3}$ sur son domaine de définition.

\textbf{Solution détaillée :}
\begin{enumerate}
    \item \textbf{Domaine :} $D_f = \mathbb{R} \setminus \{3\}$ (dénominateur non nul).
    \item \textbf{Identification :} C'est un quotient. On pose :
    \[
    u(x) = 2x+1 \quad \Rightarrow \quad u'(x) = 2
    \]
    \[
    v(x) = x-3 \quad \Rightarrow \quad v'(x) = 1
    \]
    \item \textbf{Application de la formule :}
    \[
    f'(x) = \frac{u'(x)v(x) - u(x)v'(x)}{[v(x)]^2} = \frac{2(x-3) - (2x+1)\times 1}{(x-3)^2}.
    \]
    \item \textbf{Simplification du numérateur :}
    \[
    2(x-3) - (2x+1) = 2x - 6 - 2x - 1 = -7.
    \]
    \item \textbf{Résultat final :}
    \[
    \boxed{f'(x) = \frac{-7}{(x-3)^2}, \quad \forall x \in \mathbb{R}\setminus\{3\}.}
    \]
    \textit{Remarque :} La dérivée est toujours négative (sauf en $x=3$), la fonction est strictement décroissante sur $]-\infty,3[$ et sur $]3,+\infty[$.
\end{enumerate}
\end{exoresolu}

\courssub{Applications : polynômes et fonctions rationnelles}

\begin{propriete}[Dérivée d'un polynôme]
Soit $P(x) = a_n x^n + a_{n-1}x^{n-1} + \dots + a_1 x + a_0$.
Alors $P'(x) = n a_n x^{n-1} + (n-1)a_{n-1}x^{n-2} + \dots + a_1$.
La dérivée d'un polynôme de degré $n$ est un polynôme de degré $n-1$ (si $n\ge 1$).
\end{propriete}

\begin{savaistu}[Pourquoi dériver « simplifie » les polynômes ?]
Chaque dérivation abaisse le degré de 1. Après $n$ dérivations, un polynôme de degré $n$ devient constant, puis nul. C'est la base des \textbf{séries de Taylor} (Terminale) et de la résolution d'équations différentielles linéaires (modélisation amortissement, circuits RLC). En usine, si la position est un polynôme de degré 3, la vitesse est degré 2, l'accélération degré 1, le « jerk » (dérivée de l'accélération) est constant.
\end{savaistu}

\begin{exemplebox}[Fonction rationnelle simple : $f(x) = \frac{x^2+1}{x}$ sur $\mathbb{R}^*$]
On peut dériver directement au quotient ($u=x^2+1, v=x$) ou simplifier d'abord :
\[
f(x) = x + \frac{1}{x} = x + x^{-1}.
\]
La seconde méthode est plus rapide et moins source d'erreurs :
\[
f'(x) = 1 - x^{-2} = 1 - \frac{1}{x^2} = \frac{x^2-1}{x^2}.
\]
\textbf{Astuce :} Toujours vérifier si une simplification algébrique est possible avant d'appliquer la formule du quotient.
\end{exemplebox}

\courssub{Méthode et exercices de synthèse}

\begin{methode}[Dériver une fonction composée d'opérations usuelles]
\begin{enumerate}
    \item \textbf{Domaine :} Déterminer $D_f$ (zéros des dénominateurs, radicands négatifs, etc.).
    \item \textbf{Structure :} Identifier la dernière opération effectuée (somme ? produit ? quotient ?).
    \item \textbf{Poser $u, v$ :} Écrire explicitement $u(x), v(x)$ et calculer $u'(x), v'(x)$ à part.
    \item \textbf{Appliquer la règle :} Écrire la formule correspondante en remplaçant $u,v,u',v'$.
    \item \textbf{Simplifier :} Factoriser, réduire, mettre au même dénominateur pour obtenir une expression exploitable (signe, équation $f'=0$).
    \item \textbf{Conclure :} Redonner le domaine de $f'$ (souvent identique à $D_f$).
\end{enumerate}
\end{methode}

\begin{exoresolu}[Synthèse : produit et quotient]
\textbf{Énoncé :} Dériver $g(x) = \frac{(x^2+3)\sqrt{x}}{x+1}$ sur $D_g = [0,+\infty[\setminus\{-1\} = [0,+\infty[$.

\textbf{Solution :}
\begin{enumerate}
    \item Structure : Quotient principal. Numérateur = Produit.
    \item Poser :
    \[
    \begin{cases}
    U(x) = (x^2+3)\sqrt{x} & \text{(produit)} \\
    V(x) = x+1 & \Rightarrow V'(x)=1
    \end{cases}
    \]
    \item Dériver $U$ (produit) : $u_1=x^2+3 \Rightarrow u_1'=2x$ ; $v_1=\sqrt{x} \Rightarrow v_1'=\frac{1}{2\sqrt{x}}$.
    \[
    U' = u_1'v_1 + u_1v_1' = 2x\sqrt{x} + (x^2+3)\frac{1}{2\sqrt{x}} = \frac{4x^2 + x^2 + 3}{2\sqrt{x}} = \frac{5x^2+3}{2\sqrt{x}}.
    \]
    \item Appliquer quotient $g = U/V$ :
    \[
    g' = \frac{U'V - UV'}{V^2} = \frac{\frac{5x^2+3}{2\sqrt{x}}(x+1) - (x^2+3)\sqrt{x}\cdot 1}{(x+1)^2}.
    \]
    \item Simplifier (multiplier numérateur et dénominateur par $2\sqrt{x}$) :
    \[
    g'(x) = \frac{(5x^2+3)(x+1) - 2x(x^2+3)}{2\sqrt{x}(x+1)^2} = \frac{5x^3+5x^2+3x+3 - 2x^3-6x}{2\sqrt{x}(x+1)^2} = \frac{3x^3+5x^2-3x+3}{2\sqrt{x}(x+1)^2}.
    \]
\end{enumerate}
\end{exoresolu}

\begin{exercice}[Entraînement -- Dérivées de combinaisons]
Dériver les fonctions suivantes sur leur domaine de définition.
\begin{enumerate}
    \item $h_1(x) = (x-2)(x^3+1)$
    \item $h_2(x) = \dfrac{x^2-1}{x^2+1}$
    \item $h_3(x) = \dfrac{3x}{x^2+4}$
    \item $h_4(x) = \sqrt{x}(x - \frac{1}{x})$ pour $x>0$
    \item $h_5(x) = \dfrac{x+\sqrt{x}}{x-\sqrt{x}}$ pour $x>0, x\neq 1$
\end{enumerate}
\end{exercice}

\begin{corrige}[Exercice n°1 -- Éléments de correction]
\begin{enumerate}
    \item Produit : $h_1'(x) = 1\cdot(x^3+1) + (x-2)\cdot 3x^2 = x^3+1+3x^3-6x^2 = 4x^3-6x^2+1$.
    \item Quotient : $h_2'(x) = \frac{2x(x^2+1) - (x^2-1)2x}{(x^2+1)^2} = \frac{2x^3+2x-2x^3+2x}{(x^2+1)^2} = \frac{4x}{(x^2+1)^2}$.
    \item Quotient : $u(x)=3x \Rightarrow u'(x)=3$ ; $v(x)=x^2+4 \Rightarrow v'(x)=2x$.
    \[
    h_3'(x) = \frac{3(x^2+4) - 3x(2x)}{(x^2+4)^2} = \frac{3x^2+12-6x^2}{(x^2+4)^2} = \frac{12-3x^2}{(x^2+4)^2}.
    \]
    \item On réécrit $h_4(x) = x^{3/2} - x^{-1/2}$ pour $x>0$.
    $h_4'(x) = \frac{3}{2}x^{1/2} - \left(-\frac{1}{2}\right)x^{-3/2} = \frac{3}{2}\sqrt{x} + \frac{1}{2x\sqrt{x}}$.
    \item Simplifier par $\sqrt{x}$ : $h_5(x) = \frac{\sqrt{x}+1}{\sqrt{x}-1}$.
    Quotient : $u=\sqrt{x}+1, v=\sqrt{x}-1$. $u'=v'=\frac{1}{2\sqrt{x}}$.
    \[
    h_5' = \frac{\frac{1}{2\sqrt{x}}(\sqrt{x}-1) - (\sqrt{x}+1)\frac{1}{2\sqrt{x}}}{(\sqrt{x}-1)^2} = \frac{\frac{\sqrt{x}-1-\sqrt{x}-1}{2\sqrt{x}}}{(\sqrt{x}-1)^2} = \frac{-1/\sqrt{x}}{(\sqrt{x}-1)^2} = -\frac{1}{\sqrt{x}(\sqrt{x}-1)^2}.
    \]
\end{enumerate}
\end{corrige}

\coursec{Applications concrètes de la dérivation}

Après avoir maîtrisé les règles de dérivation des fonctions usuelles et les opérations sur les fonctions dérivées, nous disposons maintenant d'un outil puissant : la fonction dérivée $f'$. Mais à quoi sert-elle vraiment ? Au-delà du calcul formel, la dérivée est le \textbf{mathématicien du changement} : elle mesure la \emph{vitesse de variation} d'une grandeur à un instant précis. C'est le passage du « combien ? » (valeur de la fonction) au « à quelle vitesse ? » (valeur de la dérivée). Dans cette section, nous allons voir comment la dérivée modélise des situations concrètes venues de la mécanique, de l'économie et de la gestion, et comment elle permet de prendre des décisions optimales.

\courssub{Vitesse instantanée : la dérivée comme taux de variation instantané}

Imaginons un véhicule qui parcourt une route. Soit $d(t)$ la distance (en kilomètres) parcourue depuis le départ en fonction du temps $t$ (en heures). La \textbf{vitesse moyenne} sur un intervalle $[t_1 ; t_2]$ est le quotient $\frac{d(t_2)-d(t_1)}{t_2-t_1}$ : c'est le coefficient directeur de la sécante à la courbe de $d$ aux abscisses $t_1$ et $t_2$. Mais le conducteur regarde son compteur de vitesse : il veut connaître la \textbf{vitesse à l'instant $t$}, c'est-à-dire la limite des vitesses moyennes quand l'intervalle se resserre autour de $t$. C'est exactement la définition du nombre dérivé $d'(t)$.

\begin{definition}[Vitesse instantanée]
Si une fonction $d$ donne la position (ou la distance parcourue) en fonction du temps $t$, alors sa fonction dérivée $v = d'$ donne la \textbf{vitesse instantanée} à l'instant $t$ :
\[
v(t) = d'(t) = \lim_{h \to 0} \frac{d(t+h)-d(t)}{h}.
\]
L'unité de $v$ est l'unité de $d$ divisée par l'unité de $t$ (par exemple km/h si $d$ est en km et $t$ en h).
\end{definition}

\begin{popfigure}
\begin{tikzpicture}[scale=0.85]
\begin{axis}[
    width=\linewidth,
    height=0.55\linewidth,
    axis lines=middle,
    xlabel={$t$ (heures)},
    ylabel={$d(t)$ (km)},
    xmin=0, xmax=5.5,
    ymin=0, ymax=70,
    xtick={0,1,2,3,4,5},
    ytick={0,10,20,30,40,50,60,70},
    grid=major,
    grid style={popline},
    clip=false,
    samples=200,
    domain=0:5.5
]
% Courbe d(t) = 2t^2 + 5t
\addplot[popcurve, thick, popblue] {2*x^2 + 5*x};
\addlegendentry{$d(t)=2t^2+5t$}

% Points pour sécantes et tangente
\coordinate (A) at (axis cs:1,7);
\coordinate (B) at (axis cs:2,18);
\coordinate (C) at (axis cs:3,33);
\coordinate (T) at (axis cs:3,33);

% Sécante AB (vitesse moyenne entre 1h et 2h)
\draw[poporange, dashed, thick] (axis cs:0.5, -3) -- (axis cs:2.5, 28);
\node[poporange, anchor=south west] at (axis cs:0.5, -3) {Sécante $[1;2]$ : $v_{\text{moy}}=11$ km/h};

% Sécante BC (vitesse moyenne entre 2h et 3h)
\draw[popgreen, dashed, thick] (axis cs:1.5, 5) -- (axis cs:3.5, 41);
\node[popgreen, anchor=south west] at (axis cs:1.5, 5) {Sécante $[2;3]$ : $v_{\text{moy}}=15$ km/h};

% Tangente en t=3 (vitesse instantanée)
\draw[poppink, thick] (axis cs:1.5, 5.5) -- (axis cs:4.5, 60.5);
\node[poppink, anchor=south west] at (axis cs:4.6, 60) {Tangente en $t=3$ : $v(3)=17$ km/h};

% Points
\fill[popblue] (A) circle (2.5pt) node[above left] {$A(1;7)$};
\fill[popblue] (B) circle (2.5pt) node[above left] {$B(2;18)$};
\fill[popblue] (C) circle (2.5pt) node[above right] {$C(3;33)$};

% Flèches pour montrer le resserrement
\draw[<->, popmuted] (axis cs:1, -5) -- (axis cs:2, -5) node[midway, below] {$h=1$};
\draw[<->, popmuted] (axis cs:2, -8) -- (axis cs:3, -8) node[midway, below] {$h=1$};
\draw[<->, popmuted] (axis cs:2.5, -11) -- (axis cs:3, -11) node[midway, below] {$h \to 0$};
\end{axis}
\end{tikzpicture}
\legende{Distance parcourue $d(t)=2t^2+5t$ : les sécantes donnent les vitesses moyennes, la tangente donne la vitesse instantanée $v(t)=d'(t)$.}
\end{popfigure}

\begin{exemplebox}[Vitesse d'un véhicule]
Soit $d(t) = 2t^2 + 5t$ la distance (en km) parcourue par un camion en $t$ heures.
\begin{enumerate}
\item Déterminer la fonction vitesse $v(t)$.
\item Calculer la vitesse à $t = 3$ h.
\item Interpréter.
\end{enumerate}
\tcblower
\textbf{Solution.}
\begin{enumerate}
\item $v(t) = d'(t) = (2t^2)' + (5t)' = 4t + 5$ (km/h).
\item $v(3) = 4 \times 3 + 5 = 17$ km/h.
\item À l'instant $t = 3$ h, le camion roule à une vitesse instantanée de 17 km/h. C'est la pente de la tangente à la courbe de $d$ au point d'abscisse 3.
\end{enumerate}
\end{exemplebox}

\begin{attention}[Unités et interprétation]
Ne confondez pas $d(t)$ (distance, en km) et $v(t)$ (vitesse, en km/h). La dérivée change l'unité : si $d$ est en mètres et $t$ en secondes, $v$ est en m/s. Vérifiez toujours la cohérence des unités dans l'énoncé.
\end{attention}

\courssub{Coût marginal, recette et bénéfice : la dérivée en économie}

En gestion d'entreprise, on modélise souvent le \textbf{coût total} $C(q)$ pour produire $q$ unités d'un bien. La dérivée $C'(q)$ a une interprétation économique fondamentale.

\begin{definition}[Coût marginal]
Le \textbf{coût marginal} en $q$ est la dérivée $C'(q)$ de la fonction coût total. Il approxime le coût de production de \emph{l'unité supplémentaire} (la $(q+1)$-ième) lorsque $q$ est grand :
\[
C'(q) \approx C(q+1) - C(q).
\]
C'est le taux de variation instantané du coût par rapport à la quantité.
\end{definition}

\begin{savaistu}[Le coût marginal au Cameroun]
Dans les PME camerounaises (transformation agroalimentaire, menuiserie, BTP), les gestionnaires calculent le « coût marginal » pour décider d'accepter une commande supplémentaire. Si le prix de vente dépasse le coût marginal, l'opération est rentable à court terme. C'est exactement l'application de la dérivée $C'(q)$ !
\end{savaistu}

La \textbf{recette totale} pour une quantité $q$ vendue au prix unitaire $p$ est $R(q) = p \times q$. Si le prix dépend de la quantité (loi de la demande), $p$ devient une fonction de $q$ et $R(q) = p(q) \times q$. Le \textbf{bénéfice} (ou profit) est $B(q) = R(q) - C(q)$. Pour maximiser le bénéfice, on étudie $B'(q)$.

\begin{exemplebox}[Coût marginal et décision de production]
Une entreprise de fabrication de parpaings a un coût total (en milliers de FCFA) modélisé par $C(q) = 0,5q^2 + 20q + 500$ pour $q$ centaines de parpaings.
\begin{enumerate}
\item Calculer le coût marginal $C'(q)$.
\item Estimer le coût de production de la 201\textsuperscript{e} centaine de parpaings (c'est-à-dire au voisinage de $q=200$).
\item Si le prix de vente est 45 000 FCFA la centaine, l'entreprise a-t-elle intérêt à produire cette centaine supplémentaire ?
\end{enumerate}
\tcblower
\textbf{Solution.}
\begin{enumerate}
\item $C'(q) = q + 20$ (en milliers de FCFA par centaine).
\item Pour $q = 200$ : $C'(200) = 220$ milliers de FCFA = 220 000 FCFA par centaine.
\item Prix de vente = 45 000 FCFA la centaine = 45 milliers de FCFA. Comme $45 < 220$, le coût marginal dépasse largement le prix : \textbf{non}, il ne faut pas produire cette centaine supplémentaire (perte de 175 000 FCFA).
\end{enumerate}
\end{exemplebox}

\courssub{Du signe de la dérivée aux variations : annonce intuitive}

Nous allons maintenant faire le lien fondamental entre le \textbf{signe de $f'$} et le \textbf{sens de variation de $f$}. Cette propriété, admise à ce stade, sera démontrée rigoureusement dans la leçon sur les variations.

\begin{propriete}[Signe de la dérivée et variations de la fonction (admis)]
Soit $f$ une fonction dérivable sur un intervalle $I$.
\begin{itemize}
\item Si $f'(x) > 0$ pour tout $x \in I$, alors $f$ est \textbf{strictement croissante} sur $I$.
\item Si $f'(x) < 0$ pour tout $x \in I$, alors $f$ est \textbf{strictement décroissante} sur $I$.
\item Si $f'(x) = 0$ pour tout $x \in I$, alors $f$ est \textbf{constante} sur $I$.
\end{itemize}
\textit{Interprétation graphique :} $f' > 0$ signifie que la tangente a une pente positive $\Rightarrow$ la courbe monte ; $f' < 0$ signifie pente négative $\Rightarrow$ la courbe descend.
\end{propriete}

\begin{popfigure}
\begin{tikzpicture}[scale=0.9]
\begin{axis}[
    width=\linewidth,
    height=0.5\linewidth,
    axis lines=middle,
    xlabel={$x$},
    ylabel={$f(x)$},
    xmin=-3, xmax=3,
    ymin=-2, ymax=5,
    xtick={-2,-1,0,1,2},
    ytick={-1,0,1,2,3,4},
    grid=major,
    grid style={popline},
    samples=200,
    domain=-3:3
]
% f(x) = -x^2 + 2x + 3 (parabole renversée)
\addplot[popcurve, thick, popblue] {-x^2 + 2*x + 3};
\addlegendentry{$f(x)=-x^2+2x+3$}

% Tangentes avec pentes
\draw[popgreen, thick] (axis cs:-2.5, -5.25) -- (axis cs:-0.5, 2.75);
\node[popgreen, anchor=south] at (axis cs:-1.5, -1) {$f'>0$ : croissante};

\draw[poporange, thick] (axis cs:0.5, 3.75) -- (axis cs:2.5, 3.75);
\node[poporange, anchor=north] at (axis cs:1.5, 4.2) {$f'=0$ : horizontale};

\draw[poppink, thick] (axis cs:1.5, 4.75) -- (axis cs:3.5, 0.25);
\node[poppink, anchor=north] at (axis cs:2.5, 2.5) {$f'<0$ : décroissante};

% Point au sommet
\fill[popblue] (axis cs:1, 4) circle (3pt) node[above] {$x=1$, $f'=0$};
\end{axis}
\end{tikzpicture}
\legende{Illustration : $f'(x) > 0 \Rightarrow f$ croissante ; $f'(x) = 0 \Rightarrow$ tangente horizontale (extremum possible) ; $f'(x) < 0 \Rightarrow f$ décroissante.}
\end{popfigure}

\courssub{Recherche d'extremum : résoudre $f'(x) = 0$}

Un \textbf{extremum} (maximum ou minimum) d'une fonction dérivable ne peut se produire qu'en un point où la tangente est horizontale, c'est-à-dire où $f'(x) = 0$. C'est une condition \textbf{nécessaire} mais pas suffisante : il faut vérifier le changement de signe de $f'$ (ou la concavité) pour conclure.

\begin{methode}[Recherche d'un extremum d'une fonction dérivable]
\begin{enumerate}
\item Calculer la fonction dérivée $f'(x)$.
\item Résoudre l'équation $f'(x) = 0$ pour trouver les \textbf{valeurs critiques}.
\item Étudier le \textbf{signe de $f'$} de part et d'autre de chaque valeur critique (tableau de signes).
\item Conclure :
\begin{itemize}
\item Si $f'$ passe de $+$ à $-$ $\Rightarrow$ \textbf{maximum local}.
\item Si $f'$ passe de $-$ à $+$ $\Rightarrow$ \textbf{minimum local}.
\item Si $f'$ ne change pas de signe $\Rightarrow$ pas d'extremum (point d'inflexion à tangente horizontale).
\end{itemize}
\item Calculer la valeur de $f$ en ce point pour obtenir l'extremum.
\end{enumerate}
\end{methode}

\begin{attention}[Piège classique : $f'(a)=0$ ne suffit pas !]
Ne concluez \textbf{jamais} « $f$ a un maximum en $a$ » sous prétexte que $f'(a)=0$.
\begin{itemize}
\item Exemple : $f(x)=x^3$ a $f'(0)=0$ mais pas d'extremum en 0 (fonction strictement croissante).
\item Vérifiez \textbf{toujours} le changement de signe de $f'$ (tableau de signes) ou utilisez la dérivée seconde (en Terminale).
\end{itemize}
\end{attention}

\courssub{Problème guidé d'optimisation : maximiser une recette}

Appliquons la méthode complète à un problème concret d'optimisation de recette.

\begin{exoresolu}[Maximisation de la recette d'une entreprise]
Une entreprise vend un produit. La recette totale (en centaines de FCFA) en fonction du prix de vente $x$ (en centaines de FCFA) est modélisée par :
\[
R(x) = -2x^2 + 40x, \quad \text{pour } x \in [0 ; 20].
\]
On cherche le prix $x$ qui maximise la recette et la valeur de cette recette maximale.
\tcblower
\textbf{Étape 1 : Modélisation et fonction dérivée.}
$R$ est une fonction polynôme, dérivable sur $\mathbb{R}$.
$R'(x) = (-2x^2)' + (40x)' = -4x + 40$.

\textbf{Étape 2 : Valeurs critiques.}
$R'(x) = 0 \;\Leftrightarrow\; -4x + 40 = 0 \;\Leftrightarrow\; 4x = 40 \;\Leftrightarrow\; x = 10$.
La valeur critique $x=10$ appartient bien à l'intervalle $[0 ; 20]$.

\textbf{Étape 3 : Tableau de signes de $R'$.}
$R'(x) = -4x + 40 = -4(x - 10)$.
\begin{center}
\begin{tabular}{|c|ccccc|}\hline
$x$ & $0$ & & $10$ & & $20$ \\ \hline
$R'(x)$ & & $+$ & $0$ & $-$ & \\ \hline
$R(x)$ & $0$ & $\nearrow$ & $200$ & $\searrow$ & $0$ \\ \hline
\end{tabular}
\end{center}

\textbf{Étape 4 : Interprétation.}
$R'$ passe de $+$ à $-$ en $x=10$ : $R$ a un \textbf{maximum} en $x=10$.
La recette maximale est $R(10) = -2 \times 10^2 + 40 \times 10 = -200 + 400 = 200$ (centaines de FCFA).

\textbf{Étape 5 : Conclusion rédigée.}
\begin{quote}
\textbf{Le prix de vente qui maximise la recette est 1 000 FCFA (soit $x=10$ centaines de FCFA). La recette maximale atteinte est alors de 20 000 FCFA (soit $R(10)=200$ centaines de FCFA).}
\end{quote}
\end{exoresolu}

\begin{popfigure}
\begin{tikzpicture}[scale=0.85]
\begin{axis}[
    width=\linewidth,
    height=0.55\linewidth,
    axis lines=middle,
    xlabel={$x$ (centaines de FCFA)},
    ylabel={$R(x)$ (centaines de FCFA)},
    xmin=0, xmax=21,
    ymin=0, ymax=220,
    xtick={0,5,10,15,20},
    ytick={0,50,100,150,200},
    grid=major,
    grid style={popline},
    samples=200,
    domain=0:20,
    clip=false
]
% Parabole R(x) = -2x^2 + 40x
\addplot[popcurve, thick, popblue] {-2*x^2 + 40*x};
\addlegendentry{$R(x)=-2x^2+40x$}

% Tangente horizontale au sommet (10, 200)
\draw[poporange, dashed, thick] (axis cs:0, 200) -- (axis cs:20, 200);
\node[poporange, anchor=south west] at (axis cs:0.5, 202) {Tangente horizontale : $R'(10)=0$};

% Point maximum
\fill[poporange] (axis cs:10, 200) circle (4pt) node[above right] {$\text{Max } R(10)=200$};

% Axes pointillés vers le point
\draw[popmuted, thin, dashed] (axis cs:10, 0) -- (axis cs:10, 200);
\draw[popmuted, thin, dashed] (axis cs:0, 200) -- (axis cs:10, 200);
\node[popmuted, anchor=north] at (axis cs:10, -5) {$x=10$};
\node[popmuted, anchor=east] at (axis cs:-0.5, 200) {$200$};

% Zone de variation
\draw[popgreen, thick, <->] (axis cs:2, 72) -- (axis cs:8, 176) node[midway, above, sloped] {Croissante ($R'>0$)};
\draw[poppink, thick, <->] (axis cs:12, 176) -- (axis cs:18, 72) node[midway, above, sloped] {Décroissante ($R'<0$)};
\end{axis}
\end{tikzpicture}
\legende{Recette $R(x)=-2x^2+40x$ : maximum en $x=10$ (prix 1 000 FCFA) avec tangente horizontale.}
\end{popfigure}

\courssub{Rédiger une démarche complète et justifiée}

Au Probatoire, on attend de vous une \textbf{rédaction structurée} : ce n'est pas seulement le résultat final qui compte, mais le cheminement logique. Voici le canevas type pour tout problème d'optimisation ou d'étude de taux de variation.

\begin{methode}[Rédaction type d'un problème avec dérivée (Probatoire)]
\begin{enumerate}
\item \textbf{Modélisation} : Définir clairement la fonction étudiée ($f$, $C$, $R$, $B$, $d$...), sa variable ($x$, $q$, $t$...), son unité et son domaine de définition (contexte du problème).
\item \textbf{Calcul de la dérivée} : Écrire $f'(x)$ en justifiant chaque étape (règle de la somme, produit, puissance, etc.).
\item \textbf{Résolution / Évaluation} : Selon la question :
\begin{itemize}
\item Pour un extremum : résoudre $f'(x)=0$, faire le \textbf{tableau de signes de $f'$}, conclure sur la nature (max/min).
\item Pour un taux instantané : calculer $f'(a)$ et donner l'unité.
\end{itemize}
\item \textbf{Interprétation} : Traduire le résultat mathématique en phrase dans le langage du problème (prix optimal, vitesse, coût marginal, bénéfice maximal...).
\item \textbf{Conclusion} : Réponse finale encadrée, avec unités.
\end{enumerate}
\end{methode}

\begin{exoresolu}[Bénéfice maximal : rédaction complète]
Une entreprise produit et vend $q$ unités d'un article. Le coût total (en FCFA) est $C(q) = 2q^2 + 100q + 5000$. Le prix de vente unitaire est fixé à 500 FCFA.
\begin{enumerate}
\item Exprimer la recette $R(q)$ et le bénéfice $B(q)$.
\item Déterminer la quantité $q$ qui maximise le bénéfice et la valeur de ce maximum.
\end{enumerate}
\tcblower
\textbf{Solution.}

\textbf{1. Modélisation.}
$q \ge 0$ (quantité produite et vendue).
Recette : $R(q) = 500 \times q = 500q$ (FCFA).
Bénéfice : $B(q) = R(q) - C(q) = 500q - (2q^2 + 100q + 5000) = -2q^2 + 400q - 5000$.
$B$ est polynôme, dérivable sur $\mathbb{R}_+$.

\textbf{2. Dérivée et valeur critique.}
$B'(q) = -4q + 400$.
$B'(q) = 0 \;\Leftrightarrow\; -4q + 400 = 0 \;\Leftrightarrow\; q = 100$.

\textbf{3. Tableau de signes de $B'$.}
$B'(q) = -4(q - 100)$.
\begin{center}
\begin{tabular}{|c|ccccc|}\hline
$q$ & $0$ & & $100$ & & $+\infty$ \\ \hline
$B'(q)$ & & $+$ & $0$ & $-$ & \\ \hline
$B(q)$ & $-5000$ & $\nearrow$ & $15\,000$ & $\searrow$ & $-\infty$ \\ \hline
\end{tabular}
\end{center}

\textbf{4. Interprétation.}
$B'$ passe de $+$ à $-$ en $q=100$ : $B$ admet un \textbf{maximum} en $q=100$.
$B(100) = -2(100)^2 + 400(100) - 5000 = -20\,000 + 40\,000 - 5\,000 = 15\,000$ FCFA.

\textbf{5. Conclusion.}
\begin{quote}
\textbf{L'entreprise doit produire et vendre 100 unités pour maximiser son bénéfice. Le bénéfice maximal est de 15 000 FCFA.}
\end{quote}
\end{exoresolu}

\courssub{Synthèse et regard vers la Terminale}

\begin{aretenir}[Ce qu'il faut retenir des applications de la dérivation (1\textsuperscript{re} Tech)]
\begin{itemize}
\item \textbf{Vitesse instantanée :} $v(t) = d'(t)$ où $d$ est la position. Unité : unité de $d$ / unité de $t$.
\item \textbf{Coût marginal :} $C'(q)$ approxime le coût de l'unité supplémentaire. Décision : produire si $p > C'(q)$.
\item \textbf{Recette et bénéfice :} $R(q) = p(q)\cdot q$, $B(q) = R(q) - C(q)$. Optimiser $\Rightarrow$ étudier $B'(q)$.
\item \textbf{Signe de $f'$ et variations :} $f'>0 \Rightarrow f$ croissante ; $f'<0 \Rightarrow f$ décroissante (admis, prouvé en Terminale).
\item \textbf{Extremum :} Condition nécessaire $f'(a)=0$ (tangente horizontale). Condition suffisante : changement de signe de $f'$ (tableau de signes).
\item \textbf{Rédaction :} Modéliser $\rightarrow$ Dériver $\rightarrow$ Résoudre/Étudier le signe $\rightarrow$ Interpréter $\rightarrow$ Conclure avec unités.
\end{itemize}
\end{aretenir}

\begin{savaistu}[Dérivée et intelligence artificielle]
Les algorithmes d'apprentissage automatique (réseaux de neurones) utilisent massivement la \textbf{rétropropagation du gradient}, qui n'est autre que le calcul de dérivées partielles pour minimiser une fonction de coût. La dérivée que vous apprenez à manipuler aujourd'hui est le moteur mathématique de l'IA moderne !
\end{savaistu}

\begin{exercice}[Entraînement : applications variées]
\begin{enumerate}
\item Un projectile est lancé verticalement. Sa hauteur (en mètres) au bout de $t$ secondes est $h(t) = -5t^2 + 30t + 2$.
\begin{itemize}
\item[a)] Donner la vitesse $v(t)$ et l'accélération $a(t)$.
\item[b)] À quel instant le projectile atteint-il sa hauteur maximale ? Quelle est cette hauteur ?
\item[c)] Quelle est la vitesse au moment du retour au sol ($h=0$) ?
\end{itemize}
\item Une entreprise a un coût total $C(q) = 0,1q^3 - 5q^2 + 120q + 2000$ (en milliers de FCFA) pour $q$ centaines d'articles. Le prix de vente est de 200 000 FCFA la centaine d'articles.
\begin{itemize}
\item[a)] Exprimer le bénéfice $B(q)$ (en milliers de FCFA).
\item[b)] Déterminer la production optimale (maximisant le bénéfice) et le bénéfice maximal.
\end{itemize}
\item La population d'une bactérie évolue selon $P(t) = 100 e^{0,2t}$ (modèle admis, dérivée de $e^{kt}$ est $ke^{kt}$).
\begin{itemize}
\item[a)] Calculer le taux d'accroissement instantané $P'(t)$.
\item[b)] Interpréter $P'(5)$.
\end{itemize}
\end{enumerate}
\end{exercice}

\begin{corrige}[Exercice : Applications variées]
\begin{enumerate}
\item[a)] $v(t) = h'(t) = -10t + 30$ (m/s). $a(t) = v'(t) = -10$ (m/s²).
\item[b)] $v(t)=0 \Leftrightarrow -10t+30=0 \Leftrightarrow t=3$ s. $v$ passe de $+$ à $-$ $\Rightarrow$ maximum. $h(3) = -45+90+2 = 47$ m.
\item[c)] $h(t)=0 \Leftrightarrow -5t^2+30t+2=0 \Leftrightarrow t = \frac{-30 \pm \sqrt{900+40}}{-10}$. Racine positive : $t \approx 6,07$ s. $v(6,07) \approx -30,7$ m/s (vers le bas).
\item[a)] Recette : $R(q) = 200q$ (milliers de FCFA). Bénéfice : $B(q) = R(q) - C(q) = 200q - (0,1q^3 - 5q^2 + 120q + 2000) = -0,1q^3 + 5q^2 + 80q - 2000$.
\item[b)] $B'(q) = -0,3q^2 + 10q + 80$.
$B'(q)=0 \Leftrightarrow 0,3q^2 - 10q - 80 = 0 \Leftrightarrow 3q^2 - 100q - 800 = 0$.
$\Delta = 10000 + 9600 = 19600 = 140^2$.
$q = \frac{100 \pm 140}{6}$. La racine positive est $q = \frac{240}{6} = 40$ (centaines d'articles).
Tableau de signes de $B'$ : $B'(q) = -0,3(q-40)(q + \frac{20}{3})$. Sur $\mathbb{R}_+$, $B'$ passe de $+$ à $-$ en $q=40$ $\Rightarrow$ Maximum.
$B(40) = -0,1(64000) + 5(1600) + 3200 - 2000 = -6400 + 8000 + 3200 - 2000 = 2800$ milliers de FCFA = 2 800 000 FCFA.
\textbf{Production optimale : 4 000 articles. Bénéfice maximal : 2 800 000 FCFA.}
\item[a)] $P'(t) = 100 \times 0,2 e^{0,2t} = 20 e^{0,2t}$.
\item[b)] $P'(5) = 20 e^{1} \approx 54,36$ : à l'instant $t=5$, la population augmente d'environ 54 bactéries par unité de temps.
\end{enumerate}
\end{corrige}

\coursec{Synthèse, compléments et vers la Terminale}

Après avoir exploré les applications concrètes de la dérivation — optimisation de coûts, étude de mouvements, rendements en atelier — nous disposons maintenant d'une boîte à outils complète. Cette section finale a pour but de \textbf{structurer les connaissances}, de \textbf{fixer les formules indispensables}, d'entraîner le \textbf{raisonnement graphique} et de préparer le passage en Terminale technique. C'est le moment de transformer la « boîte à outils » en « méthode de travail » pour l'examen du Probatoire.

\courssub{Bilan synthétique : carte mentale de la leçon}

Toute la construction de ce chapitre repose sur un enchaînement logique : on part d'une variation moyenne pour atteindre une variation instantanée, ce qui définit la tangente, puis la fonction dérivée, dont les règles de calcul permettent de résoudre des problèmes techniques.

\begin{popfigure}
\begin{tikzpicture}[
    mindmap,
    concept color=popblue!80,
    level 1/.style={level distance=4.5cm, sibling angle=90},
    level 2/.style={level distance=3.5cm, sibling angle=45},
    level 3/.style={level distance=3cm, sibling angle=30},
    every node/.style={font=\small, text=white},
    edge from parent/.style={draw=popblue!50, thick}
]
\node[concept, scale=1.2, align=center]{Dérivation\\ (Première Tech)}
    child[concept color=poporange] {
        node[concept, align=center]{Taux d'accroissement\\ $\frac{f(x)-f(a)}{x-a}$}
        child { node[concept, scale=0.7, align=center]{Variation moyenne\\ sur $[a;x]$} }
        child { node[concept, scale=0.7, align=center]{Pente de la sécante\\ $(A;M)$} }
    }
    child[concept color=popgreen] {
        node[concept, align=center]{Nombre dérivé\\ $f'(a)=\lim_{x\to a}\frac{f(x)-f(a)}{x-a}$}
        child { node[concept, scale=0.7, align=center]{Limite du taux\\ d'accroissement} }
        child { node[concept, scale=0.7, align=center]{Pente de la tangente\\ en $A(a;f(a))$} }
        child { node[concept, scale=0.7, align=center]{Vitesse instantanée\\ si $f$=position} }
    }
    child[concept color=poppurple] {
        node[concept, align=center]{Tangente\\ $\mathcal{T}_A$}
        child { node[concept, scale=0.7] {Équation : $y=f'(a)(x-a)+f(a)$} }
        child { node[concept, scale=0.7, align=center]{Meilleure approx.\nolinebreak\\ affine locale} }
        child { node[concept, scale=0.7, align=center]{Outil d'approx.\nolinebreak\\ numérique} }
    }
    child[concept color=poppink] {
        node[concept, align=center]{Fonction dérivée\\ $f'$}
        child { node[concept, scale=0.7, align=center]{Formules usuelles\\ $(x^n, \sqrt{x}, 1/x,\dots)$} }
        child { node[concept, scale=0.7, align=center]{Opérations\\ $(u\pm v)', (uv)', (u/v)'$} }
        child { node[concept, scale=0.7, align=center]{Signe de $f'$ $\leftrightarrow$\\ variations de $f$} }
        child { node[concept, scale=0.7, align=center]{Applications :\\ optim., physique, éco.} }
    };
\end{tikzpicture}
\legende{Carte mentale : du taux d'accroissement aux applications techniques, en passant par le nombre dérivé, la tangente et la fonction dérivée.}
\end{popfigure}

\courssub{Tableau récapitulatif des formules à mémoriser}

Ce tableau constitue le « socle dur » exigible au Probatoire. Il doit être connu par cœur, y compris les conditions d'application (ensembles de définition et de dérivabilité).

\begin{center}
\begin{tabularx}{\linewidth}{|>{\columncolor{popblueL}}l|X|X|}\hline
\textbf{Catégorie} & \textbf{Formule / Règle} & \textbf{Condition / Remarque} \\\hline
\rowcolor{poporangeL}
\textbf{Fonctions usuelles} & $(c)' = 0$ & $c \in \mathbb{R}$ \\\cline{2-3}
& $(x^n)' = n x^{n-1}$ & $n \in \mathbb{Z}$ (ou $\mathbb{Q}$ si $x>0$) \\\cline{2-3}
& $(\sqrt{x})' = \frac{1}{2\sqrt{x}}$ & $x > 0$ (non dérivable en $0$) \\\cline{2-3}
& $\left(\frac{1}{x}\right)' = -\frac{1}{x^2}$ & $x \neq 0$ \\\hline
\rowcolor{popgreenL}
\textbf{Opérations} & $(u+v)' = u' + v'$ & Sur $D_u \cap D_v$ \\\cline{2-3}
& $(u-v)' = u' - v'$ & Sur $D_u \cap D_v$ \\\cline{2-3}
& $(uv)' = u'v + uv'$ & Sur $D_u \cap D_v$ \\\cline{2-3}
& $\left(\frac{u}{v}\right)' = \frac{u'v - uv'}{v^2}$ & Sur $D_u \cap D_v \setminus \{x; v(x)=0\}$ \\\hline
\rowcolor{poppurpleL}
\textbf{Tangente} & $\mathcal{T}_A : y = f'(a)(x-a) + f(a)$ & $f$ dérivable en $a$ \\\cline{2-3}
& Pente $m = f'(a)$ & $m \in \mathbb{R}$ \\\cline{2-3}
& Ordonnée à l'origine $b = f(a) - a f'(a)$ & Vérification : $y(a)=f(a)$ \\\hline
\rowcolor{poppinkL}
\textbf{Lien variations / signe} & $f' \ge 0 \iff f$ croissante & Sur un intervalle \\\cline{2-3}
& $f' \le 0 \iff f$ décroissante & Sur un intervalle \\\cline{2-3}
& $f'(a)=0$ + changement de signe & Extremum local en $a$ \\\hline
\end{tabularx}
\end{center}

\begin{attention}[Erreur fréquente : dériver sans vérifier l'ensemble de dérivabilité]
Un piège classique au Probatoire : on applique les formules mécaniquement sans vérifier où la fonction \emph{existe} et où elle est \emph{dérivable}.
\begin{itemize}
    \item \textbf{Racine carrée :} $f(x)=\sqrt{x}$ définie sur $[0;+\infty[$, mais dérivable seulement sur $]0;+\infty[$. En $0$, la tangente est verticale (pente infinie), $f'(0)$ n'existe pas.
    \item \textbf{Quotient :} Pour $f(x)=\frac{x^2+1}{x-2}$, on dérive avec la formule du quotient, mais le domaine de $f'$ est $\mathbb{R}\setminus\{2\}$ (même domaine que $f$).
    \item \textbf{Valeur absolue :} $f(x)=|x|$ n'est pas dérivable en $0$ (angle rentrant).
\end{itemize}
\textbf{Réflexe :} Avant toute dérivation, écrivez $D_f$ et $D_{f'}$.
\end{attention}

\courssub{Lecture graphique croisée : courbe de $f$ et signe de $f'$}

Au Probatoire, on vous donne souvent la courbe représentative $\mathcal{C}_f$ d'une fonction $f$ et on vous demande d'esquisser le tableau de signes de $f'$, ou l'inverse. C'est un exercice de \textbf{lecture graphique} sans calcul formel.

\begin{propriete}[Lecture graphique de la dérivée]
Soit $f$ dérivable sur un intervalle $I$ et $\mathcal{C}_f$ sa courbe.
\begin{itemize}
    \item $f'(x_0) > 0 \iff$ La tangente en $x_0$ a une \textbf{pente positive}.
    \item $f'(x_0) < 0 \iff$ La tangente en $x_0$ a une \textbf{pente négative}.
    \item $f'(x_0) = 0 \iff$ La tangente en $x_0$ est \textbf{horizontale} (extremum ou point d'inflexion à tangente horizontale).
    \item Plus la courbe est \textbf{raide}, plus $|f'|$ est grand.
\end{itemize}
\end{propriete}

\begin{popfigure}
\begin{tikzpicture}
\begin{groupplot}[
    group style={group size=1 by 2, vertical sep=1.5cm},
    width=\linewidth, height=5.5cm,
    axis lines=middle,
    xlabel=$x$, ylabel style={anchor=south},
    xtick={0,1,2,3,4}, ytick={-1,0,1,2},
    xmin=-0.5, xmax=4.5, ymin=-1.2, ymax=2.2,
    tick label style={font=\scriptsize},
    every axis x label/.style={at={(ticklabel* cs:1.02)}, anchor=west},
    every axis y label/.style={at={(ticklabel* cs:1.02)}, anchor=south},
]
% Courbe du haut : f(x) = 1/3 x^3 - 2x^2 + 3x
% f'(x) = x^2 - 4x + 3 = (x-1)(x-3) -> zeros at 1 and 3.
\nextgroupplot[ylabel=$f(x)$, title style={yshift=0.3cm}, title={\textbf{Courbe $\mathcal{C}_f$ de la fonction $f$}}]
\addplot[popcurve, domain=0:4, samples=200, thick] {x^3/3 - 2*x^2 + 3*x};
% Points remarquables : A(1, 4/3), B(3, 0)
\addplot[only marks, mark=*, poporange, mark size=2.5] coordinates {(1,1.333) (3,0)};
\node[poporange, anchor=south west, font=\small] at (axis cs:1,1.333) {$A(1;\,4/3)$};
\node[poporange, anchor=north east, font=\small] at (axis cs:3,0) {$B(3;\,0)$};
% Tangentes horizontales (segments)
\draw[poporange, dashed, thick] (axis cs:0.2,1.333) -- (axis cs:1.8,1.333);
\draw[poporange, dashed, thick] (axis cs:2.2,0) -- (axis cs:3.8,0);

% Tableau de signes du bas (simulé avec axis)
\nextgroupplot[
    ylabel=$f'(x)$,
    ytick={-1,0,1}, yticklabels={$-$ , $0$ , $+$},
    ymin=-1.5, ymax=1.5,
    xtick={0,1,2,3,4},
    xmin=-0.5, xmax=4.5,
    title style={yshift=0.3cm}, title={\textbf{Tableau de signes de $f'$ (déduit de $\mathcal{C}_f$)}}
]
% Zones de signe : + sur [0,1], - sur [1,3], + sur [3,4]
\addplot[popgreen!30, domain=0:1, samples=2] {1} \closedcycle;
\addplot[popred!30, domain=1:3, samples=2] {-1} \closedcycle;
\addplot[popgreen!30, domain=3:4, samples=2] {1} \closedcycle;
% Ligne zéro
\addplot[popline, thick, domain=0:4] {0};
% Points nuls
\addplot[only marks, mark=*, poporange, mark size=2.5] coordinates {(1,0) (3,0)};
% Flèches de variation (texte)
\node[popgreen, font=\small] at (axis cs:0.5, 0.7) {$f'>0 \;\nearrow$};
\node[popred, font=\small] at (axis cs:2, -0.7) {$f'<0 \;\searrow$};
\node[popgreen, font=\small] at (axis cs:3.5, 0.7) {$f'>0 \;\nearrow$};
\end{groupplot}
\end{tikzpicture}
\legende{En haut : courbe $\mathcal{C}_f$ de $f(x)=\frac{1}{3}x^3-2x^2+3x$. En bas : tableau de signes de $f'$ aligné verticalement. Les zéros de $f'$ ($x=1$ et $x=3$) correspondent aux tangentes horizontales de $\mathcal{C}_f$ (points $A$ et $B$). Le signe de $f'$ donne le sens de variation de $f$.}
\end{popfigure}

\begin{methode}[Esquisser le signe de $f'$ à partir de $\mathcal{C}_f$]
\begin{enumerate}
    \item Repérer les points où la tangente est \textbf{horizontale} $\rightarrow$ Zéros de $f'$ ($f'=0$).
    \item Sur chaque intervalle entre ces points, regarder si la courbe \textbf{monte} ($f'>0$) ou \textbf{descend} ($f'<0$).
    \item Évaluer la \textbf{raideur} : plus c'est raide, plus $|f'|$ est grand (utile pour esquisser la courbe de $f'$).
    \item Attention aux points où $f$ n'est pas dérivable (angles, tangentes verticales, asymptotes) : $f'$ n'existe pas (double barre $||$ au tableau).
\end{enumerate}
\end{methode}

\courssub{Compléments utiles pour la Terminale (hors exigibilité au Probatoire)}

Ces notions ne sont \textbf{pas au programme de Première} et ne tomberont pas au Probatoire. Elles sont présentées ici pour donner du sens aux formules que vous utilisez et faciliter l'entrée en Terminale technique.

\begin{savaistu}[Terminale : Dérivée de $x \mapsto (ax+b)^n$ et composée $u(ax+b)$]
En Terminale, on généralise la dérivation des puissances aux fonctions composées de la forme $u(ax+b)$ où $a,b \in \mathbb{R}$.
\begin{propriete}[Règle de la chaîne (cas affine intérieur)]
Si $u$ est dérivable sur $I$, alors $x \mapsto u(ax+b)$ est dérivable sur $J = \{x \in \mathbb{R} \mid ax+b \in I\}$ et :
\[ \big(u(ax+b)\big)' = a \times u'(ax+b) \]
\end{propriete}
\textbf{Application immédiate :} Pour $n \in \mathbb{N}^*$, $\big((ax+b)^n\big)' = a \cdot n (ax+b)^{n-1} = an(ax+b)^{n-1}$.
\end{savaistu}

\begin{exemplebox}[Exemples de dérivées composées (niveau Terminale)]
\begin{itemize}
    \item $f(x) = (3x+1)^2$. Ici $u(X)=X^2$, $a=3, b=1$. $f'(x) = 3 \times 2(3x+1) = 6(3x+1)$.
    \item $g(x) = \sqrt{5x-2} = (5x-2)^{1/2}$. $g'(x) = 5 \times \frac{1}{2}(5x-2)^{-1/2} = \frac{5}{2\sqrt{5x-2}}$ (pour $x > 2/5$).
    \item $h(x) = \frac{1}{2x+3} = (2x+3)^{-1}$. $h'(x) = 2 \times (-1)(2x+3)^{-2} = -\frac{2}{(2x+3)^2}$.
\end{itemize}
\textit{En Première, on développerait $(3x+1)^2 = 9x^2+6x+1$ pour dériver : $18x+6 = 6(3x+1)$. La règle de la chaîne évite le développement.}
\end{exemplebox}

\begin{savaistu}[Terminale : Notation différentielle $\frac{dy}{dx}$ et interprétation physique]
La notation de Leibniz $\frac{dy}{dx}$ (ou $\frac{df}{dx}$) insiste sur le rapport de deux variations infinitésimales.
\begin{itemize}
    \item Si $y = f(x)$, on note $dy = f'(x) \, dx$. C'est la \textbf{différentielle} de $f$.
    \item $\frac{dy}{dx} = f'(x)$ se lit « dérivée de $y$ par rapport à $x$ ».
\end{itemize}
\textbf{Interprétation physique (Cinématique) :}
Si $x(t)$ est la position (en m) au temps $t$ (en s) :
\begin{align*}
    v(t) &= \frac{dx}{dt} = x'(t) && \text{Vitesse (m/s)} \\
    a(t) &= \frac{dv}{dt} = \frac{d^2x}{dt^2} = x''(t) && \text{Accélération (m/s$^2$)}
\end{align*}
La notation $\frac{d}{dt}$ rappelle qu'on dérive \emph{par rapport au temps}. C'est la notation standard en physique et en génie mécanique/électrique.
\end{savaistu}

\courssub{Travail de rédaction : critères d'une copie bien justifiée au Probatoire}

Au Probatoire technique, la note ne sanctionne pas seulement le résultat final, mais la \textbf{qualité de la démarche}. Un exercice de tangente ou de dérivation bien rédigé rapporte des points même en cas d'erreur de calcul mineure.

\begin{methode}[Fiche réflexe « Tangente » — Les 5 étapes obligatoires]
Pour toute question « Déterminer l'équation de la tangente à $\mathcal{C}_f$ en $A(a;f(a))$ » :
\begin{enumerate}
    \item \textbf{Vérifier la dérivabilité :} Écrire « $f$ est dérivable sur $I$ car... » (somme/produit/quotient de fonctions usuelles dérivables) et préciser $a \in D_{f'}$.
    \item \textbf{Calculer $f(a)$ :} Coordonnée $y$ du point de contact $A$.
    \item \textbf{Calculer $f'(x)$ puis $f'(a)$ :} Pente $m$ de la tangente. Montrer le calcul de $f'(x)$ (règles de dérivation) puis l'évaluation en $a$.
    \item \textbf{Écrire l'équation réduite :} $y = f'(a)(x-a) + f(a)$. Développer et réduire pour obtenir la forme $y = mx + p$.
    \item \textbf{Vérifier (le « point final » qui sauve la note) :} Remplacer $x$ par $a$ dans l'équation trouvée ; on doit retomber sur $f(a)$. Écrire : « Vérification : pour $x=a$, $y = \dots = f(a)$. L'équation est correcte. »
\end{enumerate}
\end{methode}

\begin{aretenir}[Critères de rédaction « Probatoire »]
\begin{itemize}
    \item \textbf{Unités :} Si le problème est concret (mécanique, électricité, économie), \textbf{toujours} écrire les unités dans la conclusion (m, m/s, €, Ω, etc.).
    \item \textbf{Phrases complètes :} « La fonction $f$ est dérivable sur $\mathbb{R}$ car... » et non juste « $f$ dérivable sur $\mathbb{R}$ ».
    \item \textbf{Justification du signe :} Pour un tableau de variations, la ligne $f'$ doit être justifiée (factorisation, étude de signe de trinôme, etc.), pas seulement remplie.
    \item \textbf{Conclusion encadrée :} Terminer par « Donc l'équation de la tangente est $\boxed{y = 2x - 6}$ » ou « Le coût minimal est $\boxed{1250 \text{ FCFA}}$ ».
\end{itemize}
\end{aretenir}

\courssub{Exemple de synthèse complète}

\begin{exoresolu}[Synthèse : fonction polynôme, tangente, variations]
Soit la fonction $f$ définie sur $\mathbb{R}$ par $f(x) = x^2 - 4x + 3$.
\begin{enumerate}
    \item Étudier les variations de $f$ et dresser son tableau.
    \item Déterminer l'équation de la tangente $\mathcal{T}$ à la courbe $\mathcal{C}_f$ au point d'abscisse $a = 3$.
    \item Vérifier que $\mathcal{T}$ coupe $\mathcal{C}_f$ en un second point et déterminer ses coordonnées.
\end{enumerate}

\tcblower
\textbf{Solution détaillée}

\textbf{1. Variations de $f$.}
$f$ est un polynôme du second degré, elle est dérivable sur $\mathbb{R}$.
$f'(x) = 2x - 4 = 2(x-2)$.
\begin{itemize}
    \item $f'(x) = 0 \iff x = 2$.
    \item $f'(x) < 0 \iff x < 2$ (décroissante sur $]-\infty; 2]$).
    \item $f'(x) > 0 \iff x > 2$ (croissante sur $[2; +\infty[$).
\end{itemize}
$f(2) = 4 - 8 + 3 = -1$ (minimum global).
\begin{center}
\begin{tabular}{|c|ccccc|}\hline
$x$ & $-\infty$ & & $2$ & & $+\infty$ \\ \hline
$f'(x)$ & & $-$ & $0$ & $+$ & \\ \hline
$f(x)$ & $+\infty$ & $\searrow$ & $-1$ & $\nearrow$ & $+\infty$ \\ \hline
\end{tabular}
\end{center}

\textbf{2. Tangente en $a=3$.}
\begin{itemize}
    \item \textit{Vérification :} $3 \in \mathbb{R} = D_{f'}$. OK.
    \item $f(3) = 9 - 12 + 3 = 0$. Point $A(3;0)$.
    \item $f'(3) = 2(3) - 4 = 2$. Pente $m=2$.
    \item Équation : $y = 2(x-3) + 0 \iff y = 2x - 6$.
    \item \textit{Vérification :} Pour $x=3$, $y = 2\times 3 - 6 = 0 = f(3)$. \textbf{Correct.}
\end{itemize}
$\boxed{\mathcal{T} : y = 2x - 6}$

\textbf{3. Intersection $\mathcal{T} \cap \mathcal{C}_f$ (autre que $A$).}
On résout $f(x) = 2x - 6$.
$x^2 - 4x + 3 = 2x - 6 \iff x^2 - 6x + 9 = 0 \iff (x-3)^2 = 0$.
L'équation a une \textbf{racine double} $x=3$.
\textbf{Conclusion :} La tangente en $A(3;0)$ ne coupe la courbe en \textbf{aucun autre point} (elle est tangente au minimum ? Non, $A$ n'est pas le sommet. Ici $f(x) - (2x-6) = (x-3)^2 \ge 0$, la courbe est \emph{au-dessus} de sa tangente, ce qui est normal pour une fonction convexe).
\end{exoresolu}

\courssub{Auto-évaluation : quiz rapide de 5 questions}

Testez vos acquis en 5 minutes chrono. Les réponses sont dans le corrigé suivant.

\begin{exercice}[Quiz « Spécial Probatoire » — 5 questions]
\begin{enumerate}
    \item Soit $f(x) = 3\sqrt{x} - \frac{2}{x}$ sur $]0;+\infty[$. Écrire $f'(x)$ sous la forme la plus simple.
    \item La courbe $\mathcal{C}_f$ admet une tangente parallèle à l'axe des abscisses au point d'abscisse $x_0$. Que vaut $f'(x_0)$ ? Justifier en une phrase.
    \item On donne $f'(x) = \frac{(x-1)(x+2)}{x^2+1}$ sur $\mathbb{R}$. Dresser le tableau de variations de $f$ (inutile de calculer $f(x)$).
    \item Soit $g(x) = (2x-1)(x^2+3)$. Calculer $g'(x)$ en utilisant la formule du produit (ne pas développer $g$ avant de dériver). Factoriser le résultat.
    \item Vrai ou Faux ? « Si $f'(a) = 0$, alors $f$ admet un extremum en $a$. » Justifier par un contre-exemple si Faux.
\end{enumerate}
\end{exercice}

\begin{corrige}[Quiz « Spécial Probatoire »]
\begin{enumerate}
    \item $f(x) = 3x^{1/2} - 2x^{-1}$.
    $f'(x) = 3 \cdot \frac{1}{2}x^{-1/2} - 2(-1)x^{-2} = \frac{3}{2\sqrt{x}} + \frac{2}{x^2}$.
    \item $f'(x_0) = 0$. Une tangente parallèle à l'axe des abscisses est horizontale, sa pente est nulle, or la pente est $f'(x_0)$.
    \item Dénominateur $x^2+1 > 0$ toujours. Signe de $f'$ = signe du numérateur $(x-1)(x+2)$.
    Zéros : $-2$ et $1$.
    \begin{center}
    \begin{tabular}{|c|ccccccc|}\hline
    $x$ & $-\infty$ & & $-2$ & & $1$ & & $+\infty$ \\ \hline
    $f'(x)$ & & $+$ & $0$ & $-$ & $0$ & $+$ & \\ \hline
    $f$ & $-\infty$ & $\nearrow$ & $f(-2)$ & $\searrow$ & $f(1)$ & $\nearrow$ & $+\infty$ \\ \hline
    \end{tabular}
    \end{center}
    Max en $-2$, min en $1$.
    \item $u(x)=2x-1 \Rightarrow u'=2$. $v(x)=x^2+3 \Rightarrow v'=2x$.
    $g' = u'v + uv' = 2(x^2+3) + (2x-1)(2x) = 2x^2+6 + 4x^2-2x = 6x^2 - 2x + 6 = 2(3x^2 - x + 3)$.
    (Discriminant $\Delta = 1 - 36 < 0$, $g'>0$ partout, $g$ strictement croissante).
    \item \textbf{Faux.} $f'(a)=0$ est une condition \emph{nécessaire} mais pas \emph{suffisante}.
    Contre-exemple : $f(x)=x^3$ en $a=0$. $f'(x)=3x^2$, $f'(0)=0$. Mais $f$ est strictement croissante sur $\mathbb{R}$ (pas d'extremum en $0$, c'est un point d'inflexion à tangente horizontale).
\end{enumerate}
\end{corrige}

\begin{savaistu}[Le Probatoire et la dérivation : les statistiques parlent]
Depuis 10 ans, l'analyse de fonctions (dérivée, tangente, variations) est le \textbf{chapitre le plus fréquent} au Probatoire technique (Séries F, E, STI, etc.). Il apparaît sous forme de question indépendante (calcul de dérivée, équation de tangente) \textbf{ou} imbriqué dans un problème d'optimisation (coût, volume, rendement, mouvement).
\begin{itemize}
    \item \textbf{Sujet type 1 :} « Soit $f(x)=\dots$ Calculer $f'(x)$. En déduire les variations. Tracer $\mathcal{C}_f$. » (4 à 6 points).
    \item \textbf{Sujet type 2 :} « Une entreprise a pour fonction coût $C(x)=\dots$ Déterminer le niveau de production minimisant le coût moyen. » (6 à 8 points).
    \item \textbf{Sujet type 3 :} « Un mobile suit la loi $x(t)=\dots$ Déterminer l'instant où la vitesse est nulle. » (3 à 4 points).
\end{itemize}
Maîtriser parfaitement le tableau des formules, la méthode de la tangente en 5 étapes et le lien signe de $f'$ / variations de $f$, c'est s'assurer \textbf{au minimum 10 à 15 points} sur l'épreuve. C'est le « retour sur investissement » le plus élevé du programme.
\end{savaistu}

\courssub{Vers la Terminale : ce qui vous attend}

En Terminale technique, la dérivation change de dimension :
\begin{itemize}
    \item \textbf{Fonctions composées générales :} $(u \circ v)' = v' \cdot (u' \circ v)$ (règle de la chaîne complète).
    \item \textbf{Nouvelles fonctions usuelles :} Exponentielle $e^x$, Logarithme $\ln x$, Trigonométriques $\sin, \cos, \tan$ et leurs inverses.
    \item \textbf{Équations différentielles :} On ne cherche plus la dérivée d'une fonction connue, on cherche la fonction dont on connaît la dérivée (ex: $y' = ky$ modèle population/radioactivité/RC circuit).
    \item \textbf{Primitives et Intégrales :} L'opération inverse de la dérivation, pour calculer des aires, des travaux, des charges de condensateurs.
\end{itemize}
Tout ce que vous avez construit ici — la rigueur du nombre dérivé, la géométrie de la tangente, l'algèbre des dérivées — en est le \textbf{socle indispensable}. Travaillez vos classiques, soignez votre rédaction, et la Terminale sera une suite naturelle.

\begin{aretenir}[Check-list finale avant l'examen]
\renewcommand{\labelitemi}{$\square$}
\begin{itemize}
    \item Je sais écrire la définition du nombre dérivé $f'(a)$ comme une limite.
    \item Je connais par cœur les 4 dérivées usuelles ($x^n, \sqrt{x}, 1/x, c$) et leur domaine.
    \item J'applique sans hésiter les 4 règles d'opérations ($+$, $-$, $\times$, $/$).
    \item Je sais dresser un tableau de variations complet (lignes $x, f', f$ alignées).
    \item J'applique la méthode des 5 étapes pour la tangente (dont la vérification finale).
    \item Je sais lire le signe de $f'$ sur la courbe de $f$ (et inversement).
    \item Je rédige en phrases, avec unités et conclusion encadrée.
\end{itemize}
\end{aretenir}

\newpage

\coursec{Activité d'intégration}

\coursec{Fonctions dérivées}

\courssub{Approche intuitive du nombre dérivé}

\begin{experience}[Découverte : du taux d'accroissement au nombre dérivé]
Soit une fonction $f$ définie sur un intervalle $I$ et $a \in I$. On considère un point $M(a\,;\,f(a))$ sur la courbe $\mathcal{C}_f$.
\begin{enumerate}
\item On prend un point $N$ sur $\mathcal{C}_f$ d'abscisse $a+h$ ($h \neq 0$). Le taux d'accroissement de $f$ sur $[a\,;\,a+h]$ (ou $[a+h\,;\,a]$) est :
\[ \tau(h) = \frac{f(a+h)-f(a)}{h}. \]
Géométriquement, $\tau(h)$ est le coefficient directeur de la sécante $(MN)$.
\item On fait tendre $h$ vers $0$ : le point $N$ se rapproche de $M$, la sécante $(MN)$ tend vers une position limite : la \cle{tangente} à $\mathcal{C}_f$ en $M$.
\item Si $\tau(h)$ admet une limite finie lorsque $h \to 0$, on appelle cette limite le \cle{nombre dérivé} de $f$ en $a$, noté $f'(a)$.
\end{enumerate}
\[ f'(a) = \lim_{h \to 0} \frac{f(a+h)-f(a)}{h}. \]
\end{experience}

\begin{definition}[Nombre dérivé et fonction dérivée]
Soit $f$ une fonction définie sur un intervalle $I$ et $a \in I$.
\begin{itemize}
\item Si la limite $\lim\limits_{h \to 0} \dfrac{f(a+h)-f(a)}{h}$ existe et est finie, on dit que $f$ est \cle{dérivable} en $a$ et on note $f'(a)$ ce nombre dérivé.
\item Si $f$ est dérivable en tout point de $I$, on définit la \cle{fonction dérivée} $f' : I \to \mathbb{R}$ qui associe à $x$ le nombre dérivé $f'(x)$.
\end{itemize}
\emph{Interprétation géométrique :} $f'(a)$ est le coefficient directeur de la tangente à $\mathcal{C}_f$ au point d'abscisse $a$.
\emph{Interprétation physique/économique :} $f'(a)$ est le taux de variation instantané de $f$ en $a$ (vitesse, gain marginal, etc.).
\end{definition}

\courssub{Tangente à une courbe en un point}

\begin{propriete}[Équation de la tangente -- Démonstration exigible au Probatoire]
Soit $f$ dérivable en $a$. La tangente $\mathcal{T}$ à la courbe $\mathcal{C}_f$ au point $A(a\,;\,f(a))$ a pour équation réduite :
\[ y = f'(a)(x - a) + f(a). \]
\emph{Preuve :} La tangente passe par $A$ et a pour coefficient directeur $f'(a)$. Une droite de coefficient directeur $m$ passant par $(x_0,y_0)$ a pour équation $y - y_0 = m(x - x_0)$, d'où la formule.
\end{propriete}

\begin{methode}[Tracer la tangente en un point]
\begin{enumerate}
\item Calculer $f(a)$ (ordonnée du point de contact) et $f'(a)$ (coefficient directeur).
\item Écrire l'équation réduite $y = f'(a)(x-a) + f(a)$.
\item Pour le tracé : placer le point $A(a\,;\,f(a))$ ; partir de $A$, avancer de $1$ unité en abscisse et de $f'(a)$ unités en ordonnée pour obtenir un second point de la tangente.
\end{enumerate}
\end{methode}

\courssub{Fonction dérivée et dérivées des fonctions usuelles}

\begin{propriete}[Tableau des dérivées usuelles -- À connaître par cœur]
Soient $n \in \mathbb{N}^*$, $k \in \mathbb{R}$, $x > 0$ pour $\sqrt{x}$ et $\ln x$, $x \neq 0$ pour $1/x$.
\[
\begin{array}{|c|c|c|}
\hline
f(x) & f'(x) & \text{Domaine de dérivabilité} \\
\hline
k & 0 & \mathbb{R} \\
x & 1 & \mathbb{R} \\
x^n & n x^{n-1} & \mathbb{R} \ (\text{ou } \mathbb{R}^* \text{ si } n<0) \\
\sqrt{x} & \dfrac{1}{2\sqrt{x}} & ]0\,;\,+\infty[ \\
\dfrac{1}{x} & -\dfrac{1}{x^2} & \mathbb{R}^* \\
\hline
\end{array}
\]
\end{propriete}

\begin{attention}[Pièges classiques sur les puissances]
\begin{itemize}
\item $(x^2)' = 2x$ et non $2$.
\item $(\sqrt{x})' = \frac{1}{2\sqrt{x}}$ et non $\frac{1}{\sqrt{x}}$.
\item Pour $f(x) = \frac{1}{x^2} = x^{-2}$, on utilise la formule des puissances : $f'(x) = -2x^{-3} = -\frac{2}{x^3}$.
\item La dérivée d'une constante est $0$ : $(12)' = 0$, $(-5)' = 0$.
\end{itemize}
\end{attention}

\courssub{Opérations sur les fonctions dérivées}

\begin{propriete}[Linéarité de la dérivation -- Démonstration exigible]
Soient $u$ et $v$ deux fonctions dérivables sur un intervalle $I$, et $k \in \mathbb{R}$.
\begin{itemize}
\item \textbf{Somme :} $(u+v)' = u' + v'$ sur $I$.
\item \textbf{Multiple par une constante :} $(ku)' = k\,u'$ sur $I$.
\item \textbf{Différence :} $(u-v)' = u' - v'$ (cas particulier avec $k=-1$).
\end{itemize}
\emph{Remarque :} Les formules du produit $(uv)' = u'v + uv'$ et du quotient $\left(\frac{u}{v}\right)' = \frac{u'v - uv'}{v^2}$ font partie du programme de Terminale.
\end{propriete}

\begin{methode}[Dériver une fonction polynôme ou rationnelle simple]
\begin{enumerate}
\item Exprimer la fonction comme combinaison linéaire de fonctions usuelles (puissances, racine, inverse).
\item Appliquer la linéarité : dériver terme à terme.
\item Utiliser le tableau des dérivées usuelles pour chaque terme.
\item Simplifier l'expression obtenue.
\end{enumerate}
\end{methode}

\courssub{Applications concrètes de la dérivation}

\textbf{Objectifs de l'activité}\par
À l'issue de cette activité d'intégration, l'élève sera capable de :
\begin{itemize}
\item modéliser un problème économique simple (bénéfice = recette $-$ coût) à l'aide d'une fonction polynôme du second degré ;
\item calculer la fonction dérivée d'un polynôme à l'aide des formules usuelles et de la linéarité ;
\item résoudre l'équation $B'(x)=0$ et interpréter sa solution comme un extremum (maximum) ;
\item déterminer l'équation de la tangente à une courbe en un point et interpréter concrètement son coefficient directeur (gain marginal) ;
\item tracer une parabole et sa tangente sur papier millimétré ;
\item rédiger et justifier une démarche, une réponse et une conclusion argumentée à destination d'un non-spécialiste.
\end{itemize}

\textbf{Situation}\par
Mama Ngo Bell tient un petit stand de beignets haricots très fréquenté, situé près du carrefour Ngo Bell à Douala. Chaque matin, elle se pose la même question : \emph{combien de beignets dois-je préparer pour gagner le plus d'argent possible ?} Si elle en prépare trop peu, elle perd des clients ; si elle en prépare trop, une partie reste invendue et le coût de fabrication (farine, huile, bois de chauffe) dépasse la recette.

Son fils aîné, élève en classe de Première technique, a observé les ventes pendant plusieurs semaines et propose le modèle suivant, où $x$ désigne le nombre de \textbf{centaines de beignets} vendus par jour (avec $0 \leq x \leq 20$) :
\begin{itemize}
\item la recette journalière, en milliers de francs CFA, est donnée par
\[ R(x) = -2x^2 + 40x \ ; \]
\item le coût journalier de fabrication, en milliers de francs CFA, est donné par
\[ C(x) = 4x + 12. \]
\end{itemize}

Mama Ngo Bell veut des réponses précises : quelle quantité de beignets lui rapporte le bénéfice maximal ? Combien gagne-t-elle alors ? Et si elle vendait $500$ beignets ($x = 5$), que vaudrait le gain supplémentaire apporté par une centaine de beignets de plus ?

\textbf{Consignes}\par
Travail en groupes de 4 élèves. Chaque groupe rédige ses réponses sur une feuille qui sera restituée au tableau.

\begin{enumerate}
\item \textbf{Mise en équation du bénéfice.} Le bénéfice journalier est la différence entre la recette et le coût : $B(x) = R(x) - C(x)$. Montrer que $B(x) = -2x^2 + 36x - 12$. Préciser l'ensemble de définition de $B$.
\item \textbf{Recherche du maximum.} Calculer la fonction dérivée $B'(x)$, puis résoudre l'équation $B'(x) = 0$. Étudier le signe de $B'(x)$ sur $[0\,;\,20]$ et dresser le tableau de variations de $B$.
\item \textbf{Conclusion quantitative.} En déduire la quantité de beignets (en unités) qui maximise le bénéfice, ainsi que le bénéfice maximal en francs CFA.
\item \textbf{Tangente et gain marginal.} Déterminer l'équation réduite de la tangente à la courbe de $B$ au point d'abscisse $a = 5$. Interpréter concrètement son coefficient directeur : que représente-t-il pour Mama Ngo Bell lorsqu'elle passe de $500$ à $600$ beignets ?
\item \textbf{Représentation graphique.} Sur papier millimétré, tracer la parabole représentative de $B$ sur $[0\,;\,20]$ (échelle suggérée : $1$ cm pour $2$ centaines en abscisse, $1$ cm pour $20$ milliers de FCFA en ordonnée) ainsi que la tangente obtenue à la question 4.
\item \textbf{Conseil argumenté.} Rédiger un court paragraphe (5 à 8 lignes) à destination de Mama Ngo Bell, en langage simple, qui répond à ses trois questions et justifie chaque réponse par un résultat mathématique.
\end{enumerate}

\textbf{Ressources à mobiliser}\par
\begin{aretenir}[Boîte à outils de l'activité]
\begin{itemize}
\item \textbf{Dérivées usuelles :} $(x^2)' = 2x$, $(x)' = 1$, $(k)' = 0$ pour toute constante $k$.
\item \textbf{Opérations (linéarité) :} $(u+v)' = u'+v'$ et $(ku)' = k\,u'$ pour tout réel $k$.
\item \textbf{Extremum :} si $f'$ s'annule en $c$ en changeant de signe (de $+$ vers $-$), alors $f$ admet un maximum en $c$ ; (de $-$ vers $+$) : minimum.
\item \textbf{Tangente au point d'abscisse $a$ :}
\[ y = f'(a)(x - a) + f(a). \]
\item \textbf{Interprétation du coefficient directeur :} $f'(a)$ mesure la variation approximative de $f$ quand $x$ augmente d'une unité au voisinage de $a$ (notion de \cle{gain marginal} ou \cle{coût marginal} en économie).
\end{itemize}
\end{aretenir}

\textbf{Démarche et corrigé commenté}\par
\begin{exoresolu}[Le prix optimal des beignets de Mama Ngo Bell]
Énoncé : avec $R(x) = -2x^2 + 40x$ et $C(x) = 4x + 12$ (en milliers de FCFA, $x$ en centaines de beignets, $x \in [0\,;\,20]$), déterminer le bénéfice maximal, la tangente à la courbe du bénéfice en $x=5$, et conseiller Mama Ngo Bell.
\tcblower
\textbf{Étape 1 : Expression du bénéfice.}

Le bénéfice est la recette moins le coût :
\begin{align*}
B(x) &= R(x) - C(x) \\
&= (-2x^2 + 40x) - (4x + 12) \\
&= -2x^2 + 40x - 4x - 12 \quad \text{(attention à la distribution du signe $-$) } \\
&= -2x^2 + 36x - 12.
\end{align*}
$B$ est une fonction polynôme, définie et dérivable sur $\mathbb{R}$, donc sur l'intervalle d'étude $[0\,;\,20]$.

\textbf{Étape 2 : Dérivation et recherche du maximum.}

On dérive terme à terme en utilisant la linéarité et le tableau des dérivées usuelles :
\[ B'(x) = -2 \times (x^2)' + 36 \times (x)' - (12)' = -2 \times 2x + 36 \times 1 - 0 = -4x + 36. \]
$B'$ est une fonction affine, définie sur $[0\,;\,20]$.

Résolution de $B'(x) = 0$ :
\[ -4x + 36 = 0 \iff 4x = 36 \iff x = 9. \]

Étude du signe de $B'$ : coefficient directeur $-4 < 0$, donc $B'$ est strictement décroissante.
\begin{itemize}
\item Pour $x \in [0\,;\,9[$, $B'(x) > B'(9) = 0$ (signe $+$).
\item Pour $x \in ]9\,;\,20]$, $B'(x) < B'(9) = 0$ (signe $-$).
\end{itemize}
$B'$ change de signe de $+$ en $-$ en $x=9$ : $B$ admet un \textbf{maximum} en $9$.

Tableau de variations de $B$ sur $[0\,;\,20]$ :
\begin{center}
\begin{tabular}{|c|ccccc|}\hline
$x$ & $0$ & & $9$ & & $20$ \\ \hline
$B'(x)$ & & $+$ & $0$ & $-$ & \\ \hline
$B$ & $-12$ & $\nearrow$ & $150$ & $\searrow$ & $-92$ \\ \hline
\end{tabular}
\end{center}
\emph{Valeurs utiles :} $B(0) = -12$ (coût fixe sans vente), $B(9) = -2(81) + 36(9) - 12 = -162 + 324 - 12 = 150$, $B(20) = -800 + 720 - 12 = -92$.

\textbf{Étape 3 : Bénéfice maximal.}

$B$ admet son maximum en $x = 9$, c'est-à-dire pour $\textbf{900 beignets}$ par jour.
Le bénéfice maximal est $B(9) = 150$ milliers de FCFA, soit $\textbf{150\,000 FCFA}$ par jour.

\textbf{Étape 4 : Tangente en $x = 5$ et gain marginal.}

On calcule l'ordonnée du point de contact et le coefficient directeur :
\[ B(5) = -2(25) + 36(5) - 12 = -50 + 180 - 12 = 118 \quad ; \quad B'(5) = -4(5) + 36 = 16. \]
L'équation réduite de la tangente $\mathcal{T}$ au point d'abscisse $5$ est :
\begin{align*}
y &= B'(5)(x - 5) + B(5) \\
y &= 16(x - 5) + 118 \\
y &= 16x - 80 + 118 \\
y &= 16x + 38.
\end{align*}

\emph{Interprétation concrète :} Le coefficient directeur $B'(5) = 16$ signifie qu'au voisinage de $500$ beignets vendus, \textbf{chaque centaine de beignets supplémentaire rapporte environ $16$ milliers de FCFA de bénéfice en plus}. C'est le \cle{gain marginal} en $x=5$.
\emph{Vérification :} Le gain réel en passant de $500$ à $600$ beignets ($x$ de $5$ à $6$) est $B(6)-B(5) = 156 - 118 = 38$ milliers pour deux centaines, soit $19$ milliers par centaine. L'approximation par la tangente ($16$) est de bon ordre et sous-estime légèrement car la courbe est concave (dérivée décroissante).

\textbf{Étape 5 : Représentation graphique.}

\begin{popfigure}
\begin{tikzpicture}
\begin{axis}[
    width=0.9\linewidth,
    height=7cm,
    axis lines=middle,
    xlabel={$x$ (centaines de beignets)},
    ylabel={$B(x)$ (milliers de FCFA)},
    xmin=-1, xmax=21,
    ymin=-110, ymax=170,
    xtick={0,5,9,20},
    ytick={-12,38,118,150},
    grid=both,
    grid style={popline!40},
    tick label style={font=\small},
    label style={font=\small},
    clip=false
]
% Parabole du bénéfice
\addplot[popcurve, domain=0:20, samples=200, thick] {-2*x^2+36*x-12};
% Tangente en x=5
\addplot[poporange, thick, domain=0:10.5, samples=2] {16*x+38};
% Points remarquables
\addplot[mark=*, popdark, mark size=2.5pt] coordinates {(5,118) (9,150)};
% Légendes des points (optionnel, géré par la légende globale)
\end{axis}
\end{tikzpicture}
\legende{Parabole du bénéfice $B$ (bleu) et tangente au point d'abscisse $5$ (orange). Le sommet $(9\,;\,150)$ correspond au bénéfice maximal. L'ordonnée à l'origine de la tangente est $38$.}
\end{popfigure}

\textbf{Étape 6 : Conseil argumenté à Mama Ngo Bell.}

\emph{Mama, les calculs montrent que ton bénéfice suit une parabole tournée vers le bas : il augmente tant que tu vends moins de $900$ beignets, puis diminue au-delà, car les invendus coûtent plus cher qu'ils ne rapportent. Prépare donc environ \textbf{900 beignets} par jour : tu gagneras jusqu'à \textbf{150\,000 FCFA}. Si tu n'en vends que 500, chaque centaine de beignets en plus te rapporte encore environ \textbf{16\,000 FCFA} (gain marginal), donc n'hésite pas à produire davantage. En revanche, au-delà de 900, réduis la production : à 2\,000 beignets, tu perdrais 92\,000 FCFA !}
\end{exoresolu}

\begin{attention}[Erreurs fréquentes à signaler lors de la restitution]
\begin{itemize}
\item \textbf{Signe moins oublié :} dans $R(x) - C(x)$, écrire $-4x + 12$ au lieu de $-4x - 12$.
\item \textbf{Confusion abscisse / valeur :} donner $x=9$ comme « bénéfice maximal » au lieu de $B(9)=150$. L'équation $B'(x)=0$ donne l'abscisse de l'extremum, pas sa valeur.
\item \textbf{Formule de la tangente :} écrire $y = f'(a)x + f(a)$ au lieu de $y = f'(a)(x-a) + f(a)$. Vérifier que la tangente passe bien par le point $(a\,;\,f(a))$.
\item \textbf{Conversions d'unités :} oublier que $x=9$ signifie $900$ beignets (centaines) et que $B=150$ signifie $150\,000$ FCFA (milliers).
\item \textbf{Tableau de variations :} mettre les flèches ($\nearrow$, $\searrow$) sur la ligne de $B'$ au lieu de la ligne de $B$, ou inverser les signes.
\end{itemize}
\end{attention}

\textbf{Bilan de l'activité}\par
Cette activité a permis de mobiliser l'ensemble des savoirs de la leçon dans une situation authentique de la vie économique camerounaise :
\begin{itemize}
\item la \cle{modélisation} : traduire « bénéfice = recette $-$ coût » par une fonction polynôme du second degré ;
\item la \cle{dérivation} : calculer $B'(x)$ avec les formules usuelles ($x^2$, $x$, constante) et la linéarité ;
\item l'\cle{optimisation} : l'annulation de la dérivée avec changement de signe ($+ \to -$) caractérise le maximum, ce qui donne un sens concret au tableau de variations ;
\item la \cle{tangente} : son équation et surtout l'interprétation de $f'(a)$ comme gain marginal, notion centrale en gestion et en économie ;
\item la \cle{rédaction justifiée} : formuler un conseil argumenté, compétence transversale exigée au Probatoire et au Baccalauréat technique.
\end{itemize}
La confrontation des démarches au tableau a montré qu'un même résultat mathématique ($x = 9$, $B(9) = 150$) peut être obtenu par des rédactions différentes, mais qu'il n'a de valeur pour l'utilisateur final que s'il est correctement interprété et exprimé dans les bonnes unités (beignets, FCFA).

\courssub{Synthèse, compléments et vers la Terminale}

\begin{aretenir}[Synthèse de la leçon -- À retenir pour l'examen]
\begin{itemize}
\item \textbf{Nombre dérivé :} $f'(a) = \lim\limits_{h \to 0} \frac{f(a+h)-f(a)}{h}$ (coefficient directeur de la tangente, taux de variation instantané).
\item \textbf{Fonction dérivée :} $f'(x)$ obtenue par dérivation terme à terme grâce à la linéarité et au tableau des dérivées usuelles ($x^n$, $\sqrt{x}$, $1/x$, $k$).
\item \textbf{Tangente en $a$ :} $y = f'(a)(x-a) + f(a)$.
\item \textbf{Variations / Extrema :} $f' > 0 \Rightarrow f$ croissante ; $f' < 0 \Rightarrow f$ décroissante ; $f'$ s'annule en changeant de signe $\Rightarrow$ extremum.
\item \textbf{Applications :} Optimisation (max/min), gain marginal / coût marginal ($f'(a)$), modélisation économique (bénéfice, recette, coût).
\end{itemize}
\end{aretenir}

\begin{savaistu}[Vers la Terminale technique]
En Terminale, vous approfondirez :
\begin{itemize}
\item La dérivation des fonctions \textbf{produit} $(uv)' = u'v + uv'$ et \textbf{quotient} $\left(\frac{u}{v}\right)' = \frac{u'v - uv'}{v^2}$.
\item Les fonctions \textbf{exponentielle} ($e^x$, $a^x$), \textbf{logarithme} ($\ln x$) et \textbf{trigonométriques} ($\sin$, $\cos$, $\tan$) et leurs dérivées.
\item L'étude de fonctions plus complexes (asymptotes, convexité via $f''$) et l'optimisation sous contrainte.
\item Les \textbf{équations différentielles} simples ($y' = ay$, $y' = f(x)$) pour modéliser la croissance, la décroissance radioactive, les circuits RC/RL en électricité.
\end{itemize}
Les bases posées ici (définition, tangente, variations, polynômes) sont le socle indispensable pour aborder ces notions supérieures.
\end{savaistu}

\newpage

\coursec{Méthodes et exercices résolus}

\coursec{Fonctions dérivées}

\courssub{Approche intuitive du nombre dérivé}

\begin{definition}[Taux d'accroissement et nombre dérivé en un point]
Soit $f$ une fonction définie sur un intervalle $I$ et $a \in I$.
On appelle \textbf{taux d'accroissement} de $f$ en $a$ (pour un écart $h \neq 0$ tel que $a+h \in I$) le quotient :
\[\tau(h) = \frac{f(a+h)-f(a)}{h}.\]
Geométriquement, $\tau(h)$ est le coefficient directeur de la sécante passant par les points $A(a\,;f(a))$ et $M(a+h\,;f(a+h))$ de la courbe $\mathcal{C}_f$.

Si $\tau(h)$ admet une limite finie lorsque $h$ tend vers $0$, on dit que $f$ est \textbf{dérivable en $a$} et on appelle \textbf{nombre dérivé de $f$ en $a$} cette limite, notée $f'(a)$ :
\[f'(a) = \lim_{h \to 0} \frac{f(a+h)-f(a)}{h}.\]
$f'(a)$ est le coefficient directeur de la tangente à $\mathcal{C}_f$ au point $A(a\,;f(a))$.
\end{definition}

\begin{popfigure}
\begin{tikzpicture}[scale=0.8, >=Stealth]
\draw[->] (-0.5,0) -- (5,0) node[right] {$x$};
\draw[->] (0,-0.5) -- (0,4.5) node[above] {$y$};
\draw[popcurve, thick, domain=0:4.2, samples=100] plot (\x, {0.3*(\x-1)*(\x-1)+0.5});
\node at (4.3, 3.5) {$\mathcal{C}_f$};
\coordinate (A) at (1, 0.5);
\coordinate (M1) at (3, 1.7);
\coordinate (M2) at (2, 0.9);
\coordinate (M3) at (1.5, 0.625);
\fill[popblue] (A) circle (2pt) node[below left] {$A(a,f(a))$};
\fill[poporange] (M1) circle (2pt) node[above right] {$M_1(a+h_1)$};
\fill[poporange!70!popblue] (M2) circle (2pt) node[above right] {$M_2(a+h_2)$};
\fill[poporange!40!popblue] (M3) circle (2pt) node[above right] {$M_3(a+h_3)$};
\draw[popline, dashed] (A) -- (M1);
\draw[popline, dashed] (A) -- (M2);
\draw[popline, dashed] (A) -- (M3);
\draw[popblue, thick] (0, -0.1) -- (4.5, 3.65) node[above right] {Tangente $T$};
\draw[|<->|, popmuted] (1, -0.4) -- (3, -0.4) node[midway, below] {$h_1$};
\draw[|<->|, popmuted] (1, -0.7) -- (2, -0.7) node[midway, below] {$h_2$};
\draw[|<->|, popmuted] (1, -1.0) -- (1.5, -1.0) node[midway, below] {$h_3 \to 0$};
\end{tikzpicture}
\legende{Approche géométrique : lorsque $h \to 0$, la sécante $(AM)$ tend vers la tangente $T$ en $A$.}
\end{popfigure}

\begin{methode}[Calculer $f'(a)$ à l'aide du taux d'accroissement]
Pour calculer le nombre dérivé d'une fonction $f$ en un point $a$ :
\begin{enumerate}
\item Écrire le taux d'accroissement $\tau(h)=\dfrac{f(a+h)-f(a)}{h}$ pour $h \neq 0$.
\item Calculer $f(a+h)$ en développant soigneusement (identités remarquables).
\item Calculer le numérateur $f(a+h)-f(a)$ et factoriser par $h$.
\item Simplifier par $h$ (le facteur $h$ doit \textbf{toujours} se simplifier, sinon il y a une erreur de calcul).
\item Faire tendre $h$ vers $0$ : remplacer $h$ par $0$ dans l'expression simplifiée.
\item Conclure : $f'(a)$ est la limite obtenue.
\end{enumerate}
\textbf{Réflexe :} si la limite n'existe pas (division par $0$ impossible à lever, forme indéterminée non levée, limite infinie), $f$ n'est pas dérivable en $a$.
\end{methode}

\begin{exoresolu}[Nombre dérivé de la fonction carré en $a=3$]
Soit $f$ la fonction définie sur $\mathbb{R}$ par $f(x)=x^{2}$.
\begin{enumerate}
\item Calculer le taux d'accroissement de $f$ en $3$.
\item En déduire $f'(3)$ et interpréter graphiquement le résultat.
\end{enumerate}
\tcblower
\textbf{1. Taux d'accroissement.} Pour $h\neq 0$ :
\begin{align*}
\tau(h)&=\dfrac{f(3+h)-f(3)}{h}=\dfrac{(3+h)^{2}-9}{h}\\
&=\dfrac{9+6h+h^{2}-9}{h}=\dfrac{6h+h^{2}}{h}=\dfrac{h(6+h)}{h}=6+h.
\end{align*}
On a développé $(3+h)^2$ avec l'identité remarquable $(a+b)^2=a^2+2ab+b^2$, puis factorisé par $h$ pour pouvoir simplifier.

\textbf{2. Nombre dérivé.} Quand $h$ tend vers $0$, $6+h$ tend vers $6$. Donc :
\[f'(3)=\lim_{h\to 0}(6+h)=6.\]
\textbf{Interprétation :} la tangente à la courbe de $f$ au point d'abscisse $3$ a pour coefficient directeur $6$. Concrètement, si $f(x)$ représente l'aire (en cm$^2$) d'une plaque carrée de côté $x$ (en cm), alors au voisinage de $x=3$, l'aire augmente d'environ $6$ cm$^2$ par centimètre de côté ajouté.
\end{exoresolu}

\courssub{Tangente à une courbe en un point}

\begin{propriete}[Équation de la tangente]
Soit $f$ une fonction dérivable en $a$. La tangente $T$ à la courbe $\mathcal{C}_f$ au point $A(a\,;f(a))$ a pour équation :
\[y = f'(a)(x-a) + f(a).\]
\textbf{Démonstration exigible au Probatoire :} La tangente est une droite de coefficient directeur $f'(a)$ passant par $A(a\,;f(a))$. Une droite d'équation $y=mx+p$ passant par $A$ vérifie $f(a)=ma+p \iff p=f(a)-ma$. Avec $m=f'(a)$, on obtient $y=f'(a)x + f(a)-f'(a)a = f'(a)(x-a)+f(a)$.
\end{propriete}

\begin{methode}[Déterminer l'équation de la tangente en un point]
Pour déterminer l'équation de la tangente à la courbe de $f$ au point $A$ d'abscisse $a$ :
\begin{enumerate}
\item Vérifier que $f$ est dérivable en $a$ (appartenance au domaine de dérivabilité).
\item Calculer $f(a)$ (ordonnée du point de contact) et $f'(a)$ (coefficient directeur).
\item Appliquer la formule : $y=f'(a)(x-a)+f(a)$.
\item Développer et réduire pour obtenir la forme réduite $y=mx+p$.
\item Conclure par une phrase : « La tangente à $\mathcal{C}_f$ au point $A(a\,;f(a))$ a pour équation $y=\ldots$ ».
\end{enumerate}
\textbf{Réflexe de vérification :} la tangente doit passer par $A$ : remplacer $x$ par $a$ dans l'équation doit redonner $f(a)$.
\end{methode}

\begin{exoresolu}[Tangente à une parabole]
Soit $f$ la fonction définie sur $\mathbb{R}$ par $f(x)=x^{2}-4x+3$ et $\mathcal{C}_f$ sa courbe représentative.
\begin{enumerate}
\item Déterminer l'équation de la tangente $T$ à $\mathcal{C}_f$ au point $A$ d'abscisse $1$.
\item En quel point la tangente à $\mathcal{C}_f$ est-elle horizontale ?
\end{enumerate}
\tcblower
\textbf{1.} $f$ est un polynôme, dérivable sur $\mathbb{R}$, avec $f'(x)=2x-4$.

Calculs au point d'abscisse $a=1$ :
\[f(1)=1-4+3=0 \qquad \text{et} \qquad f'(1)=2-4=-2.\]
La tangente a pour équation :
\[y=f'(1)(x-1)+f(1)=-2(x-1)+0=-2x+2.\]
\textbf{Vérification :} pour $x=1$, $y=-2+2=0=f(1)$ : la tangente passe bien par $A(1\,;0)$.

\textbf{Conclusion :} la tangente à $\mathcal{C}_f$ au point $A(1\,;0)$ a pour équation $y=-2x+2$.

\textbf{2.} Une tangente est horizontale lorsque son coefficient directeur est nul, c'est-à-dire $f'(x)=0$ :
\[2x-4=0 \iff x=2.\]
Or $f(2)=4-8+3=-1$. \textbf{Conclusion :} la tangente est horizontale au point $B(2\,;-1)$, qui est le sommet de la parabole.
\end{exoresolu}

\courssub{Fonction dérivée et dérivées des fonctions usuelles}

\begin{definition}[Fonction dérivée]
Soit $f$ une fonction définie sur un intervalle $I$. Si $f$ est dérivable en tout point de $I$ (ou de $I$ privé de ses bornes éventuelles), on appelle \textbf{fonction dérivée} de $f$, notée $f'$, la fonction qui associe à tout $x \in I$ le nombre dérivé $f'(x)$.
\end{definition}

\begin{propriete}[Tableau des dérivées usuelles (à connaître par cœur)]
Soient $n \in \mathbb{N}^*$, $k \in \mathbb{R}$ constante, $u$ une fonction dérivable.
\begin{center}
\begin{tabularx}{\linewidth}{|l|X|X|}\hline
\textbf{Fonction $f(x)$} & \textbf{Dérivée $f'(x)$} & \textbf{Domaine de dérivabilité} \\ \hline
$k$ (constante) & $0$ & $\mathbb{R}$ \\ \hline
$x^n$ & $n x^{n-1}$ & $\mathbb{R}$ \\ \hline
$\dfrac{1}{x} = x^{-1}$ & $-\dfrac{1}{x^2}$ & $\mathbb{R}^*$ \\ \hline
$\sqrt{x} = x^{1/2}$ & $\dfrac{1}{2\sqrt{x}}$ & $]0\,;+\infty[$ \\ \hline
$u(x)+v(x)$ & $u'(x)+v'(x)$ & $D_u \cap D_v$ \\ \hline
$k\,u(x)$ & $k\,u'(x)$ & $D_u$ \\ \hline
$u(x)v(x)$ & $u'(x)v(x)+u(x)v'(x)$ & $D_u \cap D_v$ \\ \hline
$\dfrac{u(x)}{v(x)}$ & $\dfrac{u'(x)v(x)-u(x)v'(x)}{[v(x)]^2}$ & $D_u \cap D_v \setminus \{x \mid v(x)=0\}$ \\ \hline
\end{tabularx}
\end{center}
\textbf{Attention :} Pour la racine carrée, la dérivée n'existe pas en $0$ (pente infinie / tangente verticale). Le domaine est $]0,+\infty[$.
\end{propriete}

\begin{methode}[Calculer une fonction dérivée]
Pour dériver une fonction construite à partir des fonctions usuelles :
\begin{enumerate}
\item Identifier la structure : somme, produit par une constante, produit, quotient.
\item Appliquer la bonne formule de dérivation (linéarité, produit, quotient).
\item Utiliser les dérivées usuelles du tableau ci-dessus.
\item Simplifier le résultat (factoriser, réduire au même dénominateur).
\item \textbf{Préciser toujours l'ensemble sur lequel $f$ est dérivable} (domaine de définition de $f'$).
\end{enumerate}
\end{methode}

\begin{exoresolu}[Dérivées de fonctions polynômes et rationnelles]
Déterminer la fonction dérivée de chacune des fonctions suivantes, en précisant l'ensemble de dérivabilité :
\[f(x)=3x^{3}-5x^{2}+2x-7 \qquad g(x)=(2x+1)(x-3) \qquad h(x)=\dfrac{2x+1}{x-4}\]
\tcblower
\textbf{Dérivée de $f$ (somme).} $f$ est un polynôme, dérivable sur $\mathbb{R}$.
\[f'(x)=3\times 3x^{2}-5\times 2x+2\times 1-0=9x^{2}-10x+2 \quad \text{sur } \mathbb{R}.\]

\textbf{Dérivée de $g$ (produit).} $g$ est dérivable sur $\mathbb{R}$. Posons $u(x)=2x+1$ et $v(x)=x-3$ ; alors $u'(x)=2$ et $v'(x)=1$.
\begin{align*}
g'(x)&=u'(x)v(x)+u(x)v'(x)=2(x-3)+(2x+1)\times 1\\
&=2x-6+2x+1=4x-5 \quad \text{sur } \mathbb{R}.
\end{align*}
\textbf{Vérification :} en développant, $g(x)=2x^{2}-5x-3$, donc $g'(x)=4x-5$. Les deux méthodes concordent.

\textbf{Dérivée de $h$ (quotient).} $h$ est définie et dérivable sur $\mathbb{R}\setminus\{4\}$. Posons $u(x)=2x+1$ et $v(x)=x-4$ ; $u'(x)=2$, $v'(x)=1$.
\begin{align*}
h'(x)&=\dfrac{u'(x)v(x)-u(x)v'(x)}{[v(x)]^{2}}=\dfrac{2(x-4)-(2x+1)\times 1}{(x-4)^{2}}\\
&=\dfrac{2x-8-2x-1}{(x-4)^{2}}=\dfrac{-9}{(x-4)^{2}} \quad \text{sur } \mathbb{R}\setminus\{4\}.
\end{align*}
\textbf{Conclusion :} $f'(x)=9x^{2}-10x+2$ sur $\mathbb{R}$ ; $g'(x)=4x-5$ sur $\mathbb{R}$ ; $h'(x)=\dfrac{-9}{(x-4)^{2}}$ sur $\mathbb{R}\setminus\{4\}$.
\end{exoresolu}

\courssub{Opérations sur les fonctions dérivées}

\begin{propriete}[Règles de dérivation (Linéalité, Produit, Quotient)]
Soient $u$ et $v$ deux fonctions dérivables sur un intervalle $I$.
\begin{itemize}
\item \textbf{Somme :} $u+v$ est dérivable sur $I$ et $(u+v)' = u' + v'$.
\item \textbf{Produit par une constante :} $ku$ est dérivable sur $I$ et $(ku)' = k u'$.
\item \textbf{Produit :} $uv$ est dérivable sur $I$ et $(uv)' = u'v + uv'$.
\item \textbf{Quotient :} Si $v(x) \neq 0$ sur $I$, $\frac{u}{v}$ est dérivable sur $I$ et $\left(\frac{u}{v}\right)' = \frac{u'v - uv'}{v^2}$.
\end{itemize}
Ces formules s'appliquent point par point pour obtenir la fonction dérivée.
\end{propriete}

\begin{attention}[Pièges classiques sur les opérations]
\begin{enumerate}
\item \textbf{Produit :} $(uv)' \neq u'v'$ ! Exemple : $f(x)=x \cdot x = x^2$. $f'(x)=2x$. Mais $u'v' = 1 \times 1 = 1 \neq 2x$.
\item \textbf{Quotient :} $\left(\frac{u}{v}\right)' \neq \frac{u'}{v'}$. L'ordre au numérateur est crucial : $u'v - uv'$ (dérivée du numérateur $\times$ dénominateur \textbf{moins} numérateur $\times$ dérivée du dénominateur).
\item \textbf{Domaine :} Pour un quotient, on exclut \textbf{toujours} les valeurs annulant le dénominateur $v(x)$ de l'ensemble de dérivabilité.
\end{enumerate}
\end{attention}

\courssub{Applications concrètes de la dérivation}

\begin{propriete}[Signe de la dérivée et sens de variation]
Soit $f$ une fonction dérivable sur un intervalle $I$.
\begin{itemize}
\item Si $f'(x) > 0$ sur $I$, alors $f$ est \textbf{strictement croissante} sur $I$.
\item Si $f'(x) < 0$ sur $I$, alors $f$ est \textbf{strictement décroissante} sur $I$.
\item Si $f'(x) = 0$ sur $I$, alors $f$ est \textbf{constante} sur $I$.
\end{itemize}
Réciproquement, si $f$ est strictement monotone, le signe de $f'$ est constant (mais peut s'annuler en des points isolés, ex: $x^3$ en $0$).
\end{propriete}

\begin{methode}[Étudier les variations d'une fonction grâce à sa dérivée]
Pour dresser le tableau de variations d'une fonction $f$ sur un intervalle $I$ :
\begin{enumerate}
\item Calculer $f'(x)$ et préciser l'ensemble de dérivabilité (coïncide avec $I$ pour les polynômes).
\item Résoudre l'équation $f'(x)=0$ dans $I$ et étudier le signe de $f'(x)$ (tableau de signes).
\item Appliquer le théorème : $f'>0 \Rightarrow f$ croissante ; $f'<0 \Rightarrow f$ décroissante.
\item Calculer les valeurs remarquables $f(x)$ aux bornes de $I$ et aux points où $f'$ s'annule (extrema).
\item Dresser le tableau de variations complet : \textbf{les trois lignes ($x$, $f'(x)$, $f$) ont EXACTEMENT le même nombre de colonnes}.
\end{enumerate}
\end{methode}

\begin{exoresolu}[Variations d'une fonction polynôme]
Soit $f$ la fonction définie sur $[-2\,;4]$ par $f(x)=x^{2}-2x-3$.
\begin{enumerate}
\item Calculer $f'(x)$ et étudier son signe.
\item Dresser le tableau de variations de $f$ sur $[-2\,;4]$.
\item En déduire le minimum de $f$ sur cet intervalle.
\end{enumerate}
\tcblower
\textbf{1.} $f$ est dérivable sur $[-2\,;4]$ et $f'(x)=2x-2=2(x-1)$.

$f'(x)$ est une fonction affine qui s'annule en $x=1$, négative avant $1$ et positive après :
\begin{center}
\begin{tabular}{|c|ccccc|}\hline
$x$ & $-2$ & & $1$ & & $4$ \\ \hline
$f'(x)$ & & $-$ & $0$ & $+$ & \\ \hline
\end{tabular}
\end{center}

\textbf{2.} Valeurs remarquables : $f(-2)=4+4-3=5$ ; $f(1)=1-2-3=-4$ ; $f(4)=16-8-3=5$.
\begin{center}
\begin{tabular}{|c|ccccc|}\hline
$x$ & $-2$ & & $1$ & & $4$ \\ \hline
$f'(x)$ & & $-$ & $0$ & $+$ & \\ \hline
$f$ & $5$ & $\searrow$ & $-4$ & $\nearrow$ & $5$ \\ \hline
\end{tabular}
\end{center}

\textbf{3.} D'après le tableau, $f$ décroît de $5$ à $-4$ puis croît de $-4$ à $5$. \textbf{Conclusion :} le minimum de $f$ sur $[-2\,;4]$ est $-4$, atteint en $x=1$.
\end{exoresolu}

\begin{methode}[Résoudre un problème concret d'optimisation]
Pour résoudre un problème concret (coût minimal, bénéfice maximal, aire maximale, volume maximal) :
\begin{enumerate}
\item Identifier la grandeur à optimiser et la variable indépendante ; modéliser par une fonction $f$ sur un intervalle $I$ adapté au contexte (contraintes physiques : longueurs positives, quantités entières, capacité maximale...).
\item Calculer $f'(x)$ et résoudre $f'(x)=0$ dans $I$.
\item Étudier le signe de $f'$ (tableau de variations) pour vérifier qu'il s'agit bien d'un maximum ou d'un minimum (changement de signe de $f'$).
\item Calculer la valeur optimale et répondre à la question posée avec une phrase complète et les unités.
\end{enumerate}
\textbf{Réflexe :} toujours vérifier que la solution trouvée appartient à l'intervalle $I$ imposé par le contexte.
\end{methode}

\begin{exoresolu}[Bénéfice maximal d'un atelier de Douala]
Un atelier de confection de Douala produit chaque jour $q$ sacs, avec $q\in[0\,;50]$. Le bénéfice journalier, en milliers de francs CFA, est modélisé par :
\[B(q)=-q^{2}+40q-100.\]
\begin{enumerate}
\item Étudier les variations de $B$ sur $[0\,;50]$.
\item En déduire le nombre de sacs à produire pour un bénéfice maximal, et la valeur de ce bénéfice.
\end{enumerate}
\tcblower
\textbf{1.} $B$ est un polynôme, dérivable sur $[0\,;50]$, et $B'(q)=-2q+40=2(20-q)$.

$B'(q)=0 \iff q=20$. De plus $B'(q)>0$ pour $q<20$ et $B'(q)<0$ pour $q>20$.

Valeurs : $B(0)=-100$ ; $B(20)=-400+800-100=300$ ; $B(50)=-2500+2000-100=-600$.
\begin{center}
\begin{tabular}{|c|ccccc|}\hline
$q$ & $0$ & & $20$ & & $50$ \\ \hline
$B'(q)$ & & $+$ & $0$ & $-$ & \\ \hline
$B$ & $-100$ & $\nearrow$ & $300$ & $\searrow$ & $-600$ \\ \hline
\end{tabular}
\end{center}

\textbf{2.} D'après le tableau, $B$ admet un maximum en $q=20$, qui appartient bien à $[0\,;50]$.

\textbf{Conclusion :} l'atelier doit produire $20$ sacs par jour ; le bénéfice maximal est alors de $300$ milliers de francs CFA, soit $300\,000$ FCFA par jour.
\end{exoresolu}

\begin{exoresolu}[Position relative d'une courbe et de sa tangente]
Soit $f$ la fonction définie sur $\mathbb{R}$ par $f(x)=x^{2}$ et $T$ la tangente à sa courbe au point d'abscisse $1$.
\begin{enumerate}
\item Déterminer l'équation de $T$.
\item Étudier le signe de $d(x)=f(x)-(2x-1)$ et conclure sur la position de la courbe par rapport à $T$.
\end{enumerate}
\tcblower
\textbf{1.} Calculons la dérivée de $f$ : $f'(x)=2x$ pour tout réel $x$. Au point d'abscisse $a=1$ :
\[f(1)=1 \qquad \text{et} \qquad f'(1)=2.\]
D'après la formule de la tangente, $T$ a pour équation :
\[y=f'(1)(x-1)+f(1)=2(x-1)+1=2x-1.\]

\textbf{2.} Étudions le signe de la différence $d(x)=f(x)-(2x-1)$ :
\[d(x)=x^{2}-2x+1=(x-1)^{2}.\]
Or un carré est toujours positif ou nul : pour tout réel $x$, $d(x)\geq 0$, et $d(x)=0$ uniquement pour $x=1$.

Par suite, $f(x)\geq 2x-1$ pour tout réel $x$.

\textbf{Conclusion :} la courbe de $f$ est toujours située au-dessus de sa tangente $T$, et la touche uniquement au point de contact $A(1\,;1)$. C'est la propriété de convexité de la parabole (qui sera formalisée en Terminale par $f''(x) \ge 0$).
\end{exoresolu}

\courssub{Synthèse, compléments et vers la Terminale}

\begin{aretenir}[Synthèse de la leçon : Fonctions dérivées]
\begin{itemize}
\item \textbf{Nombre dérivé $f'(a)$ :} limite du taux d'accroissement $\frac{f(a+h)-f(a)}{h}$ quand $h\to 0$. Coefficient directeur de la tangente en $a$.
\item \textbf{Fonction dérivée $f'$ :} fonction qui associe $f'(x)$ à chaque $x$ du domaine de dérivabilité.
\item \textbf{Formules usuelles :} $(x^n)'=nx^{n-1}$, $(1/x)'=-1/x^2$, $(\sqrt{x})'=1/(2\sqrt{x})$, $(k)'=0$.
\item \textbf{Règles :} $(u+v)'=u'+v'$, $(ku)'=ku'$, $(uv)'=u'v+uv'$, $(u/v)'=(u'v-uv')/v^2$.
\item \textbf{Tangente en $a$ :} $y = f'(a)(x-a) + f(a)$.
\item \textbf{Variations :} $f'>0 \iff f$ croissante ; $f'<0 \iff f$ décroissante. Extrema aux points où $f'$ change de signe.
\item \textbf{Optimisation :} Modéliser $\to$ Dériver $\to$ Étudier le signe $\to$ Conclure avec unités.
\end{itemize}
\end{aretenir}

\begin{savaistu}[Vers la Terminale : Dérivées des fonctions composées et réciproques]
En Première technique, on dérive des fonctions \textbf{explicites} construites par opérations algébriques (somme, produit, quotient) à partir de $x^n$, $1/x$, $\sqrt{x}$.
En Terminale, le programme s'élargit aux :
\begin{itemize}
\item \textbf{Fonctions composées :} $(u \circ v)'(x) = u'(v(x)) \cdot v'(x)$ (Règle de la chaîne). Ex: $\sqrt{3x+1}$, $(x^2+1)^3$.
\item \textbf{Fonctions réciproques :} Si $f$ est bijective et dérivable, $(f^{-1})'(y) = \frac{1}{f'(x)}$ avec $y=f(x)$.
\item \textbf{Nouvelles fonctions usuelles :} Exponentielle $e^x$, Logarithme $\ln x$, Trigonométriques $\sin x$, $\cos x$.
\item \textbf{Seconde dérivée $f''$ :} Convexité/Concavité, points d'inflexion, étude complète de fonctions.
\end{itemize}
Les bases solides posées ici (calcul algébrique rigoureux, tableau de variations, rédaction) sont la clé de la réussite en Terminale.
\end{savaistu}

\begin{attention}[Les 5 erreurs qui coûtent des points au Probatoire]
\begin{enumerate}
\item \textbf{Oublier de simplifier par $h$ avant la limite.} Dans le calcul du nombre dérivé, on ne peut pas remplacer $h$ par $0$ dans $\dfrac{f(a+h)-f(a)}{h}$ tant que le facteur $h$ n'a pas été simplifié : on obtiendrait $\dfrac{0}{0}$.
\item \textbf{Confondre $f(a)$ et $f'(a)$ dans l'équation de la tangente.} La formule est $y=f'(a)(x-a)+f(a)$ : $f'(a)$ est le coefficient directeur, $f(a)$ est l'ordonnée du point de contact. Intervertir les deux est l'erreur la plus fréquente.
\item \textbf{Dériver un produit terme à terme.} $(uv)'\neq u'v'$ ! Il faut appliquer $(uv)'=u'v+uv'$. De même, $\left(\dfrac{u}{v}\right)'\neq\dfrac{u'}{v'}$.
\item \textbf{Conclure sur les variations sans étudier le signe de $f'$.} Affirmer « $f$ est croissante » sans tableau de signes de $f'$ ni référence au théorème coûte des points de justification. Le signe de $f'$ doit être établi proprement.
\item \textbf{Oublier le domaine et le contexte.} Pour un quotient, exclure les valeurs qui annulent le dénominateur (ex. $x\neq 4$). Dans un problème concret, vérifier que la solution est dans l'intervalle et donner la réponse avec l'unité (FCFA, sacs, cm\ldots).
\end{enumerate}
\end{attention}

\newpage

\coursec{Exercices}

\coursec{Leçon 04 : Fonctions dérivées}

\courssub{Approche intuitive du nombre dérivé}

\begin{definition}[Taux d'accroissement]
Soit $f$ une fonction définie sur un intervalle $I$ et $a \in I$. Pour tout $h \neq 0$ tel que $a+h \in I$, on appelle \textbf{taux d'accroissement} de $f$ sur $[a; a+h]$ (ou $[a+h; a]$) le quotient :
\[
\tau(h) = \frac{f(a+h) - f(a)}{h}
\]
Ce quotient représente la pente de la sécante passant par les points $A(a; f(a))$ et $B(a+h; f(a+h))$ de la courbe $\mathcal{C}_f$.
\end{definition}

\begin{propriete}[Nombre dérivé en un point]
Si la fonction $f$ admet une limite finie lorsque $h \to 0$ pour le taux d'accroissement, on dit que $f$ est \textbf{dérivable en $a$}. Cette limite est appelée \textbf{nombre dérivé de $f$ en $a$} et se note $f'(a)$ :
\[
f'(a) = \lim_{h \to 0} \frac{f(a+h) - f(a)}{h}
\]
Géométriquement, $f'(a)$ est la pente de la tangente à $\mathcal{C}_f$ au point d'abscisse $a$.
\end{propriete}

\begin{methode}[Calculer un nombre dérivé par la définition]
\begin{enumerate}
    \item Écrire l'expression de $f(a+h)$.
    \item Calculer la différence $f(a+h) - f(a)$ et factoriser par $h$ au numérateur.
    \item Simplifier par $h \neq 0$ pour obtenir une expression sans dénominateur $h$.
    \item Faire tendre $h$ vers $0$ pour obtenir $f'(a)$.
\end{enumerate}
\end{methode}

\begin{exemplebox}[Nombre dérivé de $x \mapsto x^2$ en $a=3$]
Soit $f(x)=x^2$. Calculons $f'(3)$.
\begin{align*}
f(3+h) &= (3+h)^2 = 9 + 6h + h^2 \\
f(3+h) - f(3) &= (9+6h+h^2) - 9 = 6h + h^2 = h(6+h) \\
\frac{f(3+h)-f(3)}{h} &= 6+h \quad (h \neq 0) \\
f'(3) &= \lim_{h \to 0} (6+h) = 6
\end{align*}
La tangente en $x=3$ a pour pente $6$.
\end{exemplebox}

\begin{attention}[Pièges fréquents]
\begin{itemize}
    \item Ne pas écrire $\frac{f(a+h)-f(a)}{h}$ pour $h=0$ (division par zéro interdite). On simplifie \textbf{avant} de faire tendre $h$ vers $0$.
    \item Vérifier que $f$ est bien définie au voisinage de $a$ (ex: $\sqrt{x}$ non dérivable en $0$ car non définie pour $h<0$).
    \item La dérivée n'existe pas si la limite est infinie (tangente verticale) ou si les limites à gauche et à droite sont différentes (point anguleux, ex: $|x|$ en $0$).
\end{itemize}
\end{attention}

\courssub{Tangente à une courbe en un point}

\begin{definition}[Tangente à une courbe]
Soit $f$ dérivable en $a$. La \textbf{tangente} à la courbe $\mathcal{C}_f$ au point $A(a; f(a))$ est la droite passant par $A$ de coefficient directeur $f'(a)$.
\end{definition}

\begin{propriete}[Équation réduite de la tangente]
L'équation réduite de la tangente $\mathcal{T}_A$ en $A(a; f(a))$ est :
\[
y = f'(a)(x - a) + f(a)
\]
Soit sous la forme $y = mx + p$ avec $m = f'(a)$ et $p = f(a) - a f'(a)$.
\end{propriete}

\begin{methode}[Déterminer l'équation d'une tangente]
\begin{enumerate}
    \item Calculer $f'(x)$ (fonction dérivée) puis évaluer $f'(a)$.
    \item Calculer $f(a)$ (ordonnée du point de contact).
    \item Écrire l'équation : $y = f'(a)(x - a) + f(a)$ et la réduire sous la forme $y = mx + p$.
\end{enumerate}
\end{methode}

\begin{exemplebox}[Tangente à $f(x) = -2x^2 + 4x - 1$ en $a=1$]
$f'(x) = -4x + 4 \Rightarrow f'(1) = 0$.
$f(1) = -2 + 4 - 1 = 1$.
Tangente : $y = 0 \cdot (x-1) + 1 \Rightarrow \boxed{y = 1}$ (tangente horizontale).
Intersection avec l'axe des ordonnées ($x=0$) : $y=1$, point $B(0;1)$.
\end{exemplebox}

\courssub{Fonction dérivée et dérivées des fonctions usuelles}

\begin{definition}[Fonction dérivée]
Si $f$ est dérivable en tout point d'un intervalle $I$, on définit la \textbf{fonction dérivée} $f'$ sur $I$ par : $\forall x \in I,\ f'(x) = \lim_{h \to 0} \frac{f(x+h)-f(x)}{h}$.
\end{definition}

\begin{aretenir}[Tableau des dérivées usuelles (Première Technique)]
Pour $x$ dans l'ensemble de définition indiqué :
\begin{center}
\begin{tabular}{|c|c|c|c|}\hline
$f(x)$ & $f'(x)$ & Domaine de dérivabilité & Note \\ \hline
$k$ (constante) & $0$ & $\mathbb{R}$ &  \\ \hline
$x^n$ ($n \in \mathbb{N}^*$) & $n x^{n-1}$ & $\mathbb{R}$ & $n=1 \Rightarrow 1$ \\ \hline
$\dfrac{1}{x} = x^{-1}$ & $-\dfrac{1}{x^2}$ & $\mathbb{R}^*$ &  \\ \hline
$\sqrt{x}$ & $\dfrac{1}{2\sqrt{x}}$ & $]0; +\infty[$ & Non dérivable en $0$ \\ \hline
$k f(x)$ & $k f'(x)$ & $\text{Dom } f'$ & Linéarité \\ \hline
$f(x)+g(x)$ & $f'(x)+g'(x)$ & $\text{Dom } f' \cap \text{Dom } g'$ & Somme \\ \hline
$f(x)g(x)$ & $f'(x)g(x)+f(x)g'(x)$ & $\text{Dom } f' \cap \text{Dom } g'$ & Produit \\ \hline
$\dfrac{f(x)}{g(x)}$ & $\dfrac{f'(x)g(x)-f(x)g'(x)}{g(x)^2}$ & $\text{Dom } f' \cap \text{Dom } g' \cap \{g \neq 0\}$ & Quotient \\ \hline
\end{tabular}
\end{center}
\emph{Rappel :} Pour $u(x)=x^n$, la formule $(u^n)' = n u' u^{n-1}$ (dérivée de fonction composée puissance) est au programme de Terminale. En Première, on développe $(u(x))^n$ si $n$ est petit ou on utilise le tableau ci-dessus sur des formes simples.
\end{aretenir}

\begin{exemplebox}[Dérivée de fonctions usuelles]
\begin{enumerate}
    \item $f(x) = 3x^2 - 5x + 2 \Rightarrow f'(x) = 6x - 5$.
    \item $g(x) = \dfrac{x+2}{x-1}$ ($x \neq 1$).
    $g'(x) = \dfrac{1\cdot(x-1) - (x+2)\cdot 1}{(x-1)^2} = \dfrac{-3}{(x-1)^2}$.
    \item $h(x) = \sqrt{x} \ (x>0) \Rightarrow h'(x) = \dfrac{1}{2\sqrt{x}}$. $h'(4) = \dfrac{1}{4}$.
\end{enumerate}
\end{exemplebox}

\courssub{Opérations sur les fonctions dérivées}

\begin{propriete}[Règles de dérivation -- Synthèse]
Soient $u$ et $v$ dérivables sur un intervalle $I$.
\begin{itemize}
    \item \textbf{Somme :} $(u+v)' = u' + v'$
    \item \textbf{Produit :} $(uv)' = u'v + uv'$ \quad (\textbf{pas} $u'v'$ !)
    \item \textbf{Quotient :} $\left(\dfrac{u}{v}\right)' = \dfrac{u'v - uv'}{v^2}$ \quad ($v(x) \neq 0$ sur $I$)
    \item \textbf{Inverse :} $\left(\dfrac{1}{v}\right)' = -\dfrac{v'}{v^2}$ \quad ($v(x) \neq 0$)
\end{itemize}
\end{propriete}

\begin{attention}[Erreurs classiques à éviter]
\begin{itemize}
    \item Confondre $(uv)'$ et $u'v'$.
    \item Oublier le carré au dénominateur pour le quotient : $(u/v)' \neq (u'v - uv')/v$.
    \item Inverser les termes au numérateur du quotient : c'est $u'v - uv'$ (dérivée du numérateur $\times$ dénominateur \textbf{moins} numérateur $\times$ dérivée du dénominateur).
    \item Ne pas préciser l'ensemble de définition de la dérivée (exclusion des zéros du dénominateur, de $0$ pour $\sqrt{x}$, etc.).
\end{itemize}
\end{attention}

\courssub{Applications concrètes de la dérivation}

\courssub{Étude de variations et tableau de variations}

\begin{propriete}[Signe de la dérivée et variations]
Soit $f$ dérivable sur un intervalle $I$.
\begin{itemize}
    \item Si $f'(x) \ge 0$ sur $I$, alors $f$ est croissante sur $I$.
    \item Si $f'(x) \le 0$ sur $I$, alors $f$ est décroissante sur $I$.
    \item Si $f'(x) = 0$ sur $I$, alors $f$ est constante sur $I$.
\end{itemize}
Un changement de signe de $f'$ correspond à un extremum (maximum si $+ \to -$, minimum si $- \to +$).
\end{propriete}

\begin{methode}[Dresser un tableau de variations]
\begin{enumerate}
    \item Calculer $f'(x)$ et la factoriser si possible.
    \item Résoudre $f'(x) = 0$ pour trouver les valeurs critiques.
    \item Étudier le signe de $f'(x)$ sur chaque intervalle (tableau de signes).
    \item Construire le tableau de variations : ligne $x$, ligne $f'(x)$ (signes + flèches $\nearrow/\searrow$), ligne $f(x)$ (valeurs aux bornes et extrema, flèches).
\end{enumerate}
\end{methode}

\begin{exoresolu}[Variations de $f(x) = x^3 - 3x + 1$]
$f'(x) = 3x^2 - 3 = 3(x-1)(x+1)$.
$f'(x)=0 \Leftrightarrow x=-1$ ou $x=1$.
\begin{center}
\begin{tabular}{|c|ccccccc|}\hline
$x$ & $-\infty$ & & $-1$ & & $1$ & & $+\infty$ \\ \hline
$f'(x)$ & & $+$ & $0$ & $-$ & $0$ & $+$ & \\ \hline
$f(x)$ & $-\infty$ & $\nearrow$ & $3$ & $\searrow$ & $-1$ & $\nearrow$ & $+\infty$ \\ \hline
\end{tabular}
\end{center}
$f$ admet un maximum $3$ en $-1$ et un minimum $-1$ en $1$.
\tcblower
\textbf{Solution détaillée} : Voir calculs ci-dessus. Le tableau respecte la règle : 3 lignes, même nombre de cases. Les flèches sont \textbf{uniquement} sur la ligne de $f$.
\end{exoresolu}

\courssub{Problèmes d'optimisation (Maximum / Minimum)}

\begin{methode}[Résoudre un problème d'optimisation]
\begin{enumerate}
    \item Définir la variable $x$ (longueur, quantité, etc.) et son ensemble de définition $D$ (contraintes physiques : $x>0$, aire fixée, etc.).
    \item Exprimer la grandeur à optimiser $P(x)$ (périmètre, coût, bénéfice, volume...) en fonction de $x$ uniquement.
    \item Étudier les variations de $P$ sur $D$ (dérivée, signe, tableau de variations).
    \item Conclure : la valeur optimale est atteinte pour $x_0$ (extremum ou borne de $D$). Donner la valeur optimale $P(x_0)$ et répondre à la question posée (dimensions, coût, etc.).
\end{enumerate}
\end{methode}

\begin{exoresolu}[Clôture contre un mur -- Aire $200\,\text{m}^2$]
Soit $x$ la largeur (perp. au mur), $L$ la longueur (par. au mur).
\begin{enumerate}
    \item $x \cdot L = 200 \Rightarrow L = \frac{200}{x}$. $D = ]0; +\infty[$.
    \item Longueur clôture $P(x) = 2x + L = 2x + \frac{200}{x}$.
    \item $P'(x) = 2 - \frac{200}{x^2} = \frac{2(x^2-100)}{x^2} = \frac{2(x-10)(x+10)}{x^2}$.
    Signe de $P'$ sur $]0;+\infty[$ : $-$ sur $]0;10]$, $+$ sur $[10;+\infty[$.
    \item Minimum pour $x=10\,\text{m}$. $L = 20\,\text{m}$. $P_{\min} = 40\,\text{m}$.
    \item Coût minimal = $40 \times 2\,500 = 100\,000$ FCFA.
\end{enumerate}
\tcblower
\textbf{Solution détaillée} : Voir étapes ci-dessus. Vérification : $P(10)=40$, $P(5)=50$, $P(20)=50$. C'est bien un minimum global.
\end{exoresolu}

\courssub{Applications économiques : Coût marginal, Revenu, Bénéfice}

\begin{definition}[Fonctions économiques usuelles]
Soit $x$ la quantité produite/vendue (souvent en centaines ou milliers d'unités).
\begin{itemize}
    \item \textbf{Coût total :} $C(x)$ (souvent $C(x) = \text{fixes} + \text{variables}$).
    \item \textbf{Coût marginal :} $C'(x)$ = coût approximatif de production de l'unité supplémentaire au voisinage de $x$. Unité : unité de $C$ / unité de $x$.
    \item \textbf{Prix unitaire :} $p(x)$ (fonction décroissante généralement).
    \item \textbf{Revenu total :} $R(x) = x \cdot p(x)$ (attention aux unités : si $x$ en centaines et $p$ en kFCFA/unité, $R$ en kFCFA $\times$ centaine ? Convention courante : $R(x)=x \cdot p(x)$ donne le revenu en kFCFA si $x$ est en centaines et $p$ en kFCFA).
    \item \textbf{Bénéfice :} $B(x) = R(x) - C(x)$.
\end{itemize}
\end{definition}

\begin{exemplebox}[Coût marginal -- Usine de ciment]
$C(q) = q^2 + 10q + 500$ (kFCFA), $q$ en centaines de sacs.
$C'(q) = 2q + 10$ (kFCFA par centaine de sacs).
$C'(20) = 50$ kFCFA/centaine = $500$ FCFA/sac.
Interprétation : Produire la 2001ème centaine de sacs coûte environ $50\,000$ FCFA de plus (coût marginal).
\end{exemplebox}

\courssub{Tangentes : cas particuliers}

\begin{itemize}
    \item \textbf{Tangente parallèle à une droite $y=mx+p$ :} Résoudre $f'(a) = m$, puis trouver $f(a)$ et l'équation.
    \item \textbf{Tangente passant par un point $A(x_A; y_A)$ hors courbe :} L'inconnue est l'abscisse du contact $a$. Résoudre $y_A = f'(a)(x_A - a) + f(a)$.
    \item \textbf{Tangente horizontale :} Résoudre $f'(a) = 0$ (extremum de $f$).
\end{itemize}

\begin{exemplebox}[Tangente parallèle à $y=2x-7$ pour $f(x)=x^2-6x+5$]
$f'(x)=2x-6$. Parallèle $\Rightarrow 2x-6=2 \Rightarrow x=4$.
$f(4)=16-24+5=-3$. Point $M(4;-3)$.
Tangente : $y+3=2(x-4) \Rightarrow \boxed{y=2x-11}$.
\end{exemplebox}

\courssub{Synthèse, compléments et vers la Terminale}

\begin{aretenir}[Check-list pour le Probatoire]
\begin{itemize}
    \item \textbf{Savoir calculer} $f'(a)$ par la définition (limite du taux d'accroissement).
    \item \textbf{Connaître par cœur} le tableau des dérivées usuelles ($x^n, 1/x, \sqrt{x}$) et les règles de dérivation (somme, produit, quotient).
    \item \textbf{Savoir écrire} l'équation de la tangente : $y = f'(a)(x-a) + f(a)$.
    \item \textbf{Savoir dresser} un tableau de variations complet (3 lignes, même nombre de colonnes, flèches sur ligne $f$).
    \item \textbf{Savoir modéliser} un problème concret (optimisation, économie) : variable, fonction, domaine, variations, conclusion \textbf{en phrase}.
    \item \textbf{Justifier} chaque étape : "car $f'$ est de signe constant sur $I$", "car $f'(x)=0$ et $f'$ change de signe", etc.
\end{itemize}
\end{aretenir}

\begin{savaistu}[Vers la Terminale Technique]
En Terminale, vous étendrez la dérivation aux fonctions :
\begin{itemize}
    \item Exponentielle : $(e^u)' = u' e^u$, $(a^x)' = a^x \ln a$.
    \item Logarithme népérien : $(\ln u)' = \frac{u'}{u}$.
    \item Trigonométriques : $(\sin u)' = u'\cos u$, $(\cos u)' = -u'\sin u$.
    \item Réciproques : $(\arcsin x)' = \frac{1}{\sqrt{1-x^2}}$, etc.
\end{itemize}
Vous verrez aussi les \textbf{primitives} (opération inverse) et les \textbf{intégrales} (calcul d'aires, valeurs moyennes). La dérivée reste l'outil principal pour l'étude de fonctions et l'optimisation.
\end{savaistu}

\courssub{Je vérifie mes connaissances}

\begin{exercice}[QCM : Dérivées des fonctions usuelles]
Parmi les quatre propositions suivantes, une seule est exacte. La recopier sur la copie en justifiant brièvement.
\begin{enumerate}
    \item Soit $f(x) = 3x^2 - 5x + 2$. La dérivée $f'(x)$ est égale à :
    \begin{itemize}
        \item[A.] $6x - 5$
        \item[B.] $6x + 5$
        \item[C.] $3x - 5$
        \item[D.] $6x^2 - 5$
    \end{itemize}
    \item Soit $g(x) = \dfrac{1}{x}$ sur $\mathbb{R}^*$. La dérivée $g'(x)$ est égale à :
    \begin{itemize}
        \item[A.] $\dfrac{1}{x^2}$
        \item[B.] $-\dfrac{1}{x^2}$
        \item[C.] $\ln|x|$
        \item[D.] $-\dfrac{1}{x}$
    \end{itemize}
    \item Soit $h(x) = \sqrt{x}$ sur $[0; +\infty[$. Le nombre dérivé de $h$ en $4$ vaut :
    \begin{itemize}
        \item[A.] $\dfrac{1}{2}$
        \item[B.] $\dfrac{1}{4}$
        \item[C.] $2$
        \item[D.] $4$
    \end{itemize}
    \item La tangente à la courbe $y = x^3$ au point d'abscisse $1$ a pour équation réduite :
    \begin{itemize}
        \item[A.] $y = 3x - 2$
        \item[B.] $y = 3x + 2$
        \item[C.] $y = x - 2$
        \item[D.] $y = 3x - 1$
    \end{itemize}
\end{enumerate}
\end{exercice}

\begin{exercice}[Vrai ou Faux ? Justifier]
Dire si les affirmations suivantes sont vraies ou fausses. Justifier chaque réponse par un calcul, un contre-exemple ou une propriété du cours.
\begin{enumerate}
    \item Si $f'(a) = 0$ alors $f$ admet un maximum en $a$.
    \item Si $u$ et $v$ sont dérivables sur un intervalle $I$, alors $(uv)' = u'v'$ sur $I$.
    \item La fonction $x \mapsto \sqrt{x}$ est dérivable en $0$.
    \item Si la tangente à la courbe $\mathcal{C}_f$ en $a$ est horizontale, alors $f'(a) = 0$.
    \item La fonction $f(x) = |x|$ est dérivable sur $\mathbb{R}$.
\end{enumerate}
\end{exercice}

\begin{exercice}[Nombre dérivé par la définition]
Soit la fonction $f$ définie sur $\mathbb{R}$ par $f(x) = x^2 + 3x$.
\begin{enumerate}
    \item Exprimer le taux d'accroissement $\dfrac{f(2+h) - f(2)}{h}$ pour $h \neq 0$ et le simplifier.
    \item En faisant tendre $h$ vers $0$, déterminer le nombre dérivé $f'(2)$.
    \item Vérifier ce résultat en utilisant les formules usuelles de dérivation sur $f$.
\end{enumerate}
\end{exercice}

\begin{exercice}[Équation de tangente : application directe]
Soit la fonction $f$ définie sur $\mathbb{R}$ par $f(x) = -2x^2 + 4x - 1$. On note $\mathcal{C}_f$ sa courbe représentative.
\begin{enumerate}
    \item Calculer $f'(x)$.
    \item Déterminer l'équation réduite de la tangente $\mathcal{T}$ à $\mathcal{C}_f$ au point d'abscisse $a = 1$.
    \item Calculer les coordonnées du point d'intersection de $\mathcal{T}$ avec l'axe des ordonnées.
\end{enumerate}
\end{exercice}

\courssub{Je m'entraîne}

\begin{exercice}[Dérivée et tangente pour une fonction rationnelle]
Soit la fonction $f$ définie sur $\mathbb{R} \setminus \{1\}$ par $f(x) = \dfrac{x+2}{x-1}$.
\begin{enumerate}
    \item Calculer $f'(x)$ pour tout $x \neq 1$.
    \item Déterminer le coefficient directeur de la tangente $\mathcal{T}$ à la courbe $\mathcal{C}_f$ au point d'abscisse $a = 2$.
    \item Donner l'équation réduite de $\mathcal{T}$.
    \item Résoudre l'équation $f'(x) = 0$ et interpréter graphiquement le résultat.
\end{enumerate}
\end{exercice}

\begin{exercice}[Lecture graphique de la dérivée]
On donne ci-dessous la courbe $\mathcal{C}_f$ d'une fonction $f$ dérivable sur $[-3; 3]$. Trois tangentes $\mathcal{T}_{-1}$, $\mathcal{T}_0$, $\mathcal{T}_2$ ont été tracées aux points d'abscisses $-1$, $0$ et $2$.
\begin{popfigure}
\begin{tikzpicture}[scale=0.7]
\begin{axis}[
    axis lines=middle,
    xlabel=$x$, ylabel=$y$,
    xmin=-3.5, xmax=3.5, ymin=-3.5, ymax=3.5,
    xtick={-3,-2,-1,1,2,3}, ytick={-3,-2,-1,1,2,3},
    grid=major, grid style={popline},
    width=\linewidth, height=8cm,
    clip=false
]
% Function: f(x) = 0.2*x^3 - 0.6*x
\addplot[popcurve, domain=-3:3, samples=200, thick] {0.2*x^3 - 0.6*x};
% Tangent at -1: f(-1)=0.4, f'(-1)=0 -> y=0.4
\addplot[poporange, dashed, domain=-2.5:0.5, samples=2] {0.4};
% Tangent at 0: f(0)=0, f'(0)=-0.6 -> y=-0.6x
\addplot[popgreen, dashed, domain=-1.5:1.5, samples=2] {-0.6*x};
% Tangent at 2: f(2)=0.4, f'(2)=1.8 -> y = 1.8x - 3.2
\addplot[poppurple, dashed, domain=1:3, samples=2] {1.8*x - 3.2};
% Points
\addplot[only marks, mark=*, popdark] coordinates {(-1,0.4) (0,0) (2,0.4)};
\node[popdark, anchor=south west] at (axis cs:-1,0.4) {$A(-1;0,4)$};
\node[popdark, anchor=north west] at (axis cs:0,0) {$O$};
\node[popdark, anchor=south west] at (axis cs:2,0.4) {$B(2;0,4)$};
\node[poporange, anchor=south east] at (axis cs:-2.5,0.4) {$\mathcal{T}_{-1}$};
\node[popgreen, anchor=north east] at (axis cs:1.5,-0.9) {$\mathcal{T}_0$};
\node[poppurple, anchor=south west] at (axis cs:2.5,1.3) {$\mathcal{T}_2$};
\end{axis}
\end{tikzpicture}
\legende{Courbe $\mathcal{C}_f$ et tangentes en $-1$, $0$ et $2$}
\end{popfigure}
\begin{enumerate}
    \item En utilisant le graphique, donner une valeur approchée de $f'(-1)$, $f'(0)$ et $f'(2)$.
    \item Résoudre graphiquement l'équation $f'(x) = 0$ sur $[-3; 3]$.
    \item En déduire les intervalles de variations de $f$ sur $[-3; 3]$.
\end{enumerate}
\end{exercice}

\begin{exercice}[Tangente parallèle à une droite donnée]
Soit $f(x) = x^2 - 6x + 5$ définie sur $\mathbb{R}$.
\begin{enumerate}
    \item Calculer $f'(x)$.
    \item Déterminer les points de la courbe $\mathcal{C}_f$ où la tangente est parallèle à la droite $\mathcal{D}$ d'équation $y = 2x - 7$.
    \item Écrire les équations réduites de ces tangentes.
\end{enumerate}
\end{exercice}

\begin{exercice}[Application économique : Coût marginal]
Une usine produit des sacs de ciment. Le coût total de production (en milliers de FCFA) pour $q$ centaines de sacs est modélisé par la fonction $C(q) = q^2 + 10q + 500$, pour $q \ge 0$.
\begin{enumerate}
    \item Calculer la fonction coût marginal $C'(q)$. Quelle est son unité ?
    \item Calculer $C'(20)$ et interpréter ce nombre dans le contexte de l'exercice.
    \item Le prix de vente unitaire (au sac) est fixé à $150$ FCFA. Exprimer le revenu $R(q)$ et le bénéfice $B(q)$ en milliers de FCFA.
    \item Déterminer la quantité $q$ qui maximise le bénéfice. Calculer ce bénéfice maximal.
\end{enumerate}
\end{exercice}

\courssub{Je me prépare au Probatoire}

\begin{exercice}[Étude complète et tangente -- Type Probatoire]
Soit la fonction $f$ définie sur $\mathbb{R}$ par $f(x) = x^3 - 3x + 1$. On note $\mathcal{C}_f$ sa courbe représentative dans un repère orthonormé.
\begin{enumerate}
    \item Étudier les variations de $f$ sur $\mathbb{R}$. Dresser son tableau de variations.
    \item Déterminer l'équation réduite de la tangente $\mathcal{T}$ à $\mathcal{C}_f$ au point d'abscisse $1$.
    \item Montrer que $\mathcal{T}$ coupe $\mathcal{C}_f$ en un second point $B$. Déterminer les coordonnées de $B$.
    \item Tracer $\mathcal{C}_f$ et $\mathcal{T}$ dans le même repère (on prendra $1$ cm pour $1$ unité sur les deux axes).
    \item Résoudre graphiquement l'équation $f(x) = 0$ et encadrer les solutions à $10^{-1}$ près.
\end{enumerate}
\end{exercice}

\begin{exercice}[Problème d'optimisation : Clôture contre un mur -- Type Probatoire]
Un artisan souhaite clôturer un terrain rectangulaire d'aire $200\,\text{m}^2$ accolé à un mur. Seuls trois côtés doivent être clôturés (les deux largeurs et une longueur). Soit $x$ (en mètres) la longueur du côté perpendiculaire au mur.
\begin{enumerate}
    \item Exprimer la longueur $L$ du côté parallèle au mur en fonction de $x$.
    \item Exprimer la longueur totale de clôture $P(x)$ nécessaire en fonction de $x$. Préciser l'ensemble de définition de $P$.
    \item Étudier les variations de $P$ sur son ensemble de définition.
    \item Quelles dimensions le terrain doit-il avoir pour que la longueur de clôture soit minimale ? Calculer cette longueur minimale.
    \item Si le grillage coûte $2\,500$ FCFA le mètre linéaire, quel est le coût minimal de la clôture ?
\end{enumerate}
\end{exercice}

\begin{exercice}[Modélisation : Recette et Bénéfice -- Type Probatoire]
Une petite entreprise de menuiserie fabrique et vend des chaises. L'étude du marché montre que le prix de vente unitaire $p$ (en milliers de FCFA) en fonction du nombre $x$ de chaises produites et vendues par mois (en centaines) est donné par $p(x) = 10 - 0,5x$ pour $x \in [0; 20]$. Le coût total de production mensuel $C(x)$ (en milliers de FCFA) est modélisé par $C(x) = x^2 + 4x + 10$.
\begin{enumerate}
    \item Exprimer la recette totale $R(x)$ et le bénéfice $B(x)$ en fonction de $x$.
    \item Calculer $B'(x)$ et dresser le tableau de variations de $B$ sur $[0; 20]$.
    \item Déterminer le nombre de chaises à produire pour maximiser le bénéfice mensuel. Calculer ce bénéfice maximal.
    \item L'entreprise doit-elle accepter une commande de $1\,500$ chaises ($x=15$) si cela l'oblige à arrêter sa production habituelle ? Justifier par le calcul.
\end{enumerate}
\end{exercice}

\newpage

\coursec{Corrigés des exercices}

\coursec{Approche intuitive du nombre dérivé}

\begin{definition}[Taux d'accroissement]
Soit $f$ une fonction définie sur un intervalle $I$ et $a \in I$. Pour tout $h \neq 0$ tel que $a+h \in I$, le quotient
\[
\frac{f(a+h)-f(a)}{h}
\]
est appelé \textbf{taux d'accroissement} de $f$ sur l'intervalle $[a, a+h]$ (ou $[a+h, a]$ si $h<0$). Il représente la pente de la sécante passant par les points $A(a; f(a))$ et $B(a+h; f(a+h))$ de la courbe $\mathcal{C}_f$.
\end{definition}

\begin{definition}[Nombre dérivé en un point]
Si le taux d'accroissement admet une limite finie lorsque $h \to 0$, on dit que $f$ est \textbf{dérivable en $a$}. Cette limite est appelée \textbf{nombre dérivé} de $f$ en $a$ et se note $f'(a)$ :
\[
f'(a) = \lim_{h \to 0} \frac{f(a+h)-f(a)}{h}.
\]
Géométriquement, $f'(a)$ est le \textbf{coefficient directeur de la tangente} à $\mathcal{C}_f$ au point d'abscisse $a$.
\end{definition}

\begin{attention}[Dérivabilité et continuité]
Si $f$ est dérivable en $a$, alors $f$ est continue en $a$. La réciproque est fausse : une fonction peut être continue sans être dérivable (ex. : $f(x)=|x|$ en $0$, $f(x)=\sqrt{x}$ en $0$).
\end{attention}

\begin{exemplebox}[Calcul du nombre dérivé par la définition]
Soit $f(x)=x^2+3x$. Calculer $f'(2)$ par la définition.
\begin{enumerate}
    \item $f(2)=10$ et $f(2+h) = (2+h)^2+3(2+h) = h^2+7h+10$.
    \item Taux d'accroissement : $\dfrac{f(2+h)-f(2)}{h} = \dfrac{h^2+7h}{h} = h+7 \quad (h \neq 0)$.
    \item Limite : $f'(2) = \lim_{h \to 0} (h+7) = 7$.
\end{enumerate}
Vérification par les formules : $f'(x)=2x+3 \Rightarrow f'(2)=7$. \checkmark
\end{exemplebox}

\begin{savaistu}[Histoire : La tangente et le calcul infinitésimal]
La notion de tangente a occupé les mathématiciens depuis l'Antiquité (Archimède). Au XVII\textsuperscript{e} siècle, Fermat, Descartes, puis Newton et Leibniz formalisent la « méthode des tangentes » qui devient le calcul différentiel. Le symbole $f'(x)$ est dû à Lagrange (1797), tandis que $\frac{dy}{dx}$ vient de Leibniz.
\end{savaistu}

\coursec{Tangente à une courbe en un point}

\begin{propriete}[Équation de la tangente]
Soit $f$ dérivable en $a$. La tangente $\mathcal{T}_a$ à la courbe $\mathcal{C}_f$ au point $A(a; f(a))$ a pour équation :
\[
\boxed{y = f'(a)(x-a) + f(a)}.
\]
Si $f'(a)=0$, la tangente est \textbf{horizontale} (équation $y = f(a)$). Si la limite du taux d'accroissement est infinie, la tangente est \textbf{verticale} (équation $x=a$).
\end{propriete}

\begin{methode}[Déterminer l'équation d'une tangente]
\begin{enumerate}
    \item Vérifier que $f$ est dérivable en $a$ (appartenance au domaine, pas de point anguleux, pas de tangente verticale).
    \item Calculer $f'(x)$ (par les formules usuelles ou règles de dérivation).
    \item Calculer $f'(a)$ (coefficient directeur) et $f(a)$ (ordonnée du point de contact).
    \item Écrire l'équation : $y = f'(a)(x-a) + f(a)$.
\end{enumerate}
\end{methode}

\begin{exemplebox}[Tangente horizontale au sommet d'une parabole]
$f(x)=-2x^2+4x-1$. $f'(x)=-4x+4$. En $a=1$ : $f'(1)=0$, $f(1)=1$. Tangente : $y=1$ (horizontale). Le point $A(1;1)$ est le sommet (maximum).
\end{exemplebox}

\begin{exemplebox}[Tangente parallèle à une droite donnée]
$f(x)=x^2-6x+5$, droite $\mathcal{D}: y=2x-7$. $f'(x)=2x-6$.
Parallélisme $\iff$ même coefficient directeur : $f'(a)=2 \iff 2a-6=2 \iff a=4$.
$f(4)=-3$. Point $M(4;-3)$. Tangente : $y=2(x-4)-3 \Rightarrow y=2x-11$.
\end{exemplebox}

\coursec{Fonction dérivée et dérivées des fonctions usuelles}

\begin{definition}[Fonction dérivée]
Si $f$ est dérivable en tout point d'un intervalle $I$, on définit la \textbf{fonction dérivée} $f'$ sur $I$ par $x \mapsto f'(x)$.
\end{definition}

\begin{propriete}[Tableau des dérivées usuelles (à connaître par cœur)]
\begin{center}
\begin{tabularx}{\linewidth}{|l|X|l|}\hline
\textbf{Fonction $f(x)$} & \textbf{Dérivée $f'(x)$} & \textbf{Domaine de dérivabilité} \\ \hline
$k$ (constante) & $0$ & $\mathbb{R}$ \\ \hline
$x^n$ ($n \in \mathbb{N}^*$) & $n x^{n-1}$ & $\mathbb{R}$ \\ \hline
$\dfrac{1}{x}$ & $-\dfrac{1}{x^2}$ & $\mathbb{R}^*$ \\ \hline
$\sqrt{x}$ & $\dfrac{1}{2\sqrt{x}}$ & $]0;+\infty[$ \\ \hline
\end{tabularx}
\end{center}
\emph{Remarque :} La dérivée de $\ln|x|$ est $\frac{1}{x}$ mais le logarithme népérien n'est \textbf{pas au programme de Première technique} (c'est une primitive, vue en Terminale).
\end{propriete}

\begin{aretenir}[Opérations sur les fonctions dérivées]
Soit $f, g$ dérivables sur $I$, $\lambda \in \mathbb{R}$.
\begin{itemize}
    \item \textbf{Somme :} $(\lambda f + g)' = \lambda f' + g'$ \quad (Linéarité)
    \item \textbf{Produit :} $(fg)' = f'g + fg'$
    \item \textbf{Quotient :} $\left(\dfrac{f}{g}\right)' = \dfrac{f'g - fg'}{g^2}$ \quad ($g(x) \neq 0$)
\end{itemize}
\end{aretenir}

\begin{attention}[Pièges classiques]
\begin{itemize}
    \item $(fg)' \neq f'g'$ et $(f/g)' \neq f'/g'$.
    \item Pour le quotient : l'ordre au numérateur est crucial : $f'g - fg'$ (dérivée du haut $\times$ bas $-$ haut $\times$ dérivée du bas).
    \item Domaine de la dérivée = domaine de la fonction \emph{moins} les points où $f$ n'est pas dérivable (ex. $x=0$ pour $\sqrt{x}$ ou $1/x$).
\end{itemize}
\end{attention}

\begin{exemplebox}[Dérivée d'une fonction rationnelle]
$f(x)=\dfrac{x+2}{x-1}$ sur $\mathbb{R}\setminus\{1\}$.
$u=x+2 \Rightarrow u'=1$ ; $v=x-1 \Rightarrow v'=1$.
$f'(x) = \dfrac{1\cdot(x-1) - (x+2)\cdot 1}{(x-1)^2} = \dfrac{-3}{(x-1)^2}$.
$f'(x) < 0$ sur chaque intervalle de définition $\Rightarrow f$ strictement décroissante sur $]-\infty;1[$ et $]1;+\infty[$. Pas de tangente horizontale.
\end{exemplebox}

\coursec{Opérations sur les fonctions dérivées}

\begin{methode}[Dériver une fonction composée d'opérations]
\begin{enumerate}
    \item Identifier la structure : somme, produit, quotient de fonctions usuelles.
    \item Appliquer la règle correspondante (linéarité, produit, quotient).
    \item Simplifier l'expression obtenue (factoriser si possible pour l'étude de signe).
    \item Préciser le domaine de la dérivée.
\end{enumerate}
\end{methode}

\begin{exoresolu}[Entraînement : Dérivées et opérations]
\begin{enumerate}
    \item $f(x)=3x^2-5x+2 \Rightarrow f'(x)=6x-5$.
    \item $g(x)=\dfrac{1}{x} \Rightarrow g'(x)=-\dfrac{1}{x^2}$.
    \item $h(x)=\sqrt{x} \Rightarrow h'(x)=\dfrac{1}{2\sqrt{x}}$, $h'(4)=\dfrac{1}{4}$.
    \item $f(x)=x^3 \Rightarrow f'(x)=3x^2$, tangente en $1$ : $y=3x-2$.
\end{enumerate}
\end{exoresolu}

\coursec{Applications concrètes de la dérivation}

\begin{savaistu}[Vocabulaire économique]
\begin{itemize}
    \item $C(q)$ : Coût total de production (en kFCFA). $C'(q)$ : \textbf{Coût marginal} (coût d'une unité supplémentaire).
    \item $R(q) = q \times p(q)$ : Revenu total (Recette). $R'(q)$ : Revenu marginal.
    \item $B(q) = R(q) - C(q)$ : Bénéfice (Profit). $B'(q)$ : Bénéfice marginal.
    \item Optimisation : On cherche $q$ tel que $B'(q)=0$ et $B'$ passe de $+$ à $-$ (maximum).
\end{itemize}
\end{savaistu}

\begin{exoresolu}[Exercice 8 — Coût marginal et optimisation du bénéfice]
$C(q)=q^2+10q+500$ (kFCFA), $q$ en centaines de sacs.
\begin{enumerate}
    \item $C'(q)=2q+10$. Unité : \textbf{kFCFA par centaine de sacs}.
    \item $C'(20)=50$ kFCFA/centaine de sacs = $50\,000$ FCFA pour 100 sacs $\approx 500$ FCFA/sac.
    \item Prix de vente $150$ FCFA/sac $\Rightarrow$ Revenu $R(q) = 150 \times 100q = 15\,000q$ FCFA $= 15q$ kFCFA.
    $B(q)=R(q)-C(q) = -q^2+5q-500$.
    \item $B'(q)=-2q+5$. $B'(q)=0 \iff q=2,5$. $B'$ passe de $+$ à $-$ $\Rightarrow$ Maximum en $q=2,5$.
    $B(2,5)=-493,75$ kFCFA. \textbf{Perte minimale} de $493\,750$ FCFA pour $250$ sacs.
    \textbf{Conclusion :} L'entreprise est structurellement déficitaire à ce prix ; elle doit augmenter son prix de vente ou baisser ses coûts fixes.
\end{enumerate}
\end{exoresolu}

\begin{exoresolu}[Exercice 10 — Problème d'optimisation : Clôture contre un mur]
\begin{enumerate}
    \item Aire $= x \times L = 200 \Rightarrow L = \dfrac{200}{x}$.
    \item Trois côtés clôturés : $P(x) = 2x + L = 2x + \dfrac{200}{x}$. Domaine $D_P = ]0;+\infty[$.
    \item $P'(x) = 2 - \dfrac{200}{x^2} = \dfrac{2(x^2-100)}{x^2} = \dfrac{2(x-10)(x+10)}{x^2}$.
    Sur $]0;+\infty[$, $x^2>0$ et $x+10>0$, signe de $P'$ = signe de $x-10$.
    \begin{center}
    \begin{tabular}{|c|ccccc|}\hline
    $x$ & $0$ & & $10$ & & $+\infty$ \\ \hline
    $P'(x)$ & & $-$ & $0$ & $+$ & \\ \hline
    $P$ & $+\infty$ & $\searrow$ & $40$ & $\nearrow$ & $+\infty$ \\ \hline
    \end{tabular}
    \end{center}
    \item Minimum en $x=10$ m. $L = 20$ m. Longueur minimale $P(10)=40$ m.
    \item Coût minimal : $40 \times 2\,500 = \boxed{100\,000\ \text{FCFA}}$.
\end{enumerate}
\end{exoresolu}

\begin{exoresolu}[Exercice 11 — Modélisation : Recette et Bénéfice (Type Probatoire)]
$p(x)=10-0,5x$ (kFCFA), $C(x)=x^2+4x+10$ (kFCFA), $x\in[0;20]$ (centaines de chaises).
\begin{enumerate}
    \item $R(x)=x(10-0,5x)=-0,5x^2+10x$.
    $B(x)=R(x)-C(x)=-1,5x^2+6x-10$.
    \item $B'(x)=-3x+6=-3(x-2)$. $B'(x)=0 \iff x=2$.
    Signe : $B'(x)>0$ sur $[0;2]$, $B'(x)<0$ sur $[2;20]$.
    \begin{center}
    \begin{tabular}{|c|ccccc|}\hline
    $x$ & $0$ & & $2$ & & $20$ \\ \hline
    $B'(x)$ & & $+$ & $0$ & $-$ & \\ \hline
    $B$ & $-10$ & $\nearrow$ & $-4$ & $\searrow$ & $-490$ \\ \hline
    \end{tabular}
    \end{center}
    \item Maximum pour $x=2$ soit $\boxed{200\ \text{chaises/mois}}$. Bénéfice max $B(2)=-4$ kFCFA ($-4\,000$ FCFA).
    \item Pour $x=15$ ($1\,500$ chaises) : $B(15)=-257,5$ kFCFA. Perte de $257\,500$ FCFA $\ll -4\,000$ FCFA.
    \textbf{Refuser la commande} ou négocier un prix plus élevé.
\end{enumerate}
\end{exoresolu}

\coursec{Synthèse, compléments et vers la Terminale}

\begin{aretenir}[Check-list pour le Probatoire]
\begin{itemize}
    \item Calculer un nombre dérivé par la définition (limite du taux d'accroissement).
    \item Déterminer $f'(x)$ à l'aide du tableau des usuelles et des règles de dérivation (somme, produit, quotient).
    \item Écrire l'équation de la tangente en un point : $y = f'(a)(x-a) + f(a)$.
    \item Résoudre $f'(x)=0$ pour trouver les extrema (tangentes horizontales).
    \item Construire un tableau de variations complet (valeurs aux bornes, extrema, flèches).
    \item Lire graphiquement $f'(a)$ (pente de la tangente) et le signe de $f'$ (croissance/décroissance).
    \item Modéliser un problème concret (économie, géométrie) : exprimer la fonction, la dériver, étudier ses variations, conclure \textbf{en phrase} avec les unités.
\end{itemize}
\end{aretenir}

\begin{exoresolu}[Exercice 2 — Vrai ou Faux ? Justifier]
\begin{enumerate}
    \item \textbf{Faux.} $f'(a)=0$ n'implique pas un extremum. Contre-exemple : $f(x)=x^3$ en $0$. $f'(x)=3x^2 \geq 0$, $f$ strictement croissante, pas d'extremum en $0$ (point d'inflexion à tangente horizontale).
    \item \textbf{Faux.} $(uv)' = u'v + uv' \neq u'v'$. Contre-exemple : $u=v=x \Rightarrow (x^2)'=2x \neq 1$.
    \item \textbf{Faux.} $f(x)=\sqrt{x}$ en $0$. Taux d'accroissement $\frac{\sqrt{h}}{h}=\frac{1}{\sqrt{h}} \xrightarrow[h\to 0^+]{} +\infty$. Tangente verticale, pas de nombre dérivé fini.
    \item \textbf{Vrai.} Tangente horizontale $\iff$ coefficient directeur nul $\iff f'(a)=0$.
    \item \textbf{Faux.} $f(x)=|x|$ non dérivable en $0$ (point anguleux). Taux d'accroissement $1$ à droite, $-1$ à gauche. Dérivable sur $\mathbb{R}^*$ seulement.
\end{enumerate}
\end{exoresolu}

\begin{popfigure}
\centering
\begin{tikzpicture}[scale=0.8, domain=-3:3]
    \draw[->] (-3.5,0) -- (3.5,0) node[right] {$x$};
    \draw[->] (0,-1.5) -- (0,1.5) node[above] {$y$};
    % Function f(x) = 0.2*x^3 - 0.6*x  (scaled for visibility: f'(0)=-0.6, f'(2)=1.8)
    % Actually use f(x) = 1/5 x^3 - 3/5 x
    \draw[popcurve, thick, domain=-3:3, samples=200] plot (\x, {0.2*\x*\x*\x - 0.6*\x});
    % Tangents
    % at x=-1: slope 0, y=0.4
    \draw[poporange, dashed] (-3, 0.4) -- (3, 0.4) node[right] {$\mathcal{T}_{-1}$};
    % at x=0: slope -0.6, y=0
    \draw[popgreen, dashed] (-3, 1.8) -- (3, -1.8) node[right] {$\mathcal{T}_0$};
    % at x=2: slope 1.8, y=0.4
    \draw[poppurple, dashed] (-1, -2.0) -- (3, 3.4) node[right] {$\mathcal{T}_2$};
    % Points
    \fill[popblue] (-1, 0.4) circle (2pt) node[above left] {$A(-1;0,4)$};
    \fill[popblue] (0, 0) circle (2pt) node[below right] {$O$};
    \fill[popblue] (2, 0.4) circle (2pt) node[above right] {$B(2;0,4)$};
    \fill[popblue] (1, -0.4) circle (2pt) node[below right] {$C(1;-0,4)$};
\end{tikzpicture}
\legende{Courbe $y=\frac{1}{5}x^3-\frac{3}{5}x$ et tangentes en $-1$, $0$, $2$ pour l'exercice 6.}
\end{popfigure}

\begin{exoresolu}[Exercice 6 — Lecture graphique de la dérivée]
D'après la figure ci-dessus (courbe $f(x)=\frac{1}{5}x^3-\frac{3}{5}x$) :
\begin{enumerate}
    \item $\mathcal{T}_{-1}$ horizontale $\Rightarrow f'(-1)=0$. $\mathcal{T}_0$ pente $-0,6 \Rightarrow f'(0)\approx -0,6$. $\mathcal{T}_2$ pente $1,8 \Rightarrow f'(2)\approx 1,8$.
    \item $f'(x)=0$ pour tangentes horizontales : $x=-1$ et $x=1$. $\mathcal{S}=\{-1; 1\}$.
    \item $f$ croissante sur $[-3;-1]$ ($f'>0$), décroissante sur $[-1;1]$ ($f'<0$), croissante sur $[1;3]$ ($f'>0$). Max local en $-1$, min local en $1$.
\end{enumerate}
\end{exoresolu}

\begin{popfigure}
\centering
\begin{tikzpicture}[scale=0.6, domain=-2.5:2.5]
    \draw[->] (-2.5,0) -- (2.5,0) node[right] {$x$};
    \draw[->] (0,-2) -- (0,3.5) node[above] {$y$};
    % f(x) = x^3 - 3x + 1
    \draw[popcurve, thick, domain=-2.2:2.2, samples=200] plot (\x, {\x*\x*\x - 3*\x + 1});
    % Tangent y=-1
    \draw[poporange, dashed] (-2.5, -1) -- (2.5, -1) node[right] {$\mathcal{T}: y=-1$};
    % Points
    \fill[popblue] (-1, 3) circle (2pt) node[above left] {$(-1;3)$ Max};
    \fill[popblue] (1, -1) circle (2pt) node[below right] {$(1;-1)$ Min};
    \fill[poporange] (-2, -1) circle (2pt) node[left] {$B(-2;-1)$};
    \fill[popblue] (0, 1) circle (2pt) node[above] {$(0;1)$};
    \fill[popblue] (2, 3) circle (2pt) node[above right] {$(2;3)$};
    % x-axis intersections approx
    \draw[popmuted] (-1.88, -0.2) -- (-1.88, 0.2) node[below] {$x_1$};
    \draw[popmuted] (0.35, -0.2) -- (0.35, 0.2) node[below] {$x_2$};
    \draw[popmuted] (1.53, -0.2) -- (1.53, 0.2) node[below] {$x_3$};
\end{tikzpicture}
\legende{Courbe $\mathcal{C}_f$ de $f(x)=x^3-3x+1$, tangente au minimum et points remarquables (Exercice 9).}
\end{popfigure}

\begin{exoresolu}[Exercice 9 — Étude complète et tangente (Type Probatoire)]
$f(x)=x^3-3x+1$ sur $\mathbb{R}$.
\begin{enumerate}
    \item $f'(x)=3x^2-3=3(x-1)(x+1)$. $f'(x)=0 \iff x=-1$ ou $x=1$.
    Signe de $f'$ (trinôme $3>0$) : $+$ sur $]-\infty;-1[$, $-$ sur $]-1;1[$, $+$ sur $]1;+\infty[$.
    $f(-1)=3$, $f(1)=-1$.
    \begin{center}
    \begin{tabular}{|c|ccccccc|}\hline
    $x$ & $-\infty$ & & $-1$ & & $1$ & & $+\infty$ \\ \hline
    $f'(x)$ & & $+$ & $0$ & $-$ & $0$ & $+$ & \\ \hline
    $f$ & $-\infty$ & $\nearrow$ & $3$ & $\searrow$ & $-1$ & $\nearrow$ & $+\infty$ \\ \hline
    \end{tabular}
    \end{center}
    Maximum local $3$ en $-1$ ; minimum local $-1$ en $1$.
    \item $f'(1)=0$, $f(1)=-1$. Tangente : $\boxed{y=-1}$ (horizontale au minimum).
    \item Intersection $\mathcal{T} \cap \mathcal{C}_f$ : $f(x)=-1 \iff x^3-3x+2=0$.
    $x=1$ est racine (contact, racine double). Factorisation : $x^3-3x+2 = (x-1)^2(x+2)$.
    Seconde intersection : $x+2=0 \iff x=-2$, $y=-1$. Point $\boxed{B(-2;-1)}$.
    \item Tracé : voir figure ci-dessus. Points placés : $(-1;3)$, $(1;-1)$, $(-2;-1)$, $(0;1)$, $(2;3)$.
    \item Zéros de $f$ (encadrements au $10^{-1}$) :
    $f(-2)=-1$, $f(-1,8)\approx 0,57 \Rightarrow -1,9 < x_1 < -1,8$.
    $f(0)=1$, $f(0,4)\approx -0,14 \Rightarrow 0,3 < x_2 < 0,4$.
    $f(1,5)=-0,125$, $f(1,6)\approx 0,30 \Rightarrow 1,5 < x_3 < 1,6$.
\end{enumerate}
\textbf{Barème indicatif (10 pts) :} dérivée et factorisation $1,5$ ; signe et tableau $2$ ; tangente $1$ ; intersection et factorisation $(x-1)^2$ $2$ ; point $B$ $0,5$ ; tracé $1,5$ ; encadrements calculés $1,5$.
\end{exoresolu}

\begin{exercice}[Exercices d'entraînement supplémentaires]
\begin{enumerate}
    \item Soit $f(x)=2x^3-3x^2-12x+5$. Étudier les variations de $f$ et tracer sa courbe. Déterminer l'équation de la tangente au point d'abscisse $2$.
    \item Une entreprise a pour fonction coût $C(x)=0,5x^2-20x+500$ (en kFCFA) pour $x$ centaines d'articles. Le prix de vente unitaire est $p(x)=100-0,2x$ (en FCFA). Déterminer la production maximisant le bénéfice et calculer ce bénéfice maximal.
    \item Soit $f(x)=\dfrac{2x-1}{x+3}$ sur $]-3;+\infty[$. Étudier les variations, déterminer l'équation de la tangente en $x=1$ et tracer la courbe.
\end{enumerate}
\end{exercice}

\begin{corrige}[Exercice d'entraînement 1]
$f'(x)=6x^2-6x-12=6(x^2-x-2)=6(x-2)(x+1)$. Racines $-1, 2$.
Tableau : $f'$ $+$ sur $]-\infty;-1[$, $-$ sur $]-1;2[$, $+$ sur $]2;+\infty[$.
$f(-1)=12$ (max), $f(2)=-15$ (min).
Tangente en $2$ : $f'(2)=0$, $f(2)=-15 \Rightarrow y=-15$.
\end{corrige}

\begin{corrige}[Exercice d'entraînement 2]
$R(x) = x \times \frac{p(x)}{1000}$ (conversion FCFA $\to$ kFCFA). $p(x)=100-0,2x$ FCFA $\Rightarrow 0,1-0,0002x$ kFCFA.
$R(x) = x(0,1-0,0002x) = 0,1x - 0,0002x^2$.
$B(x) = R-C = -0,5002x^2 + 20,1x - 500$.
$B'(x) = -1,0004x + 20,1 = 0 \Rightarrow x \approx 20,09$ centaines $\approx 2009$ articles.
$B_{\max} \approx -5002 + 403,8 - 500 \dots$ (Calcul exact laissé à l'élève).
\end{corrige}

\begin{corrige}[Exercice d'entraînement 3]
$f(x)=\frac{2x-1}{x+3}$. $u=2x-1, v=x+3$. $u'=2, v'=1$.
$f'(x) = \frac{2(x+3)-(2x-1)}{(x+3)^2} = \frac{7}{(x+3)^2} > 0$ sur $]-3;+\infty[$.
$f$ strictement croissante.
$f(1)=\frac{1}{4}$, $f'(1)=\frac{7}{16}$.
Tangente : $y = \frac{7}{16}(x-1) + \frac{1}{4} = \frac{7}{16}x - \frac{3}{16}$.
Asymptote verticale $x=-3$, horizontale $y=2$.
\end{corrige}

\newpage

\coursec{Fiche bilan}

\begin{popfigure}
\begin{tikzpicture}[mindmap, grow cyclic, every node/.style={concept, text=white, font=\small},
  root concept/.append style={concept color=popblue, font=\bfseries},
  level 1/.append style={level distance=3.6cm, sibling angle=60},
  level 1 concept/.append style={font=\small}]
\node[root concept]{Fonctions dérivées}
  child[concept color=poporange]{ node {Nombre dérivé $f'(a)$ : pente de la tangente} }
  child[concept color=popgreen]{ node {Tangente : $y=f'(a)(x-a)+f(a)$} }
  child[concept color=poppurple]{ node {Dérivées usuelles : $x^n$, $\sqrt{x}$, $\dfrac{1}{x}$} }
  child[concept color=popgold]{ node {Opérations : $u+v$, $ku$, $uv$, $\dfrac{u}{v}$} }
  child[concept color=poppink]{ node {Signe de $f'$ $\Rightarrow$ sens de variation} }
  child[concept color=popblue!70]{ node {Applications : vitesse, coût marginal, optimisation} };
\end{tikzpicture}
\legende{Carte mentale de la leçon : les six idées clés autour de la dérivation.}
\end{popfigure}

\begin{bilan}
\textbf{Définitions clés}
\begin{itemize}
  \item Le \cle{nombre dérivé} de $f$ en $a$ est la limite (quand elle existe) :
  \[ f'(a)=\lim_{h\to 0}\frac{f(a+h)-f(a)}{h}. \]
  C'est le \cle{coefficient directeur} de la tangente à la courbe au point d'abscisse $a$.
  \item La \cle{tangente} en $A\big(a\,;\,f(a)\big)$ a pour équation :
  \[ y=f'(a)(x-a)+f(a). \]
  \item Si $f'(a)$ existe pour tout $a$ d'un intervalle $I$, on définit la \cle{fonction dérivée} $f'$ sur $I$.
\end{itemize}

\textbf{Formules à connaître par c\oe ur}
\begin{center}
\begin{tabularx}{\linewidth}{|l|X|X|}\hline
\rowcolor{popblueL} \textbf{Fonction $f$} & \textbf{Dérivée $f'$} & \textbf{Conditions} \\ \hline
$f(x)=k$ (constante) & $f'(x)=0$ & sur $\mathbb{R}$ \\ \hline
$f(x)=x$ & $f'(x)=1$ & sur $\mathbb{R}$ \\ \hline
$f(x)=x^n$ ($n\in\mathbb{N}^*$) & $f'(x)=n\,x^{n-1}$ & sur $\mathbb{R}$ \\ \hline
$f(x)=\dfrac{1}{x}$ & $f'(x)=-\dfrac{1}{x^2}$ & sur $\mathbb{R}^*$ \\ \hline
$f(x)=\sqrt{x}$ & $f'(x)=\dfrac{1}{2\sqrt{x}}$ & sur $]0\,;\,+\infty[$ \\ \hline
\end{tabularx}
\end{center}

\begin{center}
\begin{tabularx}{\linewidth}{|l|X|}\hline
\rowcolor{popblueL} \textbf{Opération} & \textbf{Dérivée} \\ \hline
Somme & $(u+v)'=u'+v'$ \\ \hline
Produit par un réel & $(ku)'=k\,u'$ \\ \hline
Produit & $(uv)'=u'v+uv'$ \\ \hline
Quotient & $\left(\dfrac{u}{v}\right)'=\dfrac{u'v-uv'}{v^2}$ \quad ($v\neq 0$) \\ \hline
\end{tabularx}
\end{center}

\textbf{Méthodes express}
\begin{itemize}
  \item \cle{Calculer $f'(a)$ par la définition} : former le taux d'accroissement $\dfrac{f(a+h)-f(a)}{h}$, simplifier, puis faire tendre $h$ vers $0$.
  \item \cle{Équation de tangente} : calculer $f(a)$, calculer $f'(x)$ puis $f'(a)$, remplacer dans $y=f'(a)(x-a)+f(a)$, développer.
  \item \cle{Sens de variation} : dériver $f$, étudier le signe de $f'(x)$, conclure : $f'>0 \Rightarrow f$ croissante ; $f'<0 \Rightarrow f$ décroissante.
  \item \cle{Situation concrète} : la dérivée mesure une \emph{vitesse de variation} (vitesse d'un mobile, coût marginal, débit).
\end{itemize}
\end{bilan}

\begin{aretenir}[Checklist avant le Probatoire]
\begin{itemize}
  \item[$\square$] Je sais écrire la définition du nombre dérivé avec le taux d'accroissement.
  \item[$\square$] Je sais interpréter graphiquement $f'(a)$ comme la pente de la tangente.
  \item[$\square$] Je connais par c\oe ur l'équation de la tangente : $y=f'(a)(x-a)+f(a)$.
  \item[$\square$] Je dérive sans hésiter $x^2$, $x^3$, $x^n$, $\dfrac{1}{x}$ et $\sqrt{x}$.
  \item[$\square$] Je n'oublie pas que la dérivée d'une constante est nulle.
  \item[$\square$] J'applique correctement $(uv)'=u'v+uv'$ (et non $u'v'$).
  \item[$\square$] J'applique correctement $\left(\dfrac{u}{v}\right)'=\dfrac{u'v-uv'}{v^2}$ en vérifiant $v\neq 0$.
  \item[$\square$] Je sais dériver un polynôme terme à terme.
  \item[$\square$] Je relie le signe de $f'$ au sens de variation de $f$.
  \item[$\square$] Je sais exploiter la dérivée dans un problème concret (vitesse, coût, recette).
  \item[$\square$] Je rédige mes calculs en justifiant chaque étape et je vérifie mes conditions d'existence.
\end{itemize}
\end{aretenir}

\findecours

\end{document}
